\documentclass[11pt]{article}

\usepackage[margin=1in]{geometry}
\usepackage{amsmath,amssymb,amsthm}
\usepackage{mathtools}
\usepackage{enumitem}
\usepackage[hidelinks]{hyperref}
\usepackage[expansion=false]{microtype}
\usepackage{xr}


\newtheorem{theorem}{Theorem}[section]
\newtheorem{lemma}[theorem]{Lemma}
\newtheorem{proposition}[theorem]{Proposition}
\newtheorem{corollary}[theorem]{Corollary}
\newtheorem{conjecture}[theorem]{Conjecture}
\theoremstyle{definition}
\newtheorem{definition}[theorem]{Definition}
\newtheorem{hypothesis}[theorem]{Hypothesis}
\newtheorem{remark}[theorem]{Remark}

\DeclareMathOperator{\lcm}{lcm}
\newcommand{\NN}{\mathbb{N}}
\newcommand{\ZZ}{\mathbb{Z}}
\newcommand{\QQ}{\mathbb{Q}}
\newcommand{\PP}{\mathcal{P}}
\newcommand{\der}{{}'}
\newcommand{\dd}{\,\mathrm{d}}

\title{On the equation $n''=n$ and a problem of Erd\H{o}s and Barbeau \\
       on products of prime--reciprocal sums}
\author{Vico Bonfioli\thanks{\texttt{vicobonfioli@gmail.com}.
  See the acknowledgments for a candid note on the use of AI assistance in this work.}}
\date{\today}

\externaldocument[C-]{erdos307-companion}
\begin{document}
\maketitle

\begin{abstract}
Erd\H{o}s Problem~\#307 (posed by Barbeau, 1976; recorded in Erd\H{o}s--Graham, 1980) asks
whether there exist two finite sets $P,Q$ of primes with
$\bigl(\sum_{p\in P}\tfrac1p\bigr)\bigl(\sum_{q\in Q}\tfrac1q\bigr)=1$.
Writing $a=\prod P$, $b=\prod Q$, a solution exists if and only if the \emph{arithmetic derivative}
has a two--cycle $a'=b,\ b'=a$ with $a\neq b$; equivalently iff $n''=n$ has a solution that is not a
fixed point $p^{p}$. That reformulation, and the equivalence with the derivative, are due to
Ufnarovski--\AA hlander \cite{UfnarovskiAhlander2003}, whose Conjecture~4 asserts in exactly these
terms that $\bigl(\sum1/p_i\bigr)\bigl(\sum1/q_j\bigr)=1$ has no solution in distinct primes. Our
starting point is that \#307 \emph{is} that conjecture --- an identification neither literature
records --- which ties the problem to the Ufnarovski--\AA hlander dynamics: their conjecture that
every $n''=n$--cycle has period one would force \#307 to have answer \textsc{no}, while a
\textsc{yes} would refute it. Exploiting the equivalence we prove a \emph{rigidity lemma} (each prime--reciprocal
sum is automatically in lowest terms, so $N_P=D_Q$, $N_Q=D_P$, and $P\cap Q=\varnothing$) and a
quantitative \emph{barrier}: building on Rosen's bound $|P\cup Q|\ge 59$, any solution has
$\min(\prod P,\prod Q)\ge 2.09\times10^{56}$, for \emph{both} members. This is the first valid
unconditional bound of its kind. The prior claims are Kovi\v{c}'s: his Proposition~17 is sound but
conditional on the smaller member being odd, while Proposition~18 --- the only unconditional
cardinality and size bound in the literature, giving $r+s\ge34$ and
$\max\{m,n\}\ge\prod_{i\le17}p_i$ --- rests on an invalid step, deriving $p_1q_s<rs$ by multiplying
an upper bound against a lower one, where the available products give only $p_1q_1<rs<p_rq_s$ and no
cardinality bound at all; we give an explicit model and machine--check both points. Against the
strongest surviving unconditional bound, that paper's computer search $x>10^4$, the gain is $52$
orders of magnitude. A finite verification over the complete list of
$49{,}961$ admissible $59$--prime supports then closes the $59$--prime level: any solution in fact
has $|P\cup Q|\ge 60$ and $\min(\prod P,\prod Q)>3.50\times10^{57}$; a quadratic--reciprocity
certificate, requiring no factorisation, then proves the canonical tail family (the first $59$
primes plus any $q$) empty, and together with a factoring pass $46{,}483$ of the $49{,}961$
one--new--prime families at level $60$ ($93\%$; the residual is factoring--hard). We add a parity dichotomy (a hypothetical all--odd
cycle would need $\ge 1412$ primes and have $\min(\prod P,\prod Q)>10^{2519}$), a proof that longer
cycles are strictly harder ($k=3$ already forces $361{,}139$ primes), and a function--field theorem:
over $\mathbb{F}_q[t]$ the derivative has \emph{no} cycles, so the obstruction over $\mathbb{Z}$ is
purely archimedean. A local congruence anatomy of the cycle equation, the Erd\H{o}s--Wintner fact
that the abundance density $0.0420$ is a genuine constant, and a Th\=abit/Euler--type constructive
calculus (the \emph{Frame
Rule} and a divisor--trick) that faithfully parametrises solutions but, by an exact identity
$\Delta_0=M_0N(1-s_{M_0}s_N)$, provably does \emph{not} reduce the difficulty below \#307 itself
(the natural ``make $\Delta_0$ small'' objective is \#307 in disguise). No fixed--shape family can
help either: for families of \emph{any} growth rate, not merely polynomial ones, a product identity
holding infinitely often forces every member constant, which closes the Th\=abit/Euler programme
(primes $3\cdot2^{n}-1$ and relatives) that is the classical technique for problems of this kind. Recast as a bilinear breeder it yields a conditional reduction
$(\mathrm{H})\Rightarrow$\textsc{yes} to an explicit prime--density hypothesis, which a quantitative
ledger shows is not a feasible finite search. On the counting side we prove that both Pythagorean layers have density zero
unconditionally \emph{on the squarefree stratum}, which is where every tool in the proof lives and
where a \#307 solution lies by rigidity; the unrestricted statement is not proved. The minus layer $m'-2m=d^2$ is the set on which the arithmetic derivative takes
square values, so Heath-Brown's square sieve applies once the symbol is factorised as
$\bigl(\tfrac{m'-2m}{r}\bigr)=\bigl(\tfrac mr\bigr)\bigl(\tfrac{\sigma_r(m)-2}{r}\bigr)$, after
which the additivity of $\sigma_r$ makes the relevant sums multiplicative and Hal\'asz supplies the
cancellation. Keeping the sieve parameter a function of $N$ would give the rate
$\ll N\log\log N(\log N)^{-1/4}$, but that is stated as a \textsc{conjecture} and not proved here:
it needs a form of the cancellation lemma uniform in the modulus, which we do not establish. The
sharp exponent $N^{1/2}$ is not reached either, and we record three distinct reasons: the multiplicative decomposition scatters the cancellation across frequencies, a
bilinear decomposition is unavailable for a signed non-multiplicative summand, and the natural
large-sieve reformulation is circular at its principal term. The $34$ reciprocity-immune families at
level $60$ are upgraded from probable to proved by Atkin--Morain certificates, archived for
independent verification. The rigidity result, the barrier, and the arithmetic core of the counting
argument are formalised in Lean~4 (mathlib), extremality proved rather than assumed; the only
non--logical input is a verified evaluation of the first $59$ primes. The problem remains open. The heuristic case for existence is weaker than it appears: the constant in
the expected-count heuristic cannot be measured even in principle, since the barrier places the
entire relevant region above $10^{112.9}$ while the largest value of $\sigma(a)\sigma(a')$ observed
below $10^{7}$ is $0.5535$.
\end{abstract}


\section{Introduction}\label{sec:intro}

Let $P,Q$ be finite, nonempty sets of (distinct) primes. Erd\H{o}s Problem~\#307
\cite{ErdosProblems307}, attributed to Barbeau's 1976 ``Computer challenge corner'' problem
\#477 \cite{Barbeau1976} and listed in Erd\H{o}s--Graham \cite{ErdosGraham1980}, asks:
\begin{quote}
do there exist finite sets of primes $P,Q$ with
$\displaystyle\Bigl(\sum_{p\in P}\frac1p\Bigr)\Bigl(\sum_{q\in Q}\frac1q\Bigr)=1$\,?
\end{quote}
The problem is verifiable in the sense that a single example would settle it. No example is known.
The published partial facts are a remark of J.~Rosen on MathOverflow \cite{MO320838}, that any
solution needs at least $59$ primes in $P\cup Q$, and the problem page \cite{ErdosProblems307},
which asserts $|P\cup Q|\ge60$.

The relation between the two is worth stating precisely, since this paper proves the second. The
page's justification reads, in full, that $P$ and $Q$ are disjoint and
$\sum_{p\in P\cup Q}1/p\ge2$, ``whence $|P\cup Q|\ge60$''. That implication does not hold as
displayed: the reciprocals of the first $59$ primes already sum to $2.00235>2$, so a mass of $2$ is
attainable on $59$ primes and the stated hypotheses yield $|P\cup Q|\ge59$, which is Rosen's bound.
Reaching $60$ requires excluding the $49{,}961$ admissible $59$--element supports one at a time,
which is what Proposition~\ref{prop:close59} does and what \texttt{erdos307\_sixty} machine-checks.
We therefore claim no priority for the \emph{statement} $|P\cup Q|\ge60$, which predates this work on
the problem page, and we do claim the proof. The known
``solutions'' of Cambie and others solve only the \emph{weakened} version in which $1$ is permitted
as an element, e.g.\ $(1+\tfrac15)(\tfrac12+\tfrac13)=1$; these do not bear on the prime version.

Expanding the product exhibits \#307 as a rigidly shaped Egyptian fraction for $1$:
$1=\sum_{p\in P}\sum_{q\in Q}1/(pq)$, whose denominators are the \emph{rectangular grid} $P\times Q$
of semiprimes; a \emph{rank--one}, multiplicatively separable representation. The unrestricted
companion (representing $1$ by reciprocals of distinct two--prime--factor numbers, with no product
structure imposed) is by contrast \emph{solved}, and comprehensively so: every natural number is a
sum of distinct semiprime unit fractions \cite{Semiprime2026}, and every $a/b$ with $b$ squarefree
above an explicit threshold is likewise representable, so the $\omega=2$ problem is settled well
past the single value $1$. Johnson (1978) gave a $48$--term decomposition,
Watanabe \cite{Watanabe2020} $47$--term ones, and Graham's survey \cite{GrahamEgyptian} records exactly
this contrast, noting after Barbeau that the product form ``is not known.'' Thus \#307 is precisely the
rank--one restriction of a solved problem. Nor is the density--zero nature of the prime denominators the
obstruction in itself: Bloom \cite{BloomShifted} represents every positive rational by distinct unit
fractions drawn from the (equally sparse) shifted primes $\{p-1\}$. What resists is the rank--one
rigidity (the demand that the representation \emph{factor} as a product of two prime--reciprocal
sums) which the barrier (Theorem~\ref{thm:barrier}) shows costs at least $59$ primes, hence of order
$870$ semiprime terms against Johnson's $47$.

Our starting point is an identification of \#307 with a question about the \emph{arithmetic
derivative} $n\mapsto n'$, the unique map with $0'=1'=0$, $p'=1$ for primes $p$, and the Leibniz
rule $(ab)'=a'b+ab'$. (The map goes back to J.~Mingot Shelly (1911) and appeared as a 1950 Putnam
problem; Barbeau \cite{Barbeau1961} gave the first systematic study, with the modern account in
Ufnarovski--\AA hlander \cite{UfnarovskiAhlander2003} and OEIS \cite{OEIS_A003415}.) For
$n=\prod_i p_i^{e_i}$ one has $n'=n\sum_i e_i/p_i$. We show (Section~\ref{sec:bridge}) that \#307
is \emph{equivalent} to the existence of a two--cycle of $n\mapsto n'$ that is not a fixed point.
The two--cycle reformulation itself is not new, and it is older than is generally recorded.
Ufnarovski--\AA hlander \cite{UfnarovskiAhlander2003} prove that a period--two point of $n\mapsto n'$
has squarefree members with disjoint supports, and deduce that the period--two case of their
Conjecture~3 is equivalent to their Conjecture~4: $\bigl(\sum_{i\le k}1/p_i\bigr)
\bigl(\sum_{j\le l}1/q_j\bigr)=1$ has no solution in distinct primes. That is \#307, in the
derivative literature, in 2003. It was restated ``in disguise'' by Seva \cite{Seva323194} in 2019
(via $Q\mapsto Q\sum_{p\mid Q}1/p$, which agrees with the arithmetic derivative on squarefree
integers, where the question lives), who is the first to name it as \#307, and obtained
independently by Bado \cite{Bado2026}; see Section~\ref{sec:prior}. What is new here is the
identification of the two problems --- \#307 is Ufnarovski--\AA hlander's Conjecture~4 --- and the
resulting two--way bridge, importing their period--one conjecture and exporting our bounds.
Remarkably, Barbeau authored foundational papers on both sides (1961 and 1976
\cite{Barbeau1961,Barbeau1976}).

The bridge is productive in both directions.
\begin{itemize}[leftmargin=1.4em]
\item \emph{Importing} from the derivative literature: the Ufnarovski--\AA hlander theorem that any
two--cycle has squarefree members with disjoint supports \cite{UfnarovskiAhlander2003} makes the
``finite sets of distinct primes'' hypothesis of \#307 automatic, and links \#307 to the
Ufnarovski--\AA hlander conjecture that all derivative cycles have period one.
\item \emph{Exporting} to that literature: our barrier (Theorem~\ref{thm:barrier}) is, we believe,
the first valid unconditional bound of its kind. Kovi\v{c} \cite{Kovic2012} states three bounds on an $n''=n$ solution. A computer
search gives $x>10^4$, unconditionally. Proposition~17 gives $>3\cdot5\cdots29\approx3.23\times10^9$,
but only when the \emph{smaller} member is odd; Kovi\v{c} notes explicitly that no estimate follows
otherwise. Proposition~18 claims the unconditional $r+s\ge34$ and $\max\{m,n\}\ge\prod_{i=1}^{17}p_i
\approx1.92\times10^{21}$; its proof does not establish them, for the reason set out in
Section~\ref{sec:prior}. Against the strongest \emph{valid unconditional} bound, $10^4$, our
$\min(\prod P,\prod Q)\ge2.09\times10^{56}$ is a gain of $52$ orders of magnitude; against the
conditional Proposition~17 it is $47$, in the case where that proposition applies. On cardinality
there is no prior bound to improve: $|U|\ge60$ appears to be the first, and it re-derives
Kovi\v{c}'s claimed $r+s\ge34$ as a corollary by an unrelated route.
\end{itemize}

\paragraph{Results and organisation.}
Section~\ref{sec:rigidity} proves the Rigidity Lemma (Theorem~\ref{thm:rigidity}): the fraction
$\sum_{p\in S}1/p$ is automatically in lowest terms, which forces, for any solution, the exact
identities $N_P=D_Q$, $N_Q=D_P$, disjointness $P\cap Q=\varnothing$, and $D_Q=s\,D_P$ where
$s=\sum_{p\in P}1/p$. The disjointness core was anticipated by R.~Israel \cite{MO320838}; the full
rigidity is what powers everything later. Section~\ref{sec:bridge} establishes the bridge.
Section~\ref{sec:barrier} proves the barrier, the parity dichotomy, and that longer cycles are
strictly harder (so two--cycles are extremal). Section~\ref{C-sec:anatomy} gives a local congruence
anatomy, a heuristic model, and a proof (via Erd\H{o}s--Wintner) that the relevant abundance density
is a genuine constant. Section~\ref{C-sec:ff} adds a function--field lens: over $\mathbb{F}_q[t]$ the
derivative has \emph{no} cycles, so the obstruction in $\mathbb{Z}$ is purely archimedean.
Section~\ref{sec:frame} develops the Frame Rule and divisor--trick,
and, crucially, records the exact identity that renders the natural constructive objective
circular. Section~\ref{C-sec:breeder} recasts the rule as a bilinear breeder, records a conditional
reduction $(\mathrm{H})\Rightarrow$ \textsc{yes} to an explicit prime--density hypothesis, and shows
quantitatively why $(\mathrm{H})$ does not amount to a feasible finite search.
Section~\ref{sec:computations} documents the computations, and Section~\ref{sec:lean} the
Lean~4 formalisation. We are explicit throughout about what is proved, what is heuristic, and what
is borrowed; the problem remains open.

\paragraph{Notation.}
For a finite set $S$ of distinct primes put
\[
D_S \coloneqq \prod_{p\in S} p, \qquad N_S \coloneqq \sum_{p\in S} \frac{D_S}{p},
\qquad\text{so}\qquad \sum_{p\in S}\frac1p = \frac{N_S}{D_S}.
\]
We call $N_S$ the \emph{cofactor sum}. For a solution we write $a=D_P$, $b=D_Q$,
$s=N_P/D_P=\sum_{p\in P}1/p$ and $t=N_Q/D_Q=\sum_{q\in Q}1/q$, so $st=1$, and (without loss of
generality) $s>1>t$. We write $\omega(n)$ for the number of distinct primes dividing $n$.

\section{The Rigidity Lemma}\label{sec:rigidity}

\begin{lemma}[Rigidity / lowest terms]\label{thm:rigidity}
For a finite set $S$ of distinct primes, $\gcd(N_S,D_S)=1$. Equivalently, $\sum_{p\in S}1/p$ is
already in lowest terms, and $v_p\!\bigl(\sum_{r\in S}1/r\bigr)=-1$ for every $p\in S$.
\end{lemma}

\begin{proof}
The prime divisors of $D_S=\prod_{r\in S}r$ are exactly the elements of $S$, so it suffices to show
$p\nmid N_S$ for each $p\in S$. Fix $p\in S$ and reduce $N_S=\sum_{r\in S}D_S/r$ modulo $p$. For
$r\ne p$ the term $D_S/r=\prod_{u\in S\setminus\{r\}}u$ still contains the factor $p$, so
$p\mid D_S/r$. The single surviving term is $D_S/p=\prod_{u\in S\setminus\{p\}}u$, whence
\[
N_S \equiv \prod_{u\in S\setminus\{p\}} u \pmod p.
\]
Each $u$ in this product is a prime distinct from $p$, so $p\nmid u$; as $p$ is prime it does not
divide the product. Thus $N_S\not\equiv0\pmod p$, i.e.\ $p\nmid N_S$. As $p$ ranged over all prime
divisors of $D_S$, $\gcd(N_S,D_S)=1$. The valuation statement is the same computation:
$v_p(N_S/D_S)=v_p(N_S)-v_p(D_S)=0-1=-1$.
\end{proof}

\begin{remark}
Distinctness and primality are both essential: $v_p$ of the cofactor of a repeated prime need not be
$-1$, and the final step uses Euclid's lemma. The case $|S|=1$ is consistent ($N_{\{p\}}=1$).
\end{remark}

\begin{theorem}[Solution structure]\label{thm:structure}
Suppose $P,Q$ are finite sets of distinct primes with $\dfrac{N_P}{D_P}\cdot\dfrac{N_Q}{D_Q}=1$.
Then
\[
N_P=D_Q,\qquad N_Q=D_P,\qquad P\cap Q=\varnothing,
\]
$N_P$ is squarefree with prime support exactly $Q$ (and $N_Q$ squarefree with support $P$), and
$D_Q=s\,D_P$ where $s=N_P/D_P$.
\end{theorem}

\begin{proof}
Cross-multiplying, $N_PN_Q=D_PD_Q$. By Lemma~\ref{thm:rigidity} both $N_P/D_P$ and $N_Q/D_Q$ are in
lowest terms. Since $N_Q/D_Q=D_P/N_P$ and the lowest--terms representative of a positive rational is
unique, comparing $N_Q/D_Q$ with $D_P/N_P$ gives $N_Q=D_P$ and $D_Q=N_P$.

\emph{Disjointness} (the Israel mechanism \cite{MO320838}). If $r\in P\cap Q$ then $r\mid D_P$ and,
by $N_P=D_Q=\prod Q$, also $r\mid N_P$, but $\gcd(N_P,D_P)=1$ forces $r\nmid N_P$, a contradiction.
(Equivalently $v_r(s)=v_r(t)=-1$, so $v_r(st)=-2\neq0=v_r(1)$.) Hence $P\cap Q=\varnothing$.

Finally $N_P=D_Q=\prod_{q\in Q}q$ is a product of distinct primes, hence squarefree with prime
support $Q$; symmetrically for $N_Q$. And $s\,D_P=(N_P/D_P)D_P=N_P=D_Q$.
\end{proof}

The disjointness in Theorem~\ref{thm:structure} means a solution is genuinely a partition of
$U\coloneqq P\cup Q$, and $D_PD_Q=\prod_{p\in U}p$. The valuation form $v_p(\sum 1/p_i)=-1$ underlying
disjointness is due to R.~Israel \cite{MO320838}, and appears also in Seva \cite{Seva323194}; the
cross--matching form ($N_P=D_Q$, $N_Q=D_P$) was obtained independently and concurrently by Bado
\cite[Thm.~A]{Bado2026} (Section~\ref{sec:prior}).

\section{The bridge to the arithmetic derivative}\label{sec:bridge}

\begin{lemma}\label{lem:der-cofactor}
For squarefree $m$, the arithmetic derivative equals the cofactor sum: $m'=\sum_{p\mid m}m/p=N_{\{p:\,p\mid m\}}$. Moreover $\gcd(m,m')=1$.
\end{lemma}

\begin{proof}
For $m=\prod_{p\mid m}p$ squarefree, $m'=m\sum_{p\mid m}1/p=\sum_{p\mid m}m/p$, which is exactly the
cofactor sum of the prime set of $m$. Coprimality is Lemma~\ref{thm:rigidity}.
\end{proof}

\begin{theorem}[Bridge]\label{thm:bridge}
The following are equivalent.
\begin{enumerate}[label=\textup{(\arabic*)},leftmargin=2.2em]
\item Erd\H{o}s \#307 has a solution: finite prime sets $P,Q$ with
$\bigl(\sum_{p\in P}1/p\bigr)\bigl(\sum_{q\in Q}1/q\bigr)=1$.
\item The arithmetic derivative has a two--cycle: integers $a\ne b$ with $a'=b$ and $b'=a$.
\item The equation $n''=n$ has a solution $n$ that is not a fixed point (i.e.\ $n'\ne n$, so $n$ is
not of the form $p^{p}$).
\end{enumerate}
\textup{(}The implication $(2)\Rightarrow(1)$ uses the Ufnarovski--\AA hlander squarefreeness
theorem for cycles; all other implications, and everything used in the barrier of
Section~\textup{\ref{sec:barrier}}, are unconditional; see Remark~\textup{\ref{rem:cited}}.\textup{)}
\end{theorem}

\begin{proof}
$(1)\Rightarrow(2)$. Let $a=D_P,b=D_Q$. By Lemma~\ref{lem:der-cofactor},
$a'=N_P$ and $b'=N_Q$, and by Theorem~\ref{thm:structure}, $N_P=D_Q=b$ and $N_Q=D_P=a$. Thus
$a'=b$, $b'=a$, and $a\ne b$ since $s\ne1$ (a prime--reciprocal sum is never $1$, being
$N/D$ in lowest terms with $D>1$).

$(2)\Rightarrow(1)$. Let $a'=b,b'=a$ with $a\ne b$. By the Ufnarovski--\AA hlander analysis of
two--cycles (cf.\ \cite{UfnarovskiAhlander2003,Kovic2012}), both $a$ and $b$ are squarefree; their
prime supports $P,Q$ are then disjoint since $\gcd(a,b)=\gcd(a,a')=1$. By
Lemma~\ref{lem:der-cofactor}, $a'=N_P$ and $b'=N_Q$; the
cycle gives $N_P=b=D_Q$ and $N_Q=a=D_P$, so
$\bigl(\sum_{p\in P}1/p\bigr)\bigl(\sum_{q\in Q}1/q\bigr)=\frac{N_P}{D_P}\frac{N_Q}{D_Q}
=\frac{D_Q}{D_P}\frac{D_P}{D_Q}=1$.

$(2)\Leftrightarrow(3)$. A two--cycle $a\ne b$ gives $a''=a$ with $a'=b\ne a$, so $a$ is a non--fixed
solution of $n''=n$. Conversely a solution of $n''=n$ with $n'\ne n$ yields the two--cycle
$(n,n')$. The fixed points of $n\mapsto n'$ are exactly the numbers $p^{p}$
\cite{Barbeau1961,UfnarovskiAhlander2003}.
\end{proof}

\begin{remark}[Status of the cited input]\label{rem:cited}
Part $(2)\Rightarrow(1)$ uses that a derivative two--cycle has squarefree, support--disjoint
members. Squarefreeness is the content of the Ufnarovski--\AA hlander analysis of cycles (a
non--squarefree $n$ has $p\mid\gcd(n,n')$ for some $p$, and along $n\mapsto n'$ such common factors
propagate, precluding return); disjointness then follows from $\gcd(a,a')=1$. We use it as a cited
theorem. Even \emph{without} it, the implication $(1)\Rightarrow(2)$ and the entire barrier below are
unconditional, because Rigidity already delivers squarefree, disjoint $P,Q$ from a \#307 solution.

\textup{Formalised.} \texttt{Erdos307.bridge\_forward} is $(1)\Rightarrow(2)$ and cites nothing;
the distinctness $a\ne b$ is \emph{derived} there, not assumed
\textup{(}\texttt{dprod\_ne\_dprod}\textup{)}. Support--disjointness is available separately as
\texttt{disjoint\_primeFactors\_of\_cycle}, also from Rigidity; it is not consumed by
\texttt{bridge\_forward}, which does not need it.
\texttt{bridge\_backward} is $(2)\Rightarrow(1)$ with the Ufnarovski--\AA hlander squarefreeness
appearing as an explicit hypothesis, and \texttt{two\_cycle\_iff\_double\_fixed} is
$(2)\Leftrightarrow(3)$ for an arbitrary self--map, on \emph{no} axioms at all. $0$ \texttt{sorry}
\textup{(\texttt{lean/Erdos307/Bridge.lean})}.
\end{remark}

\begin{corollary}\label{cor:ua}
If the Ufnarovski--\AA hlander conjecture holds in the weak form ``every cycle of $n\mapsto n'$ has
period one,'' then Erd\H{o}s \#307 has answer \textsc{no}. Conversely, a \textsc{yes} to \#307
exhibits a period--two cycle and refutes that statement.
\end{corollary}

\section{The barrier}\label{sec:barrier}

\begin{lemma}[Strict AM--GM]\label{lem:amgm}
If $s,t>0$ with $st=1$ and $s\ne1$, then $s+t>2$.
\end{lemma}
\begin{proof}
$s+t-2=s+1/s-2=(s-1)^2/s>0$ since $s\ne1$.
\end{proof}

\begin{lemma}[The number $59$, after Rosen \cite{MO320838}]\label{lem:59}
Any finite set $U$ of distinct primes with $\sum_{p\in U}1/p>2$ has $|U|\ge 59$. Moreover, among
$k$--element prime sets, $\sum 1/p$ is maximised by the $k$ smallest primes.
\end{lemma}
\begin{proof}
Replacing any prime of $U$ by a smaller prime not in $U$ increases $\sum1/p$, so the maximum over
$k$--element sets is at the first $k$ primes. The reciprocal sum of the first $58$ primes is
$1.99874\ldots<2$ and of the first $59$ primes is $2.00235\ldots>2$, so a sum exceeding $2$ needs at
least $59$ primes.
\end{proof}

\begin{theorem}[Barrier]\label{thm:barrier}
For any solution $(P,Q)$ of \#307, with $U=P\cup Q$:
\begin{enumerate}[label=\textup{(\alph*)},leftmargin=2.2em]
\item $|U|\ge 59$ and $\displaystyle\prod_{p\in U}p \ge \Pi_{59}\coloneqq\prod_{i=1}^{59}p_i
\approx 8.77\times10^{112}$;
\item $\displaystyle\min(\textstyle\prod P,\prod Q)\ \ge\ \sqrt{\Pi_{59}/T_{59}} \ =\
2.0929\ldots\times10^{56}$, where $T_{59}=\sum_{i=1}^{59}1/p_i=2.00235\ldots$. In particular both
$\prod P$ and $\prod Q$ exceed $2.09\times10^{56}$.
\end{enumerate}
Consequently every $n''=n$ solution has both members $>2.09\times10^{56}$. The strongest valid
unconditional bound previously available is the computer search $x>10^4$ of \cite{Kovic2012}, so this
is a gain of $52$ orders of magnitude; against the conditional Proposition~17 of that paper
($3.23\times10^9$, when the smaller member is odd) it is $47$. The unconditional $r+s\ge34$ and
$\max\{m,n\}\ge1.92\times10^{21}$ claimed in Proposition~18 of \cite{Kovic2012} are not established
by their proof (Section~\ref{sec:prior}); the bound above does not rely on them and in fact implies
the first.
\end{theorem}

\begin{proof}
(a) By $st=1$, $s\ne1$ and Lemma~\ref{lem:amgm}, $\sum_{p\in U}1/p=s+t>2$, so $|U|\ge 59$ by
Lemma~\ref{lem:59}; since $P,Q$ are disjoint (Theorem~\ref{thm:structure}),
$\prod_{p\in U}p=D_PD_Q$, and a product of $\ge 59$ distinct primes is at least the product of the
$59$ smallest, $\Pi_{59}$.

(b) Take $D_P\le D_Q$ to be the minimal side. By Theorem~\ref{thm:structure}, $D_Q=s\,D_P$ with
$s>1$, hence
\[
s\,D_P^2 \;=\; D_P\,D_Q \;=\; \prod_{p\in U}p \;=\!:\, R,
\qquad\text{so}\qquad D_P^2=\frac{R}{s}.
\]
Since $s\le s+t=T(U)\coloneqq\sum_{p\in U}1/p$ (as $t>0$),
\[
D_P^2=\frac{R}{s}\ \ge\ \frac{R}{T(U)}.
\]
It remains to bound $R/T(U)$ below over all admissible $U$. Among $k$--element prime sets the
quotient $\bigl(\prod U\bigr)\big/\bigl(\sum_{p\in U}1/p\bigr)$ is minimised by the $k$ smallest
primes (smallest numerator and, by Lemma~\ref{lem:59}, largest denominator simultaneously), and
passing from $59$ to more primes only increases it (the product grows multiplicatively while the
reciprocal sum grows by a single small term). Hence the global minimum over admissible $U$ is at
$U=\{p_1,\dots,p_{59}\}$, giving $R/T(U)\ge \Pi_{59}/T_{59}$. Therefore
\[
D_P^2 \ \ge\ \frac{\Pi_{59}}{T_{59}} \;=\; \frac{\Pi_{59}^2}{N_{59}}\;=\;4.3806\ldots\times10^{112},
\qquad D_P\ \ge\ 2.0929\ldots\times10^{56},
\]
where $N_{59}=T_{59}\,\Pi_{59}$ is the integer cofactor sum of the first $59$ primes.
\end{proof}

\begin{remark}[Why the \emph{minimum} is bounded, not only the maximum]
The product identity $D_PD_Q\ge\Pi_{59}$ alone gives only $\max(\prod P,\prod Q)\ge
\sqrt{\Pi_{59}}\approx2.96\times10^{56}$. The stronger statement~(b), a lower bound on the
\emph{smaller} side, genuinely uses rigidity: it is the identity $D_Q=s\,D_P$ (equivalently
$D_P^2=R/s$) together with $s\le T(U)$ that converts the product bound into a bound on $D_P$. We
verified numerically that $\min_U R(U)/T(U)$ is attained exactly at the first $59$ primes (the
minimising configuration is unique), so the constant $2.09\times10^{56}$ is sharp for this method.
\end{remark}

\begin{proposition}[Closing the $59$--prime level]\label{prop:close59}
No solution of \#307 has $|P\cup Q|=59$. Hence $|P\cup Q|\ge60$,
$\prod_{p\in U}p\ \ge\ \Pi_{60}>2.46\times10^{115}$, and
\[ \min\bigl(\textstyle\prod P,\prod Q\bigr)\ \ge\ \sqrt{\Pi_{60}/T_{60}}\ =\
3.5053\ldots\times10^{57}. \]
\end{proposition}

\begin{proof}[Proof (finite verification)]
Let $U=P\cup Q$ with $|U|=59$; as in Theorem~\ref{thm:barrier}(a), $T(U)>2$. The candidate list is
finite and complete: \textup{(i)} every prime $p\le167$ lies in $U$; omitting $p$ caps $T(U)$ at
$T_{60}-1/p<2$, since the $59$ largest prime reciprocals avoiding $p$ are those of the first $60$
primes less $p$, and $1/p>T_{60}-2=0.00590\ldots$ exactly for $p\le167$; \textup{(ii)} every element
$p'\ge281$ satisfies $T_{58}+1/p'>2$ (the other $58$ reciprocals sum to at most $T_{58}$), so
$p'\le787$. Thus $U$ consists of the $39$ primes up to $167$ plus twenty primes from
$(167,787]$, and exactly $49{,}961$ such sets have $T(U)>2$. For each, a solution supported on $U$
forces both $N'+2N=(a+b)^2$ and $N'-2N=(a-b)^2$ to be perfect squares, where $N=\prod U$ and
$N'=a^2+b^2$ (rigidity; the $k=0$ case of Proposition~\ref{prop:pairform}); one integer test that
simultaneously excludes all $2^{58}$ splits of $U$. Exact--integer verification shows $N'+2N$ is a
non--square for \emph{every} one of the $49{,}961$ candidates, by two independent implementations
with different enumeration orders and prunings (\texttt{code/close59.py} and
\texttt{hunt/close59.rs} in the repository). Hence no $59$--element support closes, and
$|U|\ge60$. The bound on the minimal side repeats the proof of Theorem~\ref{thm:barrier}(b) with the
minimum of $R/T$ now over $|U|\ge60$, attained at the first $60$ primes:
$D_P^2\ge\Pi_{60}/T_{60}=\Pi_{60}^2/N_{60}$, where $N_{60}$ is the cofactor sum of the first $60$
primes.
\end{proof}

\begin{corollary}[The coprime variant inherits the closed level]\label{cor:coprime60}
Consider the weakened variant in which the elements of $P$ and $Q$ are merely pairwise coprime
integers $\ge2$ \textup{(}the open form of the variant, $1\notin P\cup Q$\textup{)}. Any solution
has $|P\cup Q|\ge 60$ elements, and $\min(\prod P,\prod Q)>3.50\times10^{57}$.
\end{corollary}
\begin{proof}
Rigidity holds verbatim for pairwise--coprime families (the cofactor--sum argument uses only
coprimality), so again $T(U)=s+t>2$. The least prime factors $\ell(e)$ of the elements are distinct,
and $\sum_e 1/e\le\sum_e1/\ell(e)$. Suppose $|U|=59$ and some element $e$ is not prime, so
$e\ge\ell(e)^2$. \textup{(}Formalised: \texttt{Erdos307.coprime\_loss} is the per--element bound,
\texttt{loss\_antitone} its monotonicity, and \texttt{loss\_at\_277} the numeric consequence
\textup{(}\texttt{lean/Erdos307/Coprime60.lean}\textup{)}; the slack itself is the level--59
census.\textup{)} If $\ell(e)\le277$, the loss is
$1/\ell-1/e\ge(\ell-1)/\ell^2\ge276/277^2=0.00359\ldots$, exceeding the maximal available slack
$T_{59}-2=0.00235\ldots$; if $\ell(e)\ge281$, then $e\ge281^2$ while the other $58$ reciprocals sum
to at most $T_{58}$, giving $T(U)\le T_{58}+281^{-2}=1.99875\ldots<2$. Either way $T(U)<2$, a
contradiction: at $59$ elements every element is prime, and Proposition~\ref{prop:close59} applies
verbatim. The product bound transfers because $\prod_e e\ge\prod_e\ell(e)$ and
$T(U)\le\sum_e1/\ell(e)$, so the extremal ratio bound at the first $60$ primes still applies.
\end{proof}

\begin{remark}[Verification tier, and why the ladder stops here]\label{rem:tier60}
Theorem~\ref{thm:barrier} is machine--checked in Lean~4, and so is
Proposition~\ref{prop:close59}: the theorem \texttt{erdos307\_sixty} formalises the entire
argument; the forced--primes and element bounds, the Pythagorean plus--square, and a pruned
depth--first search over the pool whose \emph{completeness is proved, not assumed}
(\texttt{dfs\_sound}), the search itself running under the kernel; on top of the
two independent unverified programs. The level cannot be climbed further by finite means: already at $|U|=60$ the family
$U=\{p_1,\dots,p_{59}\}\cup\{q\}$ qualifies for \emph{every} prime $q>277$, and on it the two square
conditions read $x^2=Aq+\Pi_{59}$, $y^2=Bq+\Pi_{59}$ with $A,B=N_{59}\pm2\Pi_{59}$, i.e.\ the
Pell--type system $Bx^2-Ay^2=-4\Pi_{59}^2$ with $113$--digit coefficients. We verified that no
accessible modulus obstructs it (mod $32$, and all odd $\ell<4000$: the conditions pin
$q\equiv5\pmod 8$ and thin each $\ell\mid\Pi_{59}$ to about half of its residue classes and each
$\ell\nmid\Pi_{59}$ to about a quarter, but never to zero), consistent with the local--solubility
theme of Section~\ref{C-sec:anatomy}. The obstruction, however, lives at the primes \emph{of} $A$
and $B$, and it is detectable \emph{without factoring}: Proposition~\ref{prop:tailkill} below
closes this family for \emph{every} $q$ by a single Jacobi--symbol computation. (An exact sweep to
$q\le10^{12}$, \texttt{hunt/tail\_sweep.rs}, $1.25\times10^{11}$ candidates, none with
$Aq+\Pi_{59}$ even a single square, which corroborates it independently.) The tail family also exposes the self--similarity of the ladder:
writing $a=qm$, so that $m\,b=\Pi_{59}$ partitions the first $59$ primes, its two closing
conditions collapse to the exact equation
\[ D(m)\,D(b)+m^2\;=\;\Pi_{59} \]
together with the divisibility $m\mid D(b)$ (then $q=D(b)/m$, required prime); verified
exhaustively on a $12$--prime toy universe. Equivalently, the deficit $1-\sigma(m)\sigma(b)$ must
equal $m^2/\Pi_{59}$ \emph{exactly}: the circularity $\Delta_0=M_0N(1-s_{M_0}s_N)$ of
Section~\ref{sec:frame} in its purest form, with the deviation forced to be the perfect square
$m^2$. Closing even this one family is thus a $2^{58}$--fold exact--coincidence search of the same
type as \#307 itself (expected count ${\sim}2^{58}/\Pi_{59}\approx10^{-95}$): each rung of the
ladder past $60$ contains the whole problem. The one case with \emph{no} free large prime (a
$60$--set drawn entirely from small primes) is instead a finite search, and empty in the
minimal range: a direct both--squares enumeration of all $60$--prime sets with $T>2$ supported in
the first $70$ primes ($4{,}011{,}289$ sets; \texttt{hunt/close60.rs}) finds none passing even the
plus--square, so no level--$60$ solution has all its primes $\le349$. (The first $68$ primes give
$1{,}258{,}448$ sets, also empty, and reproduce independently.) By
Corollary~\ref{C-cor:kdetermined} this verdict is decisive rather than provisional: at these masses a
support passing the plus-- and minus--squares, with the parity of Proposition~\ref{C-prop:ladder}(2),
\emph{is} a two--cycle, so a survivor would have been a solution and not a candidate. The
enumeration extends with compute.
\end{remark}

\begin{proposition}[The canonical tail family is empty: a reciprocity kill]\label{prop:tailkill}
No solution of \#307 has support $U=\{p_1,\dots,p_{59}\}\cup\{q\}$, for any prime $q$.
\end{proposition}
\begin{proof}
A solution supported on $U$ forces $x^2=Aq+\Pi_{59}$ with $x=a+b$ and $A=N_{59}+2\Pi_{59}$
(Remark~\ref{rem:tier60}). Every prime $\ell\mid A$ exceeds $277$: for $\ell\le277$ one has
$\ell\mid\Pi_{59}$, so $\ell\mid A$ would force $\ell\mid N_{59}$, contradicting
$\gcd(N_{59},\Pi_{59})=1$ (rigidity). Hence $\ell\nmid\Pi_{59}$, and reduction mod $\ell$ gives
$x^2\equiv\Pi_{59}\pmod\ell$, requiring the Legendre symbol $(\Pi_{59}\mid\ell)\neq-1$. But the
Jacobi symbol satisfies
\[ \left(\tfrac{\Pi_{59}}{A}\right)\;=\;-1 \]
(a finite gcd--speed computation, script \texttt{code/tailkill.py}; no factorisation of the
$114$--digit $A$ is needed), so some $\ell\mid A$ of odd multiplicity has $(\Pi_{59}\mid\ell)=-1$.
The family is empty, for every $q$.

The companion symbol at $B=N_{59}-2\Pi_{59}$ equals $-1$ as well; necessarily: for any
admissible base $S$ (writing $D=D_S$, $N=N_S$, $A_S,B_S=N\pm2D$) one has $A_S\equiv B_S\pmod 8$
(as $4D\equiv8\bmod{16}$) and $A_S\equiv B_S\equiv N\pmod p$ for every odd $p\mid D$, while the
quadratic--reciprocity sign difference carries the exponent $\tfrac{p-1}{2}\cdot\tfrac{A_S-B_S}{2}$
with $\tfrac{A_S-B_S}{2}=2D$ even; hence $(D\mid A_S)=(D\mid B_S)$ always; one binary
certificate per tail family. Running it over all $49{,}961$ admissible bases of
Proposition~\ref{prop:close59} proves $31{,}219$ of the one--new--prime families at level $60$
empty ($62.5\%$, curiously near $5/8$); the remaining $18{,}742$ are Jacobi--inconclusive: their
$-1$--Legendre primes, if any, occur to an even count, visible only through a factorisation of
$A_S$. A factoring attack on those (trial division and Brent--Pollard $\rho$ over a Montgomery
core, \texttt{hunt/survivor\_kill.rs}; the found factors are kept as a pairwise--coprime basis, so
a kill certificate (an odd--multiplicity basis piece $P$ with $(D_S\mid P)=-1$) needs no
primality testing at all), pushed to $\rho$--budget $10^6$ with a Pollard $p-1$ stage, closes a
further $15{,}264$: in total $46{,}483$ of the $49{,}961$ one--new--prime families at level $60$
are proven empty ($93.0\%$). The $3{,}478$ still open resist all of this for a structural reason,
not a want of effort: each is Baillie--PSW \emph{composite} yet has no prime factor below ${\sim}
10^{12}$, so its $A_S$ (or $B_S$) is a \emph{balanced} semiprime of two ${\sim}57$--digit primes,
splitting one is factoring at cryptographic scale. Two cheap tests confirm no shortcut survives:
a product--tree batch--GCD over all $37{,}484$ values $A_S,B_S$ finds only shared primes below
$10^9$ (already found by $\rho$, so \emph{no} free kills), and Baillie--PSW leaves exactly
$34$ families in which $A_S$ and $B_S$ are both \emph{probable} primes; these are immune
\emph{conditionally on that primality} (a prime $A_S$ with $(D_S\mid A_S)=+1$ has no inert factor,
hence no congruence obstruction for any $q$). The conditionality is not pedantry and not
removable by more of the same computation: Baillie--PSW is rigorous only in its \emph{composite}
verdicts, so the $46{,}483$ kills are unconditional while the $34$ immunities are not. Upgrading them
would require primality \emph{certificates} for thirty-four $114$--digit integers (ECPP or a
Pocklington factorisation of $A_S-1$). The elementary route does not reach: immunity needs
\emph{both} $A_S$ and $B_S$ prime, and partial factorisation of $n-1$ (trial division to
$3\times10^6$ together with the fully split part) leaves $\log F/\log n$ below $1/3$ on at least one
side of every one of the $34$, so neither Pocklington ($F>\sqrt n$) nor
Brillhart--Lehmer--Selfridge ($F>n^{1/3}$) applies to any family: $0$ of $34$ are certifiable this
way (\texttt{code/immune\_certify.py}). This is now settled. Running the Adleman--Pomerance--Rumely
--Cohen--Lenstra test on all $68$ integers proves every one of them prime in $8$ seconds of
wall clock, so the $34$ immunities upgrade from \textup{[C]} to \textup{[P]} and are
\emph{unconditional} (\texttt{code/immune\_prove.gp}, PARI/GP $2.17.4$, whose \texttt{isprime} is a
proof and not a probable--prime test). The run reproduces the whole chain independently: $49{,}961$
admissible bases, $31{,}219$ Jacobi kills, and, among the $81$ bases whose $A_S$ and $B_S$ are both
prime, exactly the $34$ with $(D_S\mid A_S)=+1$; the other $47$ carry $(D_S\mid A_S)=-1$ and were
already killed. The primality is moreover backed by
\emph{independently checkable} evidence rather than a verdict: \texttt{code/immune\_certify.gp}
emits Atkin--Morain \textup{ECPP} certificates for all $68$ integers, archived at
\texttt{certs/immune\_ecpp.txt} ($512$\,KB) so a third
party can check them without trusting the prover. That file is a \texttt{primecertexport} dump and
is not machine--readable: PARI's \texttt{read} parses its header as a polynomial, so
\texttt{primecertisvalid} cannot consume it, and earlier versions of this paper and of the
repository documented that command in error. The working check is
\texttt{code/immune\_cert\_verify.gp}, which extracts the $68$ asserted subjects and re--proves each
by \textup{ECPP} in $1.3$ seconds: $68/68$ prime, $0$ failures. The structural half of the campaign is now
formal: the identity $(D_S\mid A_S)=(D_S\mid B_S)$, which is what makes one binary test decide a
family rather than two, is proved in Lean~4 as \texttt{Erdos307.jacobiSym\_A\_eq\_B}
\textup{(\texttt{lean/Erdos307/Campaign.lean})}, on the three standard axioms and with no
\texttt{native\_decide}; it follows from periodicity, $A_S-B_S=4D_S$ being exactly the period of
$J(a\mid\cdot)$. The primality itself is not formalised and cannot presently be: mathlib's only
route to a large prime is \texttt{lucas\_primality}, which needs a complete factorisation of
$n-1$, and that is precisely what the $0$ of $34$ above rules out. Note the two counts answer different questions and both are correct: $81$ is the
number of prime pairs, $34$ the number of \emph{immune} families among them. A complete kill/immune split of the residual is in any case beyond feasible
computation, and the honest state of the single--tail sector is \emph{$46{,}483$ empty
\textup{[P]}, $34$ immune and now unconditionally \textup{[P]} by the \textup{APRCL} run above,
$3{,}478$ factoring--hard}.
\end{proof}

\begin{proposition}[Sector decomposition of level $60$]\label{prop:sectors}
Let $U$ be a $60$--element prime support with $T(U)>2$. Then exactly one of the following holds:
\textup{(i)} $T(U)-2\ge 1/\max U$, and $U=S\cup\{q\}$ for an admissible base $S$
\textup{(}one of the $49{,}961$ of Proposition~\ref{prop:close59}\textup{)}; the
\emph{single--tail sector}, or \textup{(ii)} $T(U)-2<1/\max U$, and, writing $p>q$ for the two
largest elements and $W$ for the rest, $T(W\cup\{p\})<2$ and $T(W\cup\{q\})<2$; the
\emph{pair sector}, whose members escape every single--tail family.
\end{proposition}
\begin{proof}
$U$ lies in the tail family of $U\setminus\{r\}$ iff $T(U)-1/r\ge2$, which is easiest for
$r=\max U$; (i) is that case ($T=2$ being impossible by rigidity, the base is admissible). In
case (ii) no removal reaches $2$, and a fortiori the two stated $59$--subsets stay below $2$.
\end{proof}

\begin{remark}[The campaign; completeness of the certificate; the arity barrier]\label{rem:campaign}
The reciprocity certificate captures \emph{every} $q$--independent congruence obstruction of a
tail family: at a prime $\ell\mid A_SB_S$ the system degenerates to the rank--one congruences
$x^2\equiv D_S$ or $y^2\equiv D_S\pmod\ell$, while at $\ell\nmid 2A_SB_S$ the conditions define
conics in $(x,q)$ or $(y,q)$ with $\sim\ell$ points, deleting at most half of the residue classes
of $q$ and never all (the $2$--adic layer pins $q\equiv5\bmod8$ without emptiness); by CRT and
Dirichlet, admissible $q$ survive any finite set of such conditions. Hence a tail family admits a
congruence kill iff $(D_S\mid\ell)=-1$ for some $\ell\mid A_SB_S$, which the Jacobi symbol
detects for $31{,}219$ bases, and partial factorisation (\texttt{hunt/survivor\_kill.rs}) for a
further $15{,}264$: in total $46{,}483$ of the $49{,}961$ single--tail families at level $60$ are
proven empty ($93.0\%$), $3{,}478$ remaining open (each a balanced ${\sim}57$--digit semiprime, factoring--hard). The weapon is intrinsically \emph{arity--one}, and provably fails on the pair sector. With
$s=p+q$, $t=pq$, a pair--sector member requires the two simultaneous squares $x^2=A_Wt+D_Ws$ and
$y^2=B_Wt+D_Ws$ (the would--be third condition $s^2-4t=\square$ is automatic, being $(p-q)^2$),
together with the reciprocity--linked constraints $(D_Wp\mid q)=(D_Wq\mid p)=1$ (mod $q$, a
solution with $q\mid a$ has $N_U\equiv D_Wp$, forcing $D_Wp$ to be a square). At a prime
$\ell\mid A_W$ the condition $x^2\equiv D_Ws\pmod\ell$ constrains only $s$, which the free second
prime absorbs; the pair system is in fact locally soluble at \emph{every} modulus (a theorem,
Proposition~\ref{prop:pairlocal} below, corroborated by a direct scan over all $\ell\le600$) so,
unlike the single--tail family, the pair sector admits \emph{no} congruence obstruction at all. It
is a genuine two--parameter surface per base, $\#307$ in miniature (Remark~\ref{rem:oneside});
reciprocity cannot touch it, and only a direct both--squares search or a descent argument can.
\end{remark}

\begin{proposition}[The pair sector admits no congruence kill]\label{prop:pairlocal}
Let $W$ be any pair--sector base and $A=A_W$, $B=B_W$, $D=D_W$, so that a pair--sector solution
requires $x^2=A\,\alpha\beta+D(\alpha+\beta)$ and $y^2=B\,\alpha\beta+D(\alpha+\beta)$ with
$\alpha,\beta$ the two new primes. For every prime power $\ell^j$ this system has a solution in
units $\alpha,\beta$ and integers $x,y$ at a nonsingular point. Consequently, by CRT and Dirichlet,
no finite system of congruence conditions empties any pair--sector family: the arity barrier of
Remark~\ref{rem:campaign} is a theorem, and the certificate programme provably ends at arity one.

\textup{Formalised, in part.} \texttt{Erdos307.no\_pinning\_isUnit} is the step that makes the
tail--kill mechanism unavailable, \texttt{collapse\_identity} is regime~\textup{(2)}'s
observation that the two square conditions are one, \texttt{regime\_one\_first} and
\texttt{regime\_one\_second} are the explicit unit of regime~\textup{(1)}, and
\texttt{targets\_agree\_mod\_eight} is the mod--$8$ collapse of regime~\textup{(3)}. The lift from
$\ell$ to $\ell^j$ that regime~\textup{(1)} needs is \texttt{square\_lifts\_to\_prime\_powers},
proved from \texttt{hensel\_all\_powers}: if $\ell$ is an odd prime not dividing $c$ and $c$ is a
square modulo $\ell$, it is a square modulo every power of $\ell$. Regime~\textup{(2)} carried a citation to Weil's
bound in earlier revisions; it does not need one. Its two steps are
\texttt{Erdos307.pairlocal\_inner\_sum}, the evaluation of the inner sum of $S_{fg}$ as
$-\chi\bigl((A\alpha+D)(B\alpha+D)\bigr)$ by the quadratic complete sum of
Proposition~\ref{prop:localcomplete}, with its side condition discharged as $4D^2\alpha^2\ne0$; and
\texttt{pairlocal\_count\_pos}, the arithmetic turning the resulting $O(\ell)$ bounds into a positive
count from $\ell\ge17$. $0$ \texttt{sorry} \textup{(\texttt{lean/Erdos307/PairLocal.lean},
\texttt{lean/Erdos307/PairLocalCount.lean})}. Checked directly: over $872$ random bases with
$A-B=4D$ and every odd $\ell\le400$, the count of admissible $(\alpha,\beta)$ was positive in every
case and tracks $(\ell-1)^2/4$; the crude bounds $|S_f|,|S_g|\le2\ell$, $|S_{fg}|\le3\ell$ held in
all $409$ cases tested \textup{(\texttt{code/pairlocal\_count.gp})}.
\end{proposition}
\begin{proof}
First, no $\ell$ pins either target to a constant: that needs $\ell\mid A$ (or $B$) \emph{and}
$\ell\mid D$, whence $\ell\mid N$, contradicting $\gcd(N,D)=1$ (rigidity); so the tail--kill
mechanism is unavailable, and solubility is decided in three regimes.

\emph{(1) $\ell$ odd, $\ell\mid D$.} This covers every odd $\ell\le103$: at level $60$, omitting a
prime $p$ caps $T(U)$ at $T_{61}-1/p$, and $1/p>T_{61}-2=0.00944\ldots$ exactly for $p\le103$, so
all such primes lie in $U$, hence in $W$ (the two adjoined primes are the largest). Take $\beta=1$,
$x=1$, $\alpha=(1-D)/(A+D)\bmod\ell^j$, a unit since $A+D\equiv N\not\equiv0\pmod\ell$: the first
equation holds by construction, and the second target $(B+D)\alpha+D\equiv N\cdot N^{-1}=1\pmod\ell$
is a unit congruent to a square, so $y$ exists modulo $\ell^j$.

\emph{(2) $\ell$ odd, $\ell\nmid2D$} (in particular every $\ell\ge107$ outside $W$). The two
conditions are secretly \emph{one}: setting $m=(x^2-y^2)/4D$ and $u=(x^2-Am)/D$, the relation
$A-B=4D$ gives the polynomial identity $Bm+Du=y^2$. Only \emph{existence} is wanted, and counting the
pairs $(\alpha,\beta)$ directly settles it with no curve and no Weil bound. Write $\chi$ for the
quadratic character, $f=A\alpha\beta+D(\alpha+\beta)$ and $g=B\alpha\beta+D(\alpha+\beta)$. Summing
$\tfrac14(1+\chi(f))(1+\chi(g))$ over $\alpha,\beta\in\mathbb F_\ell^{\times}$ and discarding the at
most $2(\ell-1)$ pairs with $fg=0$,
\[ 4\,\#\{\text{both nonzero squares}\}\;\ge\;(\ell-1)^2+S_f+S_g+S_{fg}-8(\ell-1). \]
For fixed $\alpha$ each target is \emph{linear} in $\beta$, so the inner sums of $S_f$ and $S_g$
vanish by the linear evaluation above, leaving $|S_f|,|S_g|\le2\ell$. And $fg$ is a \emph{product of
two linear forms} in $\beta$, so its inner sum is the quadratic evaluation, whose side condition
$ad\ne bc$ is here $D\alpha(A-B)\alpha=4D^2\alpha^2\ne0$ --- exactly the hypotheses
$\ell\nmid2D$ and $\alpha\ne0$; summing the resulting $-\chi\bigl((A\alpha+D)(B\alpha+D)\bigr)$ over
$\alpha$ is that same evaluation again, so $|S_{fg}|\le3\ell$. Hence the count is positive as soon as
$(\ell-1)^2-7\ell-8(\ell-1)>0$, i.e.\ from $\ell\ge17$, far below the $\ell\ge107$ in play. The point
is nonsingular ($\alpha\beta\ne0$; the Jacobian contains $\operatorname{diag}(2x,2y)$), so it lifts
to every $\ell^j$.

\emph{(3) $\ell=2$.} Both targets are odd and agree modulo $8$ (as $\alpha\beta$ is odd,
$D(\alpha+\beta)\equiv0\bmod4$, and $(A-B)\alpha\beta=4D\alpha\beta\equiv0\bmod8$), and a finite
table over $(A\bmod8,\,D\bmod8)$ shows that for every base some achievable pair
$(\alpha\beta,\alpha+\beta)\bmod8$ makes the common target $\equiv1\pmod8$; a $2$--adic square,
soluble modulo $2^j$ for every $j$.

All three regimes, the collapse identity, the threshold $1/103>T_{61}-2>1/107$, and the mod--$32$
layer were verified computationally on three bases and every $\ell\le600$
(\texttt{code/pairlocal.py}; the measured density of admissible $(x,y)$ in regime~(2) is
$0.497$--$0.503$ against the predicted $\tfrac12$).
\end{proof}

\begin{remark}[Anatomy of the certificate; the empirical $5/8$ law]\label{rem:fiveeighths}
Writing $D=2D_1$, the reciprocity signs in the certificate cancel
\textup{(}$\tfrac{A-1}2+\tfrac{N-1}2=N-1+2D_1$ is even\textup{)}, leaving the clean product
\[ \left(\tfrac{D}{A}\right)\;=\;\left(\tfrac{2}{A}\right)\left(\tfrac{D_1}{N}\right),\qquad
\left(\tfrac{2}{A}\right)=+1\iff N\equiv3,5\ (\mathrm{mod}\ 8), \]
and one layer down $(N\mid D_1)=\prod_{p\mid D_1}\bigl(\tfrac{D/p}{p}\bigr)$, since
$N\equiv D/p\pmod p$: the cofactor structure of rigidity recurs inside its own certificate. This product evaluates in closed form: with $D/p=2D_1/p$ and quadratic
reciprocity,
\[ (N\mid D_1)=(2\mid D_1)\!\!\prod_{\{p,q\}\subset D_1}\!\!(-1)^{\frac{p-1}2\frac{q-1}2}
=(-1)^{\,s+\binom{t}{2}},\quad s=\#\{p\mid D_1:p\equiv3,5\ (8)\},\ \ t=\#\{p\mid D_1:p\equiv3\ (4)\}, \]
a R\'edei genus character of $D_1$, \emph{independent of $N$}. This is forced, not merely observed:
$(N\mid D_1)=\prod_{p\mid D_1}(N\mid p)$, and cofactor survival gives $N\equiv D/p=2D_1/p\pmod p$, so
$(N\mid p)=(2\mid p)\,(D_1/p\mid p)$. The first factors multiply to $(-1)^{s}$ since $(2\mid p)=-1$
exactly for $p\equiv3,5\pmod 8$; and, each unordered pair $\{p,q\}$ occurring once in each of the two
cofactors,
\[ \prod_{p\mid D_1}\Bigl(\tfrac{D_1/p}{p}\Bigr)=\prod_{\{p,q\}\subset D_1}\Bigl(\tfrac qp\Bigr)
\Bigl(\tfrac pq\Bigr)=\prod_{\{p,q\}\subset D_1}(-1)^{\frac{p-1}2\frac{q-1}2}=(-1)^{\binom t2}, \]
by quadratic reciprocity, since $\tfrac{p-1}2$ is odd exactly for $p\equiv3\pmod4$. (Checked
independently, prime by prime, on $20{,}000$ bases: $0$ violations.) The inner symbol is thus not a coin but a genus invariant; its empirical fairness
($49.9\%$, $50.0\%$ over the $49{,}961$ bases) is that character's equidistribution; and
$(D_1\mid N)$ follows by the reciprocity sign $(-1)^{\frac{D_1-1}2\frac{N-1}2}$, hence depends on $N$
only through $N\bmod4$. In fact the entire certificate closes into a finite functional. Let
$v(S)=(n_1,n_3,n_5,n_7)\bmod4$ count the odd primes of $S$ in the residue classes $1,3,5,7$ modulo
$8$. Every factor above is a function of $v$ alone: $s\equiv n_3+n_5$, $t=n_3+n_7$ (with
$\binom t2$ needing only $t\bmod4$), $D_1\equiv3^{n_3}5^{n_5}7^{n_7}\pmod8$,
$S(D_1)\equiv(n_1+n_5)-(n_3+n_7)\pmod4$, and (by the mod--$8$ law of
Proposition~\ref{C-prop:mod8}, $D_1'\equiv D_1S(D_1)\pmod8$) also
$N=D_1+2D_1'\equiv D_1+2\bigl(D_1S(D_1)\bmod4\bigr)\pmod8$. Hence
\[ (D\mid A)\;=\;K\bigl(v(S)\bigr): \]
\emph{the kill is decided by sixteen residue counts modulo $4$.} The realized vectors are exactly
the coset $n_1+n_3+n_5+n_7\equiv58\equiv2\pmod4$ (all $64$ of its classes occur), $K$ is constant
on each, and the functional reproduces the direct $113$--digit Jacobi computation on every one of
the $49{,}961$ bases (\texttt{code/kill\_classvector.py}). The five--eighths law is then an exact
finite statement, and its mechanism is the \emph{reciprocity switch}: for $N\equiv1\pmod4$ the
reciprocity sign is inert and the kill reduces to the genus coin; exact conditional rate
$12{,}699/24{,}975=0.5085$; while for $N\equiv3\pmod4$ the sign couples the parity of $t$ into
the product and the rate jumps to $18{,}520/24{,}986=0.7412$ (by class modulo $8$:
$0.509,\,0.741,\,0.508,\,0.741$ at $N\equiv1,3,5,7$; the $3/4$ classes are $N\equiv3\pmod4$, not
the $(2\mid A)$ classes). The blend $\tfrac12\cdot\tfrac{24975}{49961}+\tfrac34\cdot\tfrac{24986}{49961}
=0.62503$ against the exact $31{,}219/49{,}961=0.62487$ is the whole of the ``coincidence''. The
correlation is an ensemble property, a random--subset proxy gives ${\approx}0.45/0.55$; and
the functional yields no second certificate: it recomputes the same symbol.
\end{remark}

\begin{remark}[The kill--program is one--sided]\label{rem:oneside}
Every kill is a proof that no solution hides in that family, but if \#307 has a solution, the
solution's own family is unkillable, so no accumulation of kills can overshoot the truth: the
program's asymptote \emph{is} the existence question. Together with the self--similarity of
Remark~\ref{rem:tier60} (each rung past $60$ contains the whole problem), this locates the
campaign precisely: it corners where a solution could live, one certified region at a time, and
can terminate only if \#307 has answer \textsc{no}, which the heuristics disfavour.
\end{remark}

\begin{remark}[The reciprocity kills are the local factors of the expected count]\label{rem:killcount}
The divergent naive count of Section~\ref{C-sec:anatomy} and the reciprocity sieve of
Proposition~\ref{prop:tailkill} are the \emph{same} computation. Decompose the expected number of
solutions over tail families $S\cup\{q\}$: the summand for a base $S$ is
$\kappa_S\,\bigl(D_S\sqrt{\mathrm{mass}^2-4}\bigr)^{-1}\sum_q q^{-1}$, where
$\kappa_S=\prod_\ell(\text{local solubility density at }\ell)$ is exactly the reciprocity data.
A reciprocity--killed base carries an inert $\ell\mid A_S$ with $(D_S\mid\ell)=-1$, so
$N_S+2D_S\equiv D_S\pmod\ell$ is a non--residue for \emph{every} $q$, forcing $\kappa_S=0$; the
family drops from the count entirely. An immune base ($A_S,B_S$ both prime, both residues) has no
inert prime, so $\kappa_S>0$ and its summand is a divergent $\sum_q q^{-1}$ (verified on an explicit
$114$--digit immune base: local density $1$ at every tested prime, partial sums growing without
bound). Hence the campaign of Remark~\ref{rem:campaign} is a Selberg--sieve evaluation of the
singular series: the kills remove ${\sim}93\%$ of the count's mass, and the residual
$\log$--divergence (the entire ``weakly favours \textsc{yes}'' signal) is carried by the $34$
conditionally immune families together with the pair sector. (Conditionally, because their immunity
rests on the Baillie--PSW primality of $A_S,B_S$; the pair sector's contribution, by contrast, is
unconditional, being a theorem, Proposition~\ref{prop:pairlocal}.) This does not touch the reliability of the
count (which ignores the $N(r)\le1$ rigidity), but it pins down \emph{where} the heuristic's
\textsc{yes} lives.
\end{remark}

\begin{proposition}[Local completeness of the reciprocity sieve]\label{prop:localcomplete}
For an admissible base $S$ and an odd prime $\ell$, the admissible tail residues
$Q_\ell=\{q\bmod\ell:(A_Sq+D_S\mid\ell)\ne-1\text{ and }(B_Sq+D_S\mid\ell)\ne-1\}$ fall into exactly three
regimes. \textup{(i)} $\ell\mid D_S$, which covers every $\ell\le167$ uniformly over the $49{,}961$
admissible bases, since $1/p\ge(T_{59}-2)+\tfrac1{281}$ forces every prime $p\le167$ into every base:
then $A_S\equiv B_S\equiv N_S\not\equiv0\pmod\ell$, the two conditions collapse to the single condition
$(N_Sq\mid\ell)\ne-1$, and $|Q_\ell|=(\ell+1)/2\ni0$. \textup{(ii)} $\ell\mid A_SB_S$: one condition is
the constant $(D_S\mid\ell)$, so $Q_\ell=\emptyset$ iff $(D_S\mid\ell)=-1$ (exactly the certificate of
Proposition~\ref{prop:tailkill}) and otherwise $|Q_\ell|\ge(\ell-1)/2$. \textup{(iii)}
$\ell\nmid A_SB_SD_S$ (hence $\ell\ge173$): the complete character sums give
$|Q_\ell|\ge(\ell-3)/4-2\ge40$ (the linear sums vanish, the quadratic sum is $-\chi(A_SB_S)$, boundary
terms $\le2$). \textup{(}Regimes~\textup{(i)} and~\textup{(ii)} are formalised:
\texttt{Erdos307.regime\_one\_collapse}, \texttt{regime\_one\_same\_argument} and
\texttt{regime\_two\_constant} \textup{(}\texttt{lean/Erdos307/LocalComplete.lean}\textup{)}.
Regime~\textup{(iii)}'s counting step is \texttt{regime\_three\_count}, and its two complete
character--sum hypotheses are now theorems rather than citations:
$\sum_q\chi(aq+b)=0$ for $a\ne0$ is \texttt{Erdos307.sum\_quadraticChar\_affine}, the reindexing
$q\mapsto aq+b$ followed by \texttt{quadraticChar\_sum\_zero}; and
$\sum_q\chi((aq+b)(cq+d))=-\chi(ac)$ for $ad\ne bc$ is
\texttt{Erdos307.sum\_quadraticChar\_quadratic}, proved by factoring out $\chi(ac)$, translating to
$\sum_u\chi\bigl(u(u+e)\bigr)$ with $e\ne0$, writing $u(u+e)=u^2(1+e/u)$ for $u\ne0$ so that the
character value is $\chi(1+e/u)$, and reindexing by the involution $u\mapsto e/u$ of the nonzero
elements to reach $\sum_{t\ne0}\chi(1+t)=0-\chi(1)=-1$. Both hold over any finite field of odd
characteristic. $0$ \texttt{sorry}
\textup{(\texttt{lean/Erdos307/CompleteSums.lean})}. Checked against direct evaluation for all odd
primes $\ell\le59$: the linear sum over every $(a,b)$ with $a\ne0$, and the quadratic sum over all
$33{,}008{,}952$ quadruples with $ad\ne bc$ --- $0$ violations
\textup{(\texttt{code/complete\_sums\_check.gp})}.\textup{)}
Consequently $Q_\ell\ne\emptyset$ forces $|Q_\ell|\ge2$ at every modulus, the surviving
local density is positive (e.g.\ $\prod_{\ell\le300}|Q_\ell|/\ell\approx2.2\times10^{-19}>0$ on the
canonical base), and by CRT \emph{no finite system of congruences (no covering) can annihilate an
unkilled family}: the binary certificate of Proposition~\ref{prop:tailkill} is the complete local theory
of the tail sector, there is no ``second kill layer,'' and the one--sidedness of
Remark~\ref{rem:oneside} is structural rather than an artefact of method. (All three regimes verified
exhaustively on the canonical base: collapse and $|Q_\ell|=(\ell+1)/2$ at
$\ell\in\{3,5,7,101,277\}$; the Weil floor at every prime $281\le\ell\le600$; $A_S,B_S$ carry no factor
below $10^6$, consistent with the balanced--semiprime wall of Remark~\ref{rem:campaign}.)
\end{proposition}

\begin{proposition}[The Pythagorean tail prime is bounded]\label{prop:tailbound}
Let $S$ be a finite set of primes, $D=\prod S$, $N=\sum_{p\in S}D/p$, $A=N+2D$, $B=N-2D$, and suppose
$D<N$. Let $q$ be a prime with $q\nmid 2D$ for which $Aq+D$ and $Bq+D$ are both perfect squares. Then
there are integers $m,c$ with $mc=4D$ and
\[ q^2m^2-4Nq+c^2-4D=0, \qquad \bigl(qm^2-2N\bigr)^2=4\bigl(AB+Dm^2\bigr), \]
so that $q=2(N\pm k)/m^2$ with $k^2=AB+Dm^2$ and $m\mid 4D$. Moreover $m$ is \emph{even}, say
$m=2\alpha$ with $\alpha\mid D$ and $\beta=D/\alpha$, and then $q$ is a root of
\[ \alpha^2q^2-Nq+(\beta^2-D)=0, \qquad\text{so}\qquad
   q=\frac{N+\sqrt{N^2-4D^2+4\alpha^2D}}{2\alpha^2}, \]
a quantity strictly decreasing in $\alpha$. Hence the maximum is at $\alpha=1$ and
\[ q\ \le\ \tfrac12\Bigl(N+\sqrt{N^{2}-4D^{2}+4D}\Bigr)\ =\ tD+O(1),
   \qquad t=\tfrac12\bigl(\sigma+\sqrt{\sigma^{2}-4}\bigr),\ \ \sigma=\sigma(S)=N/D. \]
The largest admissible tail is therefore the base times the very ratio $b/a$ that the mass forces.
At level $59$ this is $9.207\times10^{112}=1.0497\,D$, a factor $7.63$ below $4N$.
\end{proposition}

\begin{proof}
Write $x^2=Aq+D$ and $y^2=Bq+D$ with $x,y\ge0$. Since $A-B=4D$ \emph{exactly},
$x^2-y^2=(A-B)q=4Dq$, that is $(x+y)(x-y)=4Dq$. The prime $q$ divides exactly one of the two factors:
were it to divide both it would divide $2x$ and $2y$, hence $x$ and $y$ as $q$ is odd, hence
$q^2\mid 4Dq$ and $q\mid 4D$, against $q\nmid2D$. Write the factor divisible by $q$ as $qm$ and the
other as $c$, so that $mc=4D$ and, in either case, $2x=qm+c$. Multiplying by $m$,
\[ 2xm\;=\;qm^2+mc\;=\;qm^2+4D, \]
and squaring against $x^2=(N+2D)q+D$ gives
$4m^2\bigl((N+2D)q+D\bigr)=q^2m^4+8Dqm^2+16D^2$. Substituting $16D^2=m^2c^2$ and cancelling $m^2>0$
leaves $q^2m^2-4Nq+c^2-4D=0$; completing the square and using $AB=N^2-4D^2$ gives the second display.
For the parity, $x+y$ and $x-y$ differ by $2y$ and so have equal parity; were both odd their
product would be odd, whereas it is $4Dq$. Hence both are even, and as $q$ is odd, $m$ is even.
Finally $m^2\ge4$ and $c^2\ge0$ give $4Nq=q^2m^2+c^2-4D\ge 4q^2-4D$, that is $Nq\ge q^2-D$; if
$q\ge N+1$ then $q^2-Nq-D\ge q-D>0$ because $D<N<q$, a contradiction, so $q\le N$. (Using only
$m\ge1$ here gives $q\le4N$; the parity removes the factor four.) For the sharp form, $m=2\alpha$ and
$c=2\beta$ with $\alpha\beta=D$ turn $q^2m^2-4Nq+c^2-4D=0$ into $\alpha^2q^2-Nq+\beta^2-D=0$, whose
larger root is $\bigl(N+g(u)\bigr)/(2u)$ with $u=\alpha^2$ and $g(u)^2=N^2-4D^2+4uD$. Its derivative
in $u$ has the sign of $2uD-g(N+g)$, and $4uD=g^2-N^2+4D^2$ turns that into
$-\bigl((g+N)^2-4D^2\bigr)/2$, negative since $N>2D$. So the root decreases in $\alpha$ and is
largest at $\alpha=1$.
\end{proof}

\begin{proposition}[What a barrier improvement must supply]\label{prop:barthreshold}
Let $(a,b)$ be a solution with $a<b$ and put $t=b/a$, so that $t=\sigma(a)$ and, by
$\sigma(a)\sigma(b)=1$,
\[ \sigma(U)\;=\;\sigma(a)+\sigma(b)\;=\;t+\tfrac1t. \]
For $n$ with $T_n>2$ let $t_n>1$ be the root of $t+1/t=T_n$. Since $\sigma(U)\le T_{|U|}$,
\[ \sigma(a)>t_n\ \Longrightarrow\ |U|\ \ge\ n+1, \]
and this is the only route by which a mass argument can raise the barrier. Numerically
\[
\begin{array}{l|cccccc}
n & 59 & 60 & 61 & 62 & 63 & 64\\\hline
T_n & 2.0023502 & 2.0059089 & 2.0094424 & 2.0128554 & 2.0161127 & 2.0193282\\
t_n & 1.0496677 & 1.0798804 & 1.1020081 & 1.1199914 & 1.1352477 & 1.1490254
\end{array}
\]
\end{proposition}

So the barrier level is decided entirely by how close $\sigma(a)$ may lie to $1$, and each further
level demands a definite gap: $|U|\ge61$ needs $\sigma(a)>1.0799$, a full $8\%$ above $1$.

That is why the mass mechanism is exhausted here rather than merely difficult. The only
unconditional lower bound on the gap is $\sigma(a)-1=(a'-a)/a\ge1/a$, which is worthless at
$a>10^{56}$; and $\sigma$ genuinely comes close to $1$ from above on very few primes, the
Sylvester-type set $\{2,3,7,41\}$ giving $\sigma=1.00058\ldots$ with four. Nothing in the cycle
equations forbids $a'-a=1$. A barrier past $60$ must therefore come from a mechanism that bounds
$\sigma(a)$ away from $1$ without appealing to the size of $U$ --- which is what
Theorem~\ref{thm:barrier} does in the opposite direction --- and we know of none.

The same reduction locates the extreme case. Writing $c=b-a$, the identity
$\sigma(U)-2=(t-1)^2/t=c^2/D$ is the minus layer $D'-2D=c^2$, and $\gcd(c,D)=\gcd(b-a,ab)=1$ because
$\gcd(a,b)=1$. Since an admissible support contains every prime $\le167$
\textup{(}Remark~\ref{rem:tier60}\textup{)}, $c$ is divisible by none of them, so $c=1$ or
$c\ge173$. The branch $c=1$ is $D'=2D+1$, equivalently $a'=a+1$ and $(a+1)'=a$; it is the closest a
solution could come to $\sigma(a)=1$, and it is exactly the configuration a mass argument cannot
reach. Below $3\times10^{6}$ the equation $a'=a+1$ has only the solutions $30$, $858$, $1722$ and
$66{,}198$, and none extends: for the first three $a+1$ is prime, so $(a+1)'=1$, and for the last
$a+1$ is not squarefree. The branch is not excluded, only unsearched, since the barrier puts it above
$10^{56}$. The level--$59$
computation that gives $|U|\ge60$ is, in these terms, an exhaustive verification that
$\sigma(a)>t_{59}$, and Proposition~\ref{prop:levelbarrier} says the same verification is unavailable
one level up.

\begin{remark}[The Lehmer reading of the tail is withdrawn]\label{rem:notlehmer}
Remark~\ref{rem:lehmer} below reads the surviving question as the primality of a term $q_n$ of the
\emph{exponential} Pell orbit, and Remark~\ref{C-rem:whatisleft} names that orbit as the only live
branch. Proposition~\ref{prop:tailbound} withdraws the reading. The orbit is genuinely infinite and
both square conditions do hold along all of it; what fails is that no term past $N$ can be prime.
The single identity responsible is $A-B=4D$, which the earlier treatment used only to form the Pell
system and never to factor $4Dq$.

Three consequences, in increasing order of what they cost.

\emph{The gate is finite, not Lehmer.} The tail primes of a base lie in the explicit finite set
$\{2(N\pm k)/m^2: m\mid 4D,\ k^2=AB+Dm^2\}$. So the second of the three gates is a search over the
divisors of $4D$, not a primality question along an exponential sequence, and the classification of
the residual difficulty as Lehmer--class rather than Schinzel--class does not survive. Checked sound
\emph{and complete} against exhaustive search on $1200$ bases with zero mismatches, and stepped out
along a real orbit, in \texttt{code/tailbound.gp}: on the base $D=1$, $N=7$ (bound $N=7$) the orbit
gives $q_1=7$ prime, then $q_2=741895$, $q_3=76921173511$, and every later term composite, with both
square conditions holding throughout.

\emph{The count reverses sign, and the search has now been run.} The expected number of admissible
$m$ is $\sum_{m\mid 4D}(AB+Dm^2)^{-1/2}\approx D^{-1/2}\sum_{m\mid4D}m^{-1}$, which on the first
immune base is $2.4\times10^{-56}$ and below $10^{-38}$ summed over all $34$ immune families, the
opposite of the ``weakly favours \textsc{yes}'' of Remark~\ref{rem:killcount} on exactly the families
that reading was about. That prediction is now checked. The full divisor space is
$4\cdot2^{58}\approx1.15\times10^{18}$ per family and is not enumerable, but because the hit
probability falls like $1/(m\sqrt D)$ the expected count is carried by small $m$: restricting to
$\omega(m)\le8$ leaves $99.9998\%$ of the expected count. Over all $34$ families that slice is
$3.0777\times10^{11}$ candidates, and
\[ \textup{it contains no }m\textup{ with }AB+Dm^2\textup{ a perfect square} \]
($4{,}933$ survive the modular cascade, exact \texttt{issquare} kills all of them; the earlier
$\omega(m)\le7$ pass, $4.7081\times10^{10}$ candidates and $771$ survivors, agrees).

The same search has been run over \emph{every} admissible base, which matters because
Proposition~\ref{prop:tailbound} requires no primality of $A_S$ or $B_S$, only the algebra. The $34$
immune families are the ones left after the Jacobi kill and the factoring attack, and their immunity
is conditional on a Baillie--PSW verdict; the other $3{,}444$ of the $3{,}478$ open families are open
precisely because $A_S$ or $B_S$ is a balanced semiprime nobody can split. The tail search does not
care. Over all $49{,}961$ admissible bases at $\omega(m)\le4$, which is $99.2\%$ of the expected count
per base, the run is $9.1296\times10^{10}$ candidates, $1{,}469$ survive the modular cascade, and
exact \texttt{issquare} kills all $1{,}469$. Every survivor was additionally checked to satisfy
$m\mid 4D$, the hypothesis the proposition needs. So
\[ \textup{no admissible base admits a Pythagorean tail prime with }\omega(m)\le4, \]
unconditionally and with no primality input. Where Proposition~\ref{prop:tailkill} leaves $3{,}478$
of the one--new--prime families at level $60$ open, of which $34$ only conditionally immune, this
leaves none open on the stated slice. It is a slice result and not a proof: the untested remainder is
$0.8\%$ of the expected count per base, and $\omega(m)\ge5$ is untested.
The run is \texttt{code/tailsearch.rs} (a cascade of quadratic--residue filters modulo primes coprime
to $2D$, $771$ survivors) with every survivor then tested exactly by \texttt{issquare} in
\texttt{code/tailsearch\_verify.gp} ($0$ squares). A positive control recovers the known solutions of
the toy base $D=1$, $N=7$, and an exhaustive exact run at $\omega(m)\le2$ agrees with the filter, so
the cascade is not silently over--rejecting.

It is worth recording what the slice covers in the other measure, since the two are thirteen orders
apart. The parametrisation runs over $m\mid4D$ with $m$ even, so $m=2\alpha$ with $\alpha\mid D$, and
at level $59$ that is $2^{59}=5.7646\times10^{17}$ divisors. Splitting on the parity of $\alpha$, the
divisors with $\omega(m)\le4$ number $2\sum_{j\le3}\binom{58}{j}=65{,}136$, and those with
$\omega(m)\le8$ number $2\sum_{j\le7}\binom{58}{j}=692{,}376{,}800$. As fractions of the space the
bound $q\le N$ actually proves finite, these are
\[ 1.130\times10^{-13}\qquad\text{and}\qquad 1.201\times10^{-9}, \]
against the $99.2\%$ and $99.9998\%$ of the \emph{expected count} quoted above. Both figures are
correct and they measure different things: the expected count concentrates on small $\omega$, the hit
probability falling like $1/(m\sqrt D)$, so a slice can carry almost all of the expectation while
carrying almost none of the space. Only the second pair of numbers is a statement about what has
been checked. This is \textsc{verified} on the stated slice, not on the
full divisor set: the residual $0.0018\%$ of the expected count is untested.

The filter primes must be coprime to $2D$, which is itself a structural fact rather than a
convenience. Since $D$ contains every prime up to $167$, for any $\ell\mid D$
\[ AB+Dm^2-N^2=D\bigl(m^2-4D\bigr)\equiv0 \pmod{\ell}, \]
so the value is the square $N^2$ modulo $\ell$ for \emph{every} $m$, and as $\gcd(D,N)=1$ Hensel
lifts that root to every power of $\ell$. The tail search therefore has no local obstruction at any
prime dividing $D$, the analogue of Proposition~\ref{prop:localcomplete} for this formulation; the
classical $64/63/65/11$ pre--filter for squares rejects nothing here, measured at a $75\%$ pass rate
modulo $64$ and $100\%$ modulo $63$. Formalised as \texttt{tail\_value\_mod\_dvd} and
\texttt{hensel\_lift\_identity}; the induction from the step to all powers is routine and is not
formalised.

\emph{The paper's own height estimate now says the same thing.} Remark~\ref{rem:lehmer} puts the
least point of a nonempty fiber at height $\exp\bigl(10^{112}\bigr)$ on regulator grounds, so its
$q\approx x^2/A$ carries some $10^{112}$ digits, against a bound of $115$. Two independent estimates,
one from divisor counting and one from the class--number formula, agree that no immune family admits
a Pythagorean tail prime. Neither is a proof: the bound $q\le N$ is proved, the emptiness is
\textsc{heuristic}.

Formalised in \texttt{lean/Erdos307/TailBound.lean} (\texttt{tail\_factor}, \texttt{tail\_quadratic},
\texttt{tail\_discriminant}, \texttt{tail\_bound}, \texttt{tail\_finite}) on the three standard
axioms, with a non--vacuity witness at the base above.
\end{remark}

\begin{proposition}[The level--by--level programme terminates at $60$]\label{prop:levelbarrier}
Let $U$ be a finite set of primes with $\sum_{p\in U}1/p>2$ and let $B$ be arbitrary. Then there is a
prime $q>B$ with $q\notin U$ such that $U\cup\{q\}$ again has $\sum 1/p>2$. Consequently, for every
$n>59$ there are infinitely many $n$--element prime supports of mass exceeding $2$.
\end{proposition}

\begin{proof}
The primes are infinite, so choose a prime $q$ exceeding both $B$ and $\max U$; then $q\notin U$ and
$\sum_{p\in U\cup\{q\}}1/p=\sum_{p\in U}1/p+1/q>2$. Iterating from the first $59$ primes, whose
reciprocal sum is $N_{59}/P_{59}=2.00235\ldots>2$, produces admissible supports of every cardinality
above $59$, each containing a prime beyond any prescribed bound.
\end{proof}

\begin{remark}[What the level barrier costs, and what it does not]\label{rem:levelbarrier}
Proposition~\ref{prop:close59} closes level $59$ by enumerating its $49{,}961$ admissible supports,
and Proposition~\ref{prop:tailkill} together with the search of Remark~\ref{rem:notlehmer} closes
level $60$ by adjoining one prime to each. The obvious continuation is to enumerate the admissible
$60$--element supports and repeat. Proposition~\ref{prop:levelbarrier} says there is no such
enumeration \emph{of supports}.

The finiteness at level $59$ is an accident of one inequality being tight: $59$ primes reach mass $2$
only just, at $2.00235\ldots$, so demanding mass $>2$ from exactly $59$ primes pins them near the
smallest $59$. Allow a sixtieth and the first $59$ already carry the mass alone, leaving the sixtieth
entirely free. The mass condition, having been met, constrains nothing further. The next proposition
shows that the \emph{solutions} are nevertheless finite at every level: the haystack is infinite, the
needles are not.

\begin{proposition}[Each level carries only finitely many solutions]\label{prop:levelfinite}
Fix cardinalities $m,n\ge0$ and rationals $c_1,c_2\ge0$. The equation
\[ \Bigl(c_1+\sum_{a\in A}\frac1a\Bigr)\Bigl(c_2+\sum_{b\in B}\frac1b\Bigr)=1 \]
has only finitely many solutions in sets $A,B$ of integers $\ge2$ with $|A|=m$, $|B|=n$; in
particular, for every $(m,n)$ there are only finitely many pairs of prime sets with
$\sigma(P)\sigma(Q)=1$, $|P|=m$, $|Q|=n$, and the bound is effective. $0$ \texttt{sorry}
\textup{(\texttt{Erdos307.sols\_finite}, \texttt{erdos307\_level\_finite})}.
\end{proposition}
\begin{proof}
Induct on $m+n$; the constants absorb removed elements. Write $u,v$ for the two sums and
$\varepsilon=1-c_1c_2$. Expanding, any solution satisfies $c_1v+u(c_2+v)=\varepsilon$, and the left
side is strictly positive for $m+n\ge1$ (if $m\ge1$ then $u>0$ and $c_2+v>0$, else the product is
$0$; if $m=0$ then $c_1(c_2+v)=1$ forces $c_1>0$, $v>0$), so $\varepsilon>0$ is fixed by
$(c_1,c_2)$ alone. One of the two terms is at least $\varepsilon/2$.

If $u(c_2+v)\ge\varepsilon/2$: every $b\in B$ has $1/b\le1/2$, so $v\le n/2$, whence
$u\ge\varepsilon/(2c_2+n)$; since $u$ is a sum of $m$ terms, some $a_0\in A$ has
$1/a_0\ge u/m$, that is
\[ a_0\;\le\;\frac{m\,(2c_2+n)}{\varepsilon}. \]
Moving $a_0$ into the constant, $(A\setminus\{a_0\},B)$ solves the equation with
$(m-1,n,c_1+1/a_0,c_2)$, and the map is injective for fixed $a_0$. If instead
$c_1v\ge\varepsilon/2$, then symmetrically some $b_0\le2c_1n/\varepsilon$ moves into $c_2$. Either
way the solution set injects into finitely many strictly smaller problems, and the induction closes.
The bounds compound: each step divides by an $\varepsilon$ whose denominator is at most the product
of the elements already fixed, so the resulting height bound is doubly exponential in $m+n$; crude,
but effective.
\end{proof}

\begin{corollary}[The largest element of a support is explicitly bounded]\label{cor:maxbound}
Let $U$ be the support of a two--cycle, $q=\max U$, $S=U\setminus\{q\}$, $D=\prod S$ and
$N=\sum_{p\in S}D/p$. Then
\[ N<2D\ \Longrightarrow\ q\le\frac{D}{2D-N},
   \qquad\qquad N\ge2D\ \Longrightarrow\ q\le\tfrac12\Bigl(N+\sqrt{N^{2}-4D^{2}+4D}\Bigr). \]
In both regimes $q$ is bounded by an explicit function of $S$ alone, so the enumeration at each
level is a finite tree whose branching is known in advance. $0$ \texttt{sorry} for the mass regime
\textup{(\texttt{Erdos307.max\_le\_of\_mass\_deficit}, \texttt{max\_le\_of\_mass\_deficit'})}.
\end{corollary}
\begin{proof}
The mass of a two--cycle support is $\ge2$ \textup{(}$\sigma(P)+\sigma(Q)\ge2\sqrt{\sigma(P)\sigma(Q)}=2$\textup{)},
so $N/D+1/q\ge2$. If $N<2D$ then $1/q\ge2-N/D>0$ and the first bound follows. If $N\ge2D$ then
$D<N$ and $q\nmid2D$, and both $Aq+D$ and $Bq+D$ are squares, being $D(U)+2\prod U=(\prod P+\prod
Q)^2$ and $D(U)-2\prod U=(\prod P-\prod Q)^2$ under $\prod U=qD$ and $D(U)=D+qN$; the second bound
is then Proposition~\ref{prop:tailbound} at $\alpha=1$, the discriminant being nonnegative because
$N\ge2D$.
\end{proof}

\begin{proposition}[The mass ladder: the $59$ smallest elements are bounded at every level]\label{prop:massladder}
Let $U$ be a set of $k$ primes with $\sum_{p\in U}1/p\ge2$, written $u_1<\dots<u_k$. Then for every
$j$ with $\sum_{i<j}1/u_i<2$,
\[ u_j\;\le\;\frac{k-j+1}{\,2-\sum_{i<j}1/u_i\,}\;\le\;\frac{k-j+1}{2-T_{j-1}} . \]
The hypothesis holds for all $j\le59$, since $T_{58}=1.99874\ldots<2$. In particular, at every level
$k\ge59$,
\[ u_{59}\;\le\;\frac{k-58}{2-T_{58}}\;=\;793.67\ldots\times(k-58), \]
so $u_1,\dots,u_{59}$ are bounded by a quantity linear in $k$; only $u_{60},\dots,u_k$ escape the
mass, and those are covered by Corollary~\ref{cor:maxbound}. $0$ \texttt{sorry}
\textup{(\texttt{Erdos307.mass\_ladder\_step}, \texttt{mass\_ladder\_step'})}.
\end{proposition}
\begin{proof}
The $k-j+1$ elements $u_j,\dots,u_k$ carry reciprocal mass at least $r_j=2-\sum_{i<j}1/u_i$, and
each is at least $u_j$, so $r_j\le(k-j+1)/u_j$, which is the first bound. For the second,
$u_i\ge p_i$ gives $\sum_{i<j}1/u_i\le T_{j-1}$, and $r_j>0$ exactly when $\sum_{i<j}1/u_i<2$,
which $T_{j-1}\le T_{58}<2$ guarantees for $j\le59$.
\end{proof}

\begin{remark}[What this gives at the first open level]\label{rem:level61}
At level $61$ the ladder reads $u_1\le30.5$, $u_{10}\le103.7$, $u_{30}\le201.1$, $u_{58}\le808.0$,
$u_{59}\le2381.0$: the $59$ smallest elements of any level--$61$ support lie among the first $353$
primes, so positions $1,\dots,59$ range over at most $\binom{353}{59}\approx10^{67.9}$
configurations, with $u_{60}$ and $u_{61}$ then confined by
Corollary~\ref{cor:maxbound}. This replaces the doubly exponential height of
Proposition~\ref{prop:levelfinite} at this level by an explicit and much smaller count. It is of
course still far beyond enumeration; the point is that the first open level now has a stated size
rather than an abstract finiteness.

Two consistency checks. At level $59$ the ladder gives $u_{59}\le793.67\ldots$, and
\texttt{Erdos307.elt\_le\_795} bounds the elements of a $59$--element support by $795$ by an
unrelated argument; the same value appears again as the mass regime of
Corollary~\ref{cor:maxbound} at the base of the first $58$ primes. Falsification: over $5{,}781$
supports of mass $\ge2$ generated by perturbing initial segments, $0$ violations of the ladder, and
over $5{,}885$ supports $0$ violations of $u_{59}\le(k-58)/(2-T_{58})$
\textup{(\texttt{code/mass\_ladder.gp})}.
\end{remark}

\begin{corollary}[An explicit height at the first open level]\label{cor:level61height}
Every solution of \#307 at level $61$ has $\prod(P\cup Q)<10^{567}$.
\end{corollary}
\begin{proof}
Write $U=S\cup\{u_{60},u_{61}\}$ with $S$ the $59$ smallest elements, $D=\prod S$.

If $\sigma(S)\ge2$ then $S$ is one of the $49{,}961$ admissible $59$--element supports of
Corollary~\ref{C-cor:count59}, whose products satisfy $\log_{10}D\le113.397$
\textup{(\texttt{code/maxprod59.rs})}. Splitting $S=P_0\sqcup Q_0$ and writing $M=\prod P_0$,
$N=\prod Q_0$, $M'=D(P_0)$, $N'=D(Q_0)$, Theorem~\ref{thm:frame} determines the two remaining primes:
$u=(MN'+N^2)/\Delta$ and $v=(M'N+M^2)/\Delta$ with $\Delta=MN-M'N'$. Here
$\Delta=D\bigl(1-\sigma(P_0)\sigma(Q_0)\bigr)$ is a positive integer, since
$\sigma(P_0)\sigma(Q_0)<\sigma(P)\sigma(Q)=1$, so $\Delta\ge1$; and $MN=D$ exactly, so
$MN'=D\,\sigma(Q_0)\le tD$ with $t=1.10200\ldots$ the window root at $T_{61}$. Hence
$u,v\le D^2+tD$, giving $\log_{10}u,\log_{10}v\le226.79$ and
$\log_{10}\prod U\le113.397+2(226.79)=566.98$.

If $\sigma(S)<2$ then Proposition~\ref{prop:massladder} bounds every element of $S$ and a
branch--and--bound over the admissible chains gives $\log_{10}D\le113.526$
\textup{(\texttt{code/ladder\_maxprod.rs})}. The ladder extends one step, $u_{60}\le2/(2-\sigma(S))\le2D$
since $2-\sigma(S)\ge1/D$, and Corollary~\ref{cor:maxbound} then gives $u_{61}\le tD\,u_{60}$, so
$\log_{10}\prod U\le454.75$. The first case dominates.
\end{proof}

\begin{remark}[What the height means]\label{rem:heightmeaning}
Proposition~\ref{prop:levelfinite} makes every level finite by an argument using only the two
cardinalities, and its height compounds doubly exponentially: at level $61$ it is a tower, not a
number one can write. Corollary~\ref{cor:level61height} replaces it by $10^{567}$, because it uses
the two square conditions and the cycle equations rather than the sizes alone. Together with the
barrier $\prod(P\cup Q)\ge(2.09\times10^{56})^2=10^{112.6}$ this confines level $61$ to the window
\[ 10^{112.6}\;\le\;\prod(P\cup Q)\;<\;10^{567}. \]
Two remarks on what this is and is not. It is not an enumeration: $10^{567}$ is far past any search,
and the reduction that matters is not the height but Theorem~\ref{thm:frame}, which turns the two
unknown primes into a \emph{computed} pair once the split of the bottom $59$ is fixed, so the free
parameters at level $61$ are the support and the split, not the primes. And it is not a route to the
problem, since the level index remains unbounded; it is a statement about one level.
\end{remark}

\begin{remark}[What the window does not do]\label{rem:windownotpin}
It is worth recording what fails, since the natural guess is wrong. One might hope the mass window
pins all but the top few elements of a support, leaving a bounded search. It does not, and the
pinning weakens as the level rises: a prime $p$ is forced into a $k$--element support exactly when
$1/p>T_{k+1}-2$, which forces $39$ primes at level $59$, $27$ at level $60$ and only $21$
\textup{(}all $p\le73$\textup{)} at level $61$, leaving $40$ of the $61$ elements unconstrained by
mass. The bound of Corollary~\ref{cor:maxbound} comes from the square conditions, not from the mass,
and this is why: past level $59$ the mass condition, once met, constrains nothing further.
Consistency check: at the base of the first $58$ primes the mass regime gives $q\le793.67\ldots$,
while \texttt{Erdos307.elt\_le\_795} bounds the elements of a $59$--element support by $795$ through
an unrelated argument \textup{(\texttt{code/level\_dichotomy.gp}; $400$ bases, $0$ failures)}.
At level $61$ the second regime applies at the extremal base and gives $q\le1.1020\,D$ with
$\log_{10}D=117.84$.
\end{remark}

\begin{remark}[The level--by--level programme survives its own barrier]\label{rem:levelrestored}
Together with Proposition~\ref{prop:levelbarrier} this splits the situation cleanly: at every level
above $59$ the admissible supports are infinite, and the solutions are finite with an effective
height bound. Restricted to any fixed level, \#307 is therefore \emph{decidable}: a terminating
computation either exhibits the finitely many solutions or certifies that the level is empty. What
is genuinely unbounded in the problem is only the level itself, so \#307 is equivalent to the
halting of one explicit search, which is the precise content of the problem page's
\textup{VERIFIABLE} tag. We record the honest scale: the bound of
Proposition~\ref{prop:levelfinite} at level $61$ is astronomically beyond enumeration, so
decidability in principle does not yield a computation in practice; the value of the proposition is
structural. Validated on the inhomogeneous case $c_1=1$, where solutions are the known
$1$--containing coprime examples: the enumeration recovers $(\{5\},\{2,3\})$,
$(\{41\},\{2,3,7\})$, $(\{1805\},\{2,3,7,43\})$ and their congeners, every one containing an
element within the proven bound \textup{(\texttt{code/level\_finite\_check.gp})}.
\end{remark}

Two consequences, and the second is the one that matters. The searches of
Remark~\ref{rem:notlehmer}, however far they are pushed in $\omega(m)$, bear on level $60$ alone;
widening them is not a route upward. And any further progress on \#307 must be
\emph{non--enumerative}: exhausting supports one level at a time is not merely expensive beyond
level $60$, it is impossible, so the method behind every unconditional exclusion in this note is now
spent. That is a no--go, not a resource limit.

Formalised in \texttt{lean/Erdos307/LevelBarrier.lean} (\texttt{admissible\_extends},
\texttt{admissible\_of\_card\_add}, \texttt{level\_enumeration\_terminates}) on the three standard
axioms, with a non--vacuity witness exhibiting the first $59$ primes as a set the hypotheses apply
to.
\end{remark}
\begin{remark}[The immune signal, and where its Lehmer reading broke]\label{rem:lehmer}
Superseded by Proposition~\ref{prop:tailbound} and Remark~\ref{rem:notlehmer}: the orbit described
here is real, but the inference that primality along it is the residual question is not, because the
tail prime is bounded by $N$. What follows is the orbit, the three gates, and the structural points
that survive.

For a reciprocity--immune base the Pell system of Remark~\ref{rem:tier60},
$B_Sx^2-A_Sy^2=-4D_S^2$, is locally soluble everywhere ($\kappa_S>0$), and a Pythagorean pair in the
tail family $S\cup\{q\}$ is \emph{exactly} a solution with $q=(x^2-D_S)/A_S=(y^2-D_S)/B_S$ a
positive prime. Its solutions grow like $\varepsilon^n$ in the fundamental unit of the order of
discriminant $4A_SB_S$ ($224$ digits, $A_SB_S$ a nonsquare, on an explicit immune base). Writing
$x_{n+1}=tx_n-x_{n-1}$ with $x_n=(\alpha^n+\beta^n)/2$ and $\alpha\beta=1$, the squares $x_n^2$ have
characteristic roots $\alpha^2,\beta^2,1$ and the constant shift is absorbed by the root $1$, so
\[ q_n=\frac{x_n^2-D_S}{A_S}\qquad\text{satisfies}\qquad q_{n+3}=(t^2-1)q_{n+2}-(t^2-1)q_{n+1}+q_n, \]
a non--degenerate ternary linear recurrence (verified on four Pell parameters). A
reciprocity--killed base is the complementary degeneration: there the same system is locally
\emph{insoluble} ($\kappa_S=0$), so no candidate exists at all, and the dichotomy of
Proposition~\ref{prop:tailkill} is Pell--solubility versus Pell--insolubility.

The return is milder than ``$k=0$'' suggests: on the Pythagorean locus the deviation $k=(a'-b)/a$ is
an \emph{integer} (the identity $v(D(u)-v)=u(u-D(v))$ with $\gcd(u,v)=1$), forced into $\{0,-1\}$ by
the near--balanced immune ratio, so $k=0$ is a sign choice rather than a measure--zero coincidence,
and the naive Wieferich shortcut that would pin $q$ is vacuous, since $x^2\equiv D_S\pmod q$ holds
identically from $x^2=A_Sq+D_S$.

Neither square condition carries the difficulty alone. Imposing only the minus--square,
$d^2=B_Sq+D_S$, the admissible $d$ satisfy $d^2\equiv D_S\pmod{B_S}$, and along such a class
$d=d_0+tB_S$ one has
\[ q(t)\;=\;\frac{d_0^2-D_S}{B_S}\;+\;2d_0\,t\;+\;B_S\,t^2, \]
a quadratic in $t$ (toy bases: $3+54t+173t^2$, $7+238t+1693t^2$, $2+446t+17357t^2$): a Bunyakovsky
question, of the same class as the $p=5$ family of Proposition~\ref{C-prop:plus}\textup{(iii)}, and
conjecturally routine. It is the \emph{conjunction} of the two squares that replaces the quadratic
parameter by a fundamental unit and turns the orbit exponential.

\emph{Where the reading broke.} That exponential growth was taken to place the question in the
Lehmer class, on two grounds: no algebraic factorisation is available, $D_S$ being squarefree so
that $x^2-D_S$ is irreducible, and no covering congruence can kill the family, by
Proposition~\ref{prop:localcomplete}. Both observations are correct, and neither is the relevant
one. The factorisation that does exist is not of $x^2-D_S$ but of $x^2-y^2=4D_Sq$, available because
$A_S-B_S=4D_S$ \emph{exactly}: the identity was used here only to build the Pell system, never
multiplicatively. Proposition~\ref{prop:tailbound} draws it, and the ternary recurrence above then
supplies no prime term past $q\le N$.

Three nested gates separate an immune family from a solution: fiber nonemptiness, a tail prime, and
the $k=0$ return. The first is a class--group condition, since Hasse supplies only rational points,
and it is quantitative: immunity makes $A_SB_S$ a product of two distinct primes, hence squarefree,
so the conic lives in the near--maximal order of $\mathbb{Q}(\sqrt{A_SB_S})$ of discriminant
${\sim}A_SB_S\approx10^{224}$, and the class--number formula puts the expected regulator at
$R=\sqrt\Delta\,L(1,\chi)/(2h)\approx10^{112}$; barring an anomalously large class number, the
\emph{least} point of a nonempty fiber sits at height $\exp\bigl(10^{112}\bigr)$. That estimate is
used in Remark~\ref{rem:notlehmer} as independent corroboration that the second gate is empty. A
direct sweep to $q\le10^{12}$ was structurally incapable of reaching even the first gate.
\end{remark}

\begin{remark}[The affine sieve does not reach a rank--one orbit]\label{rem:affinesieve}
The modern instrument for primes in orbits is the affine sieve of Bourgain, Gamburd and Sarnak
\cite{BGS2010}, which produces $r$--almost--primes in orbits of thin groups. It is the one method
designed for exactly the shape of question Remark~\ref{rem:lehmer} isolates, so its failure here is
worth recording precisely rather than left implicit.

The affine sieve needs a spectral gap: the congruence quotients of the acting group must form an
expander family, and that is what supplies the level of distribution the sieve consumes. Here the
$q_n$ run along a Pell orbit, so the acting group is the cyclic group generated by the fundamental
unit of the order of discriminant $4A_SB_S$. Cyclic groups are abelian, hence amenable, and have no
spectral gap; expansion is precisely the property an abelian group lacks. The sieve therefore
returns nothing, and it returns nothing on $2^n-1$ for the same reason. The obstruction is rank, not
thinness as such: non--elementary thin orbits, Apollonian packings among them, do yield
almost--primes.

The natural repair fails for a reason worth stating too. One may embed the rank--one orbit in a
larger thin group whose orbit does expand, and sieve there. What the affine sieve then delivers is
infinitely many almost--primes in the \emph{larger} orbit, inside which our sequence has relative
density zero; the sieve does not localise to a prescribed sub--orbit. Expansion is bought at the
cost of losing the sequence.

This meets the unit--rank trichotomy of Remark~\ref{C-rem:dichotomy} from a second direction. There,
rank zero gave a classical gate and denied a continuous family, positive rank the reverse. Here,
rank one is exactly the cell in which the affine sieve has no traction. Stated as a single missing
property: what is wanted is a substitute for expansion at rank one, and that is the Mersenne problem
itself, not a lemma about \#307.
\end{remark}

\begin{remark}[The gap to the record, measured]\label{rem:sparserecord}
It is worth stating how far the required statement is from the best that is known, because the
distance is not one of degree.

Write a sequence's \emph{density exponent} $\theta$ for $\#\{n\le X\}\asymp X^{\theta+o(1)}$. The
sparsest sets currently proved to contain infinitely many primes are the values of binary forms:
$a^2+b^4$ at $\theta=3/4$, by Friedlander and Iwaniec \cite{FI1998}, and $a^3+2b^3$ at $\theta=2/3$,
by Heath-Brown \cite{HB2001}, the latter being the sparsest polynomial form known to capture its
primes. Our $q_n$ grow like $\varepsilon^{2n}$, so $\#\{n:q_n\le X\}\asymp\log X$ and
\[ \theta=0. \]

The distance from $2/3$ to $0$ is categorical rather than quantitative. Both records are bilinear
arguments: they run a sieve of dimension one and close it with a Type~II estimate obtained by
averaging over a set of size a positive power of $X$. Every gain in such an argument is a gain of a
power of $X$, and against $\log X$ terms a power of $X$ gains nothing; below $\theta=1/2$ the
standard sieve axioms already fail, and at $\theta=0$ there is no set left to average over. This is
the same wall as Remark~\ref{rem:affinesieve} approached from the analytic rather than the spectral
side, and it is why every sequence of exponent zero, Mersenne, Fibonacci, Lucas and ours alike, is
open together.

The structure our sequence does carry is the standard one and confers nothing. From
$A_Sq_n=x_n^2-D_S$ one gets $x_n^2\equiv D_S$ for every prime $p\mid q_n$ not dividing $D_S$, so all
prime factors of $q_n$ lie in the density--$\tfrac12$ set where $D_S$ is a quadratic residue. That
constrains the factorisations without excluding any, which is exactly the position the primitive
divisor theorems already left us in: divisors, never primality.
\end{remark}

\begin{remark}[The residue, atomised, and where its difficulty comes from]\label{rem:atomised}
Decomposing the question into independently known statements locates the obstruction on a single
atom and identifies its source.

Two atoms carry theorems. Every $q_n$ beyond the thirtieth has a primitive prime divisor
\cite{BHV2001}; and the greatest prime factor satisfies
$P^+(q_n)>n\exp\bigl(\log n/(104\log\log n)\bigr)$, by Stewart \cite{Stewart2013}, which settled
questions of Schinzel and Erd\H{o}s and was the first general advance on Bang and Carmichael. Every
other atom, that $q_n$ be squarefree infinitely often, that $\omega(q_n)$ be bounded infinitely
often, that $q_n$ be a $P_k$, and primality itself, is open.

The chain that would close the crux is $P^+$ large, hence $\omega$ small, hence prime. It breaks at
the first link, and by a measurable amount. Stewart's bound certifies
$\log P^+\ge(1+o(1))\log n$ against $\log q_n\approx2n\log\varepsilon$, so the fraction of the
term's size that is certified is of order $\log n/n$: $10.6\%$ at $n=10$, $0.32\%$ at $n=10^3$,
$4\times10^{-4}\%$ at $n=10^6$. Primality is the fraction $1$. No sharpening of the constant $104$
changes this, since $n^{1+o(1)}$ is the shape linear forms in logarithms produce. The Subspace
Theorem is the other available instrument and does not help: it is stronger under weaker hypotheses
but ineffective \cite{BCZ2003}, and an ineffective statement cannot exhibit the single prime term a
construction needs.

The source is visible in the comparison with the polynomial case. For $n^2+1$ the greatest prime
factor exceeds $n^{1.3}$ infinitely often \cite{Merikoski2019}, that is $q^{0.65}$, a \emph{positive
proportion} of the term's size, and on such sequences almost-primality is available too. For an
exponentially parametrised sequence only $(\log q)^{1+o(1)}$ is available, a vanishing proportion.
The gap between the two regimes is exactly the gap between polynomial and exponential
parametrisation, and Proposition~\ref{C-prop:nopoly} proves \#307 admits no polynomial one. So the
obstruction recorded in Remark~\ref{rem:affinesieve} as a question of rank, in
Remark~\ref{rem:sparserecord} as one of density exponent, and here as one of $P^+$, is a single
obstruction with three faces, and its source is Proposition~\ref{C-prop:nopoly}: forced exponentiality
is what places the residue in the class where nothing of the required strength is known.

Two things sharpen that attribution, and both make the obstruction less incidental rather than more.
First, the boundary is exactly where Proposition~\ref{C-prop:nopoly} sits. The fraction of a term's
size that Stewart's bound certifies is $\log P^+/\log q_n\approx\log n/\log q_n$, so it stays bounded
below precisely when $\log q_n\ll\log n$, that is when $q_n$ grows polynomially: degree $d$ gives the
constant $1/d$. Every super-polynomial rate sends it to $0$, and not only the exponential one:
$\exp(\sqrt n)$ and $n^{\log n}$ do as well. So polynomial growth is not one convenient case among
several, it is the whole of the favourable region, and it is exactly what
Proposition~\ref{C-prop:nopoly} excludes.

Second, the hypotheses of Proposition~\ref{C-prop:nopoly} are worth testing, since a gap there would
dissolve everything above. It assumes a \emph{fixed} number of polynomials, and its proof uses that
finiteness when it splits $\sum_i1/p_i$ into constants plus a remainder tending to $0$. A family
whose support grows with the parameter is therefore not covered. That gap is real and useless: a
growing support means a growing number of simultaneous primality conditions, which is a harder
demand than one, not an easier one, and it does not return the family to the polynomial regime. The
proposition forbids the case that would help and leaves uncovered only cases that would not.
\end{remark}

\begin{remark}[The squarefree atom, and what it reduces to]\label{rem:squarefree}
Among the four open atoms one is structurally easier, and it is worth saying why and how far it
gets. Squarefreeness of $q_n$ asks only for an \emph{upper} bound on a count of bad events, never
for the production of a prime, so it is the one atom not subject to the parity obstruction that
blocks bounded $\omega$, $P_k$, and primality.

The argument it wants is a union bound. For a prime $p\nmid2A_SD_S$, $p^2\mid q_n$ forces
$x_n^2\equiv D_S\pmod{p^2}$; the Pell orbit is purely periodic modulo $p^2$ with period
$\pi(p^2)$, there are at most two square roots of $D_S$, and each is met a bounded number of times
per period, so the density of bad $n$ attached to $p$ is at most $C/\pi(p^2)$. Where
$\pi(p^2)=p\,\pi(p)$, which is the generic behaviour, and $\pi(p)\asymp p$, this is $\asymp C/p^2$,
and the union bound closes as soon as $C\sum_pp^{-2}<1$. The constant is comfortable:
$\sum_pp^{-2}=0.4522474\ldots$, so any $C<2.21$ suffices, and a positive proportion of the $q_n$
would be squarefree, hence infinitely many.

What is missing is the hypothesis $\pi(p^2)=p\,\pi(p)$. Its failure is precisely a
Wall--Sun--Sun prime for the sequence, and there the density is $\asymp1/p$ rather than
$\asymp1/p^2$, whose sum diverges and destroys the bound. No such prime is known for any classical
sequence and their finiteness is unproved, so the atom does not close: it \emph{reduces} to a
Wall--Sun--Sun statement. That is a sharper position than ``open'', since it names the obstruction
and shows it is a single hypothesis rather than a programme, but it is a reduction and not a proof.

The unconditional pieces are formalised. \texttt{Erdos307.primeSquareSum\_certificate} gives
$\sum_pp^{-2}<\tfrac12$ in the finite form used above, the first six primes together with the bound
$\tfrac12\sum_{m\ge16}m^{-2}<\tfrac1{30}$ for the odd primes beyond; \texttt{good\_density\_pos} and
\texttt{squarefree\_positive\_density\_of\_bound} are the union bound and the conditional
implication. All carry $0$ \texttt{sorry} and the three standard axioms
\textup{(\texttt{lean/Erdos307/Squarefree.lean})}.

Two further points fix the standing of this atom, and both are negative. First, on the immune bases
the argument is blocked at its \emph{lowest} rung as well as its highest. The density attached to a
single small $p$ is $\#\{n\bmod\pi(p^2):p^2\mid q_n\}/\pi(p^2)$, read off the Pell orbit modulo
$p^2$, and that needs the fundamental solution, the object Proposition~\ref{C-prop:untestable} places
at $10^{24.7}$ operations. For a Lucas sequence with a known fundamental unit the small-prime rung is
a finite computation; here it is not. The ladder is therefore blocked at both ends, by
non-computability below and by Wall--Sun--Sun above. Second, and more to the point, squarefreeness is
not what the construction needs: Remark~\ref{rem:lehmer} requires $q$ to be \emph{prime}, and a
squarefree $q$ with several prime factors is useless there, so even the dream lemma of this atom
would leave \#307 exactly where it stands. Of the four atoms this is the easiest, and it is also the
only one whose resolution would not advance the problem.
\end{remark}

\begin{theorem}[Parity dichotomy]\label{thm:parity}
Let $a'=b,b'=a$ be a derivative two--cycle (so $a,b$ squarefree, $\gcd(a,b)=1$). Then exactly one of
the following holds.
\begin{enumerate}[label=\textup{(\roman*)},leftmargin=2.2em]
\item \emph{Mixed parity.} Exactly one of $a,b$ is even; say $a=2k$ with $k$ odd. Then $b=a'$ is odd
and $\omega(b)$ is even.
\item \emph{All odd.} Both $a,b$ are odd. Then $P\cup Q$ avoids $2$, so reaching reciprocal sum $>2$
needs $\ge 1412$ odd primes, and $\min(a,b)>10^{2519}$.
\end{enumerate}
\end{theorem}

\begin{proof}
They cannot both be even, as $\gcd(a,b)=\gcd(a,a')=1$. In case~(i), write $a=2k$, $k$ odd squarefree;
$b=a'=k+2k'$ is odd. For odd squarefree $b=\prod_{i}q_i$ one has
$b'=\sum_i b/q_i\equiv \omega(b)\pmod 2$ (each $b/q_i$ is odd). Since $b'=a$ is even,
$\omega(b)\equiv0\pmod2$. In case~(ii) the union of supports omits $2$; the reciprocal sum of the
first $1411$ odd primes is $<2$ while that of the first $1412$ exceeds $2$, so $|P\cup Q|\ge 1412$
and $\prod_{p\in P\cup Q}p\ge 10^{5040}$ (product of the $1412$ smallest odd primes, by direct
computation). Applying the same rigidity/ratio argument as in Theorem~\ref{thm:barrier}(b) (now
with $2\notin U$) gives $\min(a,b)\ge\sqrt{R/T(U)}>10^{2519}$. The all--odd case is therefore
vastly more constrained than the mixed case and, heuristically, empty.
\end{proof}

\begin{proposition}[Longer cycles are harder; two--cycles are extremal]\label{prop:kcycles}
Every member of a $k$--cycle of the arithmetic derivative is squarefree, and for $k\le3$ the members
are pairwise coprime (each pair is consecutive). Their masses $s(a_i)=a_{i+1}/a_i$ multiply to $1$,
so by AM--GM $\sum_i s(a_i)\ge k$; for $k\le3$ this is the union's reciprocal sum, whence the support
contains at least $m_k$ primes, where $m_k$ is least with $\sum_{i\le m_k}1/p_i\ge k$. Now $m_2=59$
and $m_3=361{,}139$ (last prime $5{,}195{,}977$), so a three--cycle has at least $361{,}139$ primes in
its support and product exceeding $10^{2.26\times10^{6}}$ (hence its largest member exceeds
$10^{752{,}194}$); by Mertens $m_k$ grows doubly exponentially. Two--cycles are thus the extremal,
and only practically relevant; case for that measure. For $k\ge4$ non--consecutive members may share
primes, so the displayed bound is one on the \emph{multiset} reciprocal sum; the union's own sum is
bounded in Proposition~\textup{\ref{prop:kuniform}}.
\end{proposition}

\begin{proposition}[The barrier is uniform in the period]\label{prop:kuniform}
Let $a_1\to a_2\to\dots\to a_k\to a_1$ be a cycle of the arithmetic derivative with $k\ge2$ and
$a_i\neq a_{i+1}$, and let $U=\bigcup_i\operatorname{supp}(a_i)$. Then
\[ \sum_{p\in U}\frac1p\ \ge\ \frac{k}{\lfloor k/2\rfloor}, \]
so $|U|\ge m(k)$, the least $n$ with $T_n\ge k/\lfloor k/2\rfloor$, and $\prod_{p\in U}p\ge\Pi_{m(k)}$.
Since $k/\lfloor k/2\rfloor\ge2$ always, with equality exactly for even $k$,
\[ |U|\ \ge\ 59\quad\text{and}\quad \prod_{p\in U}p\ >\ 8.77\times10^{112} \]
for a cycle of \emph{every} period.
\end{proposition}
\begin{proof}
Rigidity gives $\gcd(a_i,a_i')=1$, and $a_{i+1}=a_i'$, so $a_i$ and $a_{i+1}$ are coprime. Hence for
each prime $p$ the set $S_p=\{i:p\mid a_i\}$ contains no two cyclically consecutive indices, i.e.\ it
is an independent set in the cycle graph $C_k$, and therefore $|S_p|\le\lfloor k/2\rfloor$.
Consequently
\[ \sum_{i}\sigma(a_i)\;=\;\sum_{p\in U}\frac{|S_p|}{p}\;\le\;\Bigl\lfloor\tfrac k2\Bigr\rfloor
   \sum_{p\in U}\frac1p . \]
The masses satisfy $\prod_i\sigma(a_i)=\prod_i a_{i+1}/a_i=1$, so $\sum_i\sigma(a_i)\ge k$ by AM--GM.
Combining the two displays gives the claim, and $T$ is increasing.
\end{proof}

The bound is uniform rather than growing, and it is the union support that a search or a sieve sees.
Odd periods are genuinely harder, even ones are not:
\[
\begin{array}{l|cccccccc}
k & 2 & 3 & 4 & 5 & 6 & 7 & 8 & 9\\\hline
k/\lfloor k/2\rfloor & 2 & 3 & 2 & 5/2 & 2 & 7/3 & 2 & 9/4\\
m(k) & 59 & 361{,}139 & 59 & 1{,}413 & 59 & 400 & 59 & 231
\end{array}
\]
so $m(9)<m(7)<m(5)$: the sequence is not monotone in $k$, and $m(k)=59$ \emph{exactly} for every
even $k$, not merely in the limit. The consequence
worth naming is that Theorem~\ref{thm:barrier} is not a statement about period two at all: every
nontrivial cycle of the arithmetic derivative, of any length, carries at least $59$ primes and a
product past $8.77\times10^{112}$. The barrier therefore bounds the whole of
Ufnarovski--\AA hlander's Conjecture~3 and not merely its period--two case, which is \#307.

\begin{corollary}[What a bounded sweep rules out, for every period at once]\label{cor:sweepperiod}
Let $a_1\to\dots\to a_k\to a_1$ be a cycle with every member at most $X$. Then
\[ k\ \ge\ \frac{\log_{10}\Pi_{59}}{\log_{10}X}\ =\ \frac{112.94\ldots}{\log_{10}X}. \]
\end{corollary}
\begin{proof}
Proposition~\ref{prop:kuniform} gives $|U|\ge59$ for a cycle of any period, so
$\prod_{p\in U}p\ge\Pi_{59}>8.77\times10^{112}$. Every prime of $U$ divides some member, so
$\prod_{p\in U}p$ divides $\prod_i a_i\le X^{k}$.
\end{proof}

\[
\begin{array}{l|cccccc}
X & 10^{6} & 10^{7} & 10^{9} & 10^{12} & 10^{18} & 10^{30}\\\hline
\text{rules out every }k\le & 18 & 16 & 12 & 9 & 6 & 3
\end{array}
\]
The sweep of Proposition~\ref{C-prop:nearmiss}, now carried to $10^{8}$, therefore already excludes
every cycle of period at most $14$ with all members in range, and not merely period two. This is the practical form
of Proposition~\ref{prop:kuniform}: because the barrier does not weaken with the period, a single
bounded search settles a whole initial segment of Ufnarovski--\AA hlander's Conjecture~3 at once,
whereas the period--two reading of the same computation settles only $k=2$.

\begin{proposition}[The barrier as a function of period and smallest prime]\label{prop:twoparam}
Let $a_1\to\dots\to a_k\to a_1$ be a cycle with $k\ge2$, let $U$ be the union of the supports, and
suppose $\min U\ge P$. With $c(k)=k/\lfloor k/2\rfloor$, let $M(k,P)$ be the least $m$ for which the
$m$ smallest primes $\ge P$ have reciprocal sum at least $c(k)$. Then $|U|\ge M(k,P)$, and
$\prod_{p\in U}p$ is at least the product of those $m$ primes.
\end{proposition}
\begin{proof}
Proposition~\ref{prop:kuniform} gives $\sum_{p\in U}1/p\ge c(k)$, and among sets of $m$ primes all
$\ge P$ the reciprocal sum is maximised by the $m$ smallest such primes.
\end{proof}

\[
\begin{array}{l|ccccc}
M(k,P) & k=2 & k=3 & k=5 & k=7 & k=9\\\hline
P=2 & 59 & 361{,}139 & 1{,}413 & 400 & 231\\
P=3 & 1{,}412 & \text{---} & 361{,}138 & 40{,}260 & 15{,}473\\
P=5 & 40{,}259 & \text{---} & \text{---} & \text{---} & 1{,}266{,}094\\
P=7 & 588{,}500 & \text{---} & \text{---} & \text{---} & \text{---}
\end{array}
\]
(dashes exceed $\pi(4\times10^{7})$; \texttt{code/kcycle\_uniform.gp}). The family contains three
results proved elsewhere in this note as special cases, and reproduces each of them exactly:
$M(2,2)=59$ is Theorem~\ref{thm:barrier}; $M(2,3)=1412$, with product $10^{5040.1}$ and hence
$\min(\prod P,\prod Q)>10^{2519}$, is the all--odd branch of Theorem~\ref{thm:parity}; and
$M(3,2)=361{,}139$ is Proposition~\ref{prop:kcycles}. That three independent derivations land on the
same integers is the check the family deserves.

Two readings. The parity dichotomy is not special to period two: an all--odd cycle of \emph{any} even
period needs $1{,}412$ primes and a product past $10^{5040}$, and of period $9$ still needs
$15{,}473$. And the family supplies, in valid form, the conditional statements that
Proposition~18 of \cite{Kovic2012} asserts but does not establish: for $\min U\ge3$ that proposition
claims $r+s\ge57$ where the true bound is $M(2,3)=1{,}412$, and for $\min U\ge5$ it claims
$r+s\ge110$ where the true bound is $M(2,5)=40{,}259$ --- larger by factors of $24.8$ and $366$.

\section{A Th\=abit--type rule and its honest limits}\label{sec:frame}

The bridge recasts the search for a solution as the search for a derivative two--cycle, and the
derivative's Leibniz structure makes the cycle equations \emph{linear} in the ``closing'' primes.
This yields a constructive calculus analogous to Th\=abit ibn Qurra's rule for amicable pairs and to
Euler's divisor method.

\begin{theorem}[Frame Rule]\label{thm:frame}
Let $M,N$ be coprime squarefree integers with derivatives $M',N'$, and set $\Delta\coloneqq MN-M'N'$.
Suppose
\[
p=\frac{MN'+N^2}{\Delta},\qquad q=\frac{M'N+M^2}{\Delta}
\]
are integers, both prime, each coprime to $MN$, with $p\ne q$. Then $a\coloneqq Mp$, $b\coloneqq Nq$
form a derivative two--cycle $a'=b$, $b'=a$; equivalently $(P,Q)$ with $P$ the prime set of $Mp$ and
$Q$ that of $Nq$ solves Erd\H{o}s \#307.

\textup{Formalised.} \texttt{Erdos307.frame\_cycle} verifies both cycle equations from the two
divisibility hypotheses, and \texttt{frame\_solve} runs the derivation backwards, so the two
quotients are \emph{forced} rather than guessed and the rule is an equivalence. The Leibniz step
$a'=M'p+M$ is \texttt{csum\_insert\_prime}. What no algebra reaches is the primality of the two
quotients \textup{(\texttt{lean/Erdos307/Frame.lean})}.
\end{theorem}

\begin{proof}
With $p$ prime and $p\nmid M$, $a=Mp$ is squarefree and $a'=M'p+M$; similarly $b'=N'q+N$. The cycle
conditions $a'=b$, $b'=a$ read
\[
M'p+M=Nq,\qquad N'q+N=Mp,
\]
a linear system in $(p,q)$ with determinant $-(MN-M'N')=-\Delta$. Solving gives exactly
$p=(MN'+N^2)/\Delta$ and $q=(M'N+M^2)/\Delta$; substituting back verifies both equations
identically. (Both identities were checked symbolically.)
\end{proof}

A two--primes--per--side variant (``Rule~2'') is obtained identically: with $a=Mp_1p_2$, $b=Nq_1q_2$
and elementary symmetric functions $e_1,e_2,f_1,f_2$, the cycle equations are linear in the $f_i$,
\[
f_2=\frac{Me_1+M'e_2}{N},\qquad f_1=\frac{MNe_2-MN'e_1-M'N'e_2}{N^2},
\]
and a solution is realised whenever $f_1^2-4f_2$ is a perfect square with prime roots.

\paragraph{An Euler-style divisor trick.}
Fix $N$ and let $M=M_0r$ with $r$ a new prime, so $M'=M_0'r+M_0$. Put
$\Delta_0\coloneqq M_0N-M_0'N'$.

\begin{proposition}[Divisor trick]\label{prop:divisor}
With the above notation and $p_{\mathrm{num}}\coloneqq MN'+N^2$:
\begin{enumerate}[label=\textup{(\alph*)},leftmargin=2.2em]
\item $\Delta(r)\coloneqq MN-M'N'=r\,\Delta_0-M_0N'$ \emph{(linear in $r$)};
\item $\Delta_0\,p_{\mathrm{num}}-M_0N'\,\Delta(r)=K$ where $K\coloneqq \Delta_0N^2+(M_0N')^2$ is
\emph{constant in $r$};
\item consequently $\Delta(r)\mid p_{\mathrm{num}}\Rightarrow \Delta(r)\mid K$, and the converse
holds provided $\gcd(\Delta(r),\Delta_0)=1$ \textup{(}equivalently $\gcd(M_0N',\Delta_0)=1$, since
$\Delta(r)\equiv-M_0N'\bmod\Delta_0$\textup{)};
\item each positive divisor $d\mid K$ with $d\equiv-M_0N'\ (\mathrm{mod}\ \Delta_0)$ yields a
candidate $r=(d+M_0N')/\Delta_0$ and recovers $\Delta(r)=d$; this enumeration is complete (it
captures every genuine solution). The resulting $p=p_{\mathrm{num}}/d$ is integral when
$\gcd(M_0N',\Delta_0)=1$, and otherwise must be checked.
\end{enumerate}
\end{proposition}

\begin{proof}
(a) and (b) are direct expansions (verified symbolically; $\partial K/\partial r=0$). (c): from
(b), $\Delta(r)\mid p_{\mathrm{num}}\Rightarrow\Delta(r)\mid\Delta_0p_{\mathrm{num}}
=M_0N'\Delta(r)+K\Rightarrow\Delta(r)\mid K$. For the converse, $\Delta(r)\mid K$ gives
$\Delta(r)\mid\Delta_0p_{\mathrm{num}}$, whence $\Delta(r)\mid p_{\mathrm{num}}$ iff
$\gcd(\Delta(r),\Delta_0)=1$. (d) inverts (a). Excluded degeneracies: $\Delta_0=0$ ($r$ undefined)
and $K\le0$.
\end{proof}

Thus, exactly as in Euler's amicable--pair machinery, one factors the single number $K$ once and
reads off candidates from its divisors lying in a fixed residue class modulo $\Delta_0$; only the
primality of $r,p,q$ then remains. Two caveats deserve emphasis:
the converse in (c) and the automatic integrality of $p$ require $\gcd(M_0N',\Delta_0)=1$ (it fails,
e.g., for $M_0=2$, $N=19\cdot47$, $r=5$, where $\Delta(r)\mid K$ but $\Delta(r)\nmid
p_{\mathrm{num}}$), and the count of admissible classes is governed by $\Delta_0$, not merely by
$\tau(K)$.

When the closing side carries a \emph{single} new prime, the rule collapses to a normal form that
separates the two layers of the problem.

\begin{proposition}[One--prime normal form; the exactness stratum]\label{prop:oneprime}
\textup{(a)} Erd\H{o}s \#307 has a solution if and only if there exist squarefree $M,N$ with
\[ MN-M'N'=M^2, \]
such that $p\coloneqq N'/M$ is a prime not dividing $M$. Indeed the equation forces $M\mid M'N'$, hence
$M\mid N'$ by $\gcd(M,M')=1$, and $N=M+M'p$; then $a=Mp$, $b=N$ satisfy $a'=M'p+M=N$ and $b'=N'=Mp=a$
(coprimality and $a\ne b$ are automatic). Conversely every solution $(a,b)$ and every prime $p\mid a$
yield such a pair with $M=a/p$, so each solution carries $\omega(a)$ witnesses. The problem thus splits
into an \emph{exactness} layer (the Diophantine equation) and a single \emph{primality} layer (one prime
value of the determined quantity $N'/M$).
\textup{(b)} \emph{The exactness layer alone carries the full barrier.} Call $(M,N)$ a \emph{relaxed
witness} if $MN-M'N'=M^2$ with $p=N'/M$ composite. The identity
$\bigl(\sigma(M)+\tfrac1p\bigr)\sigma(N)=\frac{N}{Mp}\cdot\frac{Mp}{N}=1$ holds for \emph{every} witness,
so by strict AM--GM ($\sigma(N)=1$ being impossible) $\sigma(M)+\sigma(N)>2-\tfrac1p$. Rigidity forces
three structural facts. The supports of $M$ and $N$ are \emph{automatically} disjoint: a common prime
$q$ would divide $N'=Mp$ through $q\mid M$, against $\gcd(N,N')=1$. Likewise $\gcd(p,MN)=1$: a prime
shared by $p$ and $M$ propagates to $N=M+M'p$ and lands in the previous case, while one shared by $p$
and $N$ divides $N'$. And if $p$ is even, both members are odd: $2\mid N'=Mp$ forces $N$ odd by the same
lemma, while $M$ even would make $N=M+M'p$ even. The union of supports therefore avoids every prime
divisor of $p$, and for even $p$ avoids $2$; exact greedy computation then gives, for \emph{every}
composite $p$,
\[ \omega(MN)\ \ge\ 59, \qquad MN\ >\ 8.77\times10^{112} \]
the cycle barrier itself; with strictly stronger floors at every finite $p$: $\omega\ge194$,
$MN>10^{497}$ at $p=9$; $\omega\ge230$, $MN>10^{609}$ at $p=4$; the even--$p$ floors rising towards the
all--odd bound $10^{5040}$, and the infimum $8.77\times10^{112}$ approached only as $p\to\infty$ through
composites free of prime factors $\le277$. Deleting primality thus lowers the barrier by
\emph{nothing}: the $59$--prime, $10^{112}$ wall is carried entirely by the exactness equation, and the
primality of $N'/M$ is a separate demand on top of it. (No witness, prime or composite, exists with
$M\le120$, $p\le2000$, exhaustively.)
\end{proposition}

\paragraph{The decisive obstruction: a circular objective.}
The heuristic yield of the divisor trick from one frame is
\[
\mathbb{E}[\text{cycles}] \ \approx\ A\cdot\frac{1}{\ln r\,\ln p\,\ln q},
\qquad A=\#\{d\mid K:\ d\equiv -M_0N'\ (\mathrm{mod}\ \Delta_0)\}\ \approx\ \frac{\tau(K)}{\Delta_0},
\]
the last step being an equidistribution heuristic. To make $\mathbb{E}\ge1$ at the barrier scale
($1/\ln^3\approx 5\times10^{-7}$) one needs $A\gtrsim 2\times10^6$ admissible divisors per frame,
hence $\Delta_0=O(1)$ with $\tau(K)\sim10^6$. The natural objective is therefore ``engineer frames
with $\Delta_0$ small.'' But here is the obstruction, an exact identity:
\[
\boxed{\ \Delta_0 \;=\; M_0N\bigl(1-s_{M_0}\,s_N\bigr),\qquad
s_{M_0}=\!\!\sum_{p\mid M_0}\!\frac1p,\ \ s_N=\sum_{p\mid N}\frac1p.\ }
\]
So $\Delta_0$ is small \emph{relative to} $M_0N$ exactly when $s_{M_0}s_N\approx1$ (which is the
Erd\H{o}s~\#307 condition for the frames themselves) and $\Delta_0=O(1)$ in absolute terms at
scale requires $s_{M_0}s_N=1-O(1/M_0N)$, i.e.\ a reciprocal--product within $O(1/M_0N)$ of $1$,
\emph{strictly harder} than \#307. The frame--engineering objective is thus circular: the Frame Rule
is a faithful \emph{parametrisation} of solutions, not a reduction in difficulty.

\begin{remark}[Honest status of the constructive program]
The Frame Rule (Theorem~\ref{thm:frame}) and divisor trick (Proposition~\ref{prop:divisor}) are
exact theorems; they make the existence of a solution \emph{falsifiable in principle by finite
computation per frame}, and they place \#307 in the same constructive lineage as Th\=abit's and
Euler's rules. They do not, as they stand, lower the difficulty below \#307 itself: by the boxed
identity the bottleneck (small $\Delta_0$) is the original problem in disguise, and a complete
small--frame sweep ($2.2\times10^5$ frames) produces zero cycles with total expected count $<1$, as
the barrier demands. We regard the calculus as an exact parametrisation of solutions whose central
obstruction we have now identified precisely; not a pathway that lowers the difficulty, and not
evidence that a solution is within reach.
\end{remark}

\section{Computations}\label{sec:computations}

All verdicts use exact integer/rational arithmetic; floating point is used only for pre--screening.
\begin{itemize}[leftmargin=1.4em]
\item \emph{Exhaustive non--existence below $10^8$.} Enumerating all prime sets $P$ with
$\sum1/p>1$ and $\prod P\le10^8$ (forcing $Q$ as the prime support of $N_P$ and testing
$N_Q=D_P$): $1{,}075{,}419$ candidate sets, no solution. Subsumed by Theorem~\ref{thm:barrier} but an
independent sanity check.
\item \emph{Derivative sieve.} No solution of $n''=n$ with $n'\ne n$ and both members $\le
2\times10^7$ (smallest--prime--factor sieve of the derivative).
\item \emph{Function field.} The arithmetic derivative on $\mathbb{F}_q[t]$ has no two--cycles and no
nonzero fixed points, verified exhaustively over $\mathbb{F}_2$ (degree $\le7$), $\mathbb{F}_3$
($\le5$), $\mathbb{F}_5$ ($\le4$), confirming Theorem~\ref{C-thm:ff}.
\item \emph{Density.} The sieve gives density $0.042029$ to $2\times10^7$ ($0.017623$ among
squarefree); the independent Erd\H{o}s--Wintner convolution over primes $\le10^5$ gives
$\delta=0.04202$, agreeing to four places (Proposition~\ref{C-prop:density}).
\item \emph{Conditional mass.} Of $43{,}087$ qualifying squarefree $a$ sampled, $b=a'$ is squarefree
$95.0\%$ of the time; on the resulting $40{,}920$ squarefree $b$ the mean reciprocal mass is $0.047$
vs.\ the required $1/s\approx0.934$, with empirical tail probability $\Pr[\sum_{q\mid b}1/q\ge1/s]=
0/40{,}920$.
\item \emph{Frame/ledger sweep.} Over $2.2\times10^5$ factorable small frames: the divisor--trick
identities hold on every instance; the smallest $\Delta_0$ observed is $22$; the equidistribution
$A\approx\tau(K)/\Delta_0$ holds only in aggregate ($\sum A=127$ vs.\ $\sum\tau(K)/\Delta_0=39.8$)
and fails per frame ($A=0$ for $99.95\%$ of frames); total expected cycles $<1$; zero cycles, as the
barrier requires.
\item \emph{$\mu$--Sondow line census.} Enumerating the lines $n'=n\pm d$ (squarefree members coprime to
$d$) finds no adjacent opposite--sign pair at additive distance $d$: exhaustively to $2\times10^7$ for
$|d|\le200$, and along the small--$|\delta|$ genealogy tree out to members of more than a thousand digits.
Every $\delta=\pm1$ member encountered is a known Giuga or primary pseudoperfect number; the closest
opposite--sign approach recorded is $|y-(x+d)|=11$ (at $d=1$, $30$ vs.\ $42$, and $d=37$, $210$ vs.\ $258$),
attained only among small members; the two trees do not come close at height.
\end{itemize}
Code and logs accompany this note.

\section{Formalisation in Lean~4}\label{sec:lean}

The rigidity and barrier results are formalised against mathlib (Lean~4, toolchain
\texttt{v4.30.0}). The following are machine--checked and \texttt{sorry}--free, depending only on the
standard axioms \texttt{propext}, \texttt{Classical.choice}, \texttt{Quot.sound}:
\begin{itemize}[leftmargin=1.4em]
\item Lemma~\ref{thm:rigidity} (\texttt{rigidity\_coprime}) and Theorem~\ref{thm:structure}
(\texttt{solution\_structure});
\item Lemma~\ref{lem:amgm} (\texttt{reciprocal\_sum\_gt\_two}); the exact thresholds
$\sum_{\text{first }58}1/p<2<\sum_{\text{first }59}1/p$ grounding $|U|\ge59$
(\texttt{recipSum58\_lt\_two}, \texttt{recipSum59\_gt\_two}); the numeric core
$\Pi_{59}^2\ge(4\times10^{112})N_{59}$ (\texttt{barrier\_numeric}), and the algebraic core of the
barrier (\texttt{barrier}), giving $D_P^2\ge4\times10^{112}$;
\item Theorem~\ref{C-thm:halflyap} (\texttt{log\_lt\_half\_sub}, \texttt{delta\_decreasing},
\texttt{two\_le\_add\_of\_mul\_eq\_one}) and Theorem~\ref{C-thm:lyapfalse}
(\texttt{criterion\_fails\_of\_pos}, \texttt{criterion\_fails\_of\_nonpos},
\texttt{no\_lambda\_works}, \texttt{lyapunov\_criterion\_false}), neither of which uses
\texttt{native\_decide};
\item Proposition~\ref{C-prop:masssumthreshold} (\texttt{half\_works}, \texttt{no\_f\_above\_two},
\texttt{mass\_sum\_threshold}) and Proposition~\ref{C-prop:lyapuniversal}
(\texttt{exists\_f\_of\_injective}, \texttt{increment\_iff},
\texttt{no\_two\_cycle\_of\_decreasing}), which together machine-check the closure of the sign
branch in Theorem~\ref{C-thm:noinvariant};
\item Proposition~\ref{C-prop:infzero} (\texttt{exists\_block\_sum\_near}, \texttt{infimum\_zero});
\item Theorem~\ref{C-thm:noinvariant}(1) in its mechanism (\texttt{Erdos307.Congruential}) and
Proposition~\ref{C-prop:liveslot} (\texttt{live\_slot\_threshold}, \texttt{threshold\_bracket}).
\end{itemize}
The single non--elementary ingredient, that the $k$ smallest primes minimise $\prod$ and maximise
$\sum1/p$ among $k$--element prime sets, is isolated as an explicit hypothesis of \texttt{barrier}
(the ratio bound $R/T(U)\ge\Pi_{59}/T_{59}$), so the dependency is transparent rather than hidden.
It is no longer an open dependency: \texttt{Erdos307.Extremal} proves it
(\texttt{prod\_first\_primes\_le}, \texttt{recipSum\_le\_first\_primes},
\texttt{hRatio\_of\_extremal}), which is why Theorem~\ref{thm:barrier} appears above in the
hypothesis-free form \texttt{erdos307\_barrier\_closed}.

On the trust boundary: there is none beyond the logic. Every declaration in the development,
without exception, depends only on \texttt{propext}, \texttt{Classical.choice} and
\texttt{Quot.sound}; there is no \texttt{native\_decide} site and no \texttt{ofReduceBool}.

The one that resisted was the pruned search behind Proposition~\ref{prop:close59}. \texttt{dfs}
compares rationals, and the Lean kernel cannot reduce a \texttt{Rat} comparison at all, since
\texttt{Nat.gcd} is well-founded recursion rather than a \textup{GMP} primitive; \texttt{decide}
does not merely run slowly on it, it fails to evaluate. Three changes move the execution into the
kernel. Scaling by $M=\prod(\texttt{forced39}\cup\texttt{pool99})$, a $327$--digit integer, clears
every denominator \emph{exactly}: each $1/p$ becomes $M/p$ and $\mathrm{thr}$ becomes the integer
$M\cdot\mathrm{thr}$, so the scaled search branches identically and the soundness theorem
\texttt{dfs\_sound} is reused unchanged. The leaf test replaces \texttt{Nat.sqrt} on a $130$--digit
number, some hundreds of structural steps, by a residue certificate: eleven moduli
$\{8,19,167,127,149,67,59,101,73,191,41\}$, with their quadratic residues stored as bitmasks, cover
all $49{,}961$ supports, and a non--residue is a proof of non--squareness
\textup{(\texttt{lean/Erdos307/NonSquare.lean})}. Finally the product and cofactor sum are carried
down the recursion rather than rebuilt at each leaf. The kernel then runs the whole search in about
a minute \textup{(\texttt{dfsA\_run}, \texttt{dfs\_of\_dfsA})}.

Proposition~\ref{prop:close59}'s product bounds are proved in this paper and not in Lean:
\texttt{erdos307\_sixty} establishes $|P\cup Q|\ge60$ alone.

\section{Why every known coprime example uses $1$}\label{sec:coprime}

Corollary~\ref{cor:coprime60} shows the weakened variant is no smaller than \#307 itself: with
$1\notin P\cup Q$ a solution needs $60$ pairwise coprime elements, and at $59$ every element is
forced prime. What that leaves unexplained is the shape of the examples that \emph{are} known ---
$1=(1+\frac15)(\frac12+\frac13)$, and the family generated by primary pseudoperfect numbers such as
$1806=2\cdot3\cdot7\cdot43$ giving $1=(1+\frac1{1805})(\frac12+\frac13+\frac17+\frac1{43})$. Every
one of them uses the element $1$. That is not an accident of search.

\begin{proposition}[No coprime family has reciprocal sum $1$]\label{prop:coprimeone}
Let $A$ be a finite set of pairwise coprime integers $\ge2$. Then $\sum_{a\in A}1/a\ne1$.
$0$ \texttt{sorry} \textup{(\texttt{Erdos307.Coprime.sum\_inv\_ne\_one})}.
\end{proposition}
\begin{proof}
Suppose $\sum_{a\in A}1/a=1$ and fix $a\in A$; put $R=\prod_{b\in A,\,b\ne a}b$. Multiplying by $R$
gives $R/a+\sum_{b\ne a}R/b=R$. Each term of the sum is an integer, since $b\mid R$ for $b\ne a$, and
so is $R$; hence $R/a\in\mathbb{Z}$, i.e.\ $a\mid R$. But $a$ is coprime to every $b\ne a$, so
$\gcd(a,R)=1$ and $a=1$, contradicting $a\ge2$.
\end{proof}

\begin{remark}\label{rem:coprimeedge}
Every known example of the weakened variant has one factor equal to $1+1/n$ and the other equal to
$n/(n+1)$, i.e.\ one side is a coprime family whose reciprocal sum is $1-1/(n+1)$ and which is
completed to $1$ only by adjoining the element $1$. Proposition~\ref{prop:coprimeone} says the
adjunction is unavoidable: the element $1$ contributes a full unit to the sum while dividing nothing,
and it is the only element that can, because for every other element the divisibility above closes.
So the known examples are not near misses for the open form --- they are on the other side of a
divisibility, and no refinement of them approaches $1\notin P\cup Q$.

The rigidity step that Corollary~\ref{cor:coprime60} inherits from the prime case is the same
divisibility in cofactor form: for a coprime family $A$ with product $\alpha$, the numerator
$M=\sum_{a\in A}\alpha/a$ satisfies $\gcd(M,a)=1$ for every $a\in A$, because every term but
$\alpha/a$ is divisible by $a$ while $\alpha/a=\prod_{b\ne a}b$ is coprime to it. This is
\texttt{Erdos307.Coprime.coprime\_numerator}, and it closes the one step of
Corollary~\ref{cor:coprime60}'s proof that was previously carried in prose.
\end{remark}

\subsection{A rejection sieve that factors nothing}\label{ssec:coprimesieve}

The obvious search for a $1$--free coprime example fixes one side and reads off the other: a coprime
family $A$ with product $\alpha$ determines its partner $\beta=\sum_{a\in A}\alpha/a$, and $B$ must
then be a coprime factorisation of $\beta$ with $\sum_{b\in B}1/b=\alpha/\beta$. The recorded
obstacle is that testing a candidate appears to require factorising $\beta$, which at the sizes
Corollary~\ref{cor:coprime60} forces is out of reach. It does not.

\begin{proposition}[A necessary condition, computable in small modular arithmetic]\label{prop:coprimesieve}
Let $A$ be a coprime family with $\alpha=\prod A$, $x=\sum_{a\in A}1/a$ and
$\beta=\sum_{a\in A}\alpha/a$. If some coprime family $B$ satisfies $\bigl(\sum_A1/a\bigr)
\bigl(\sum_B1/b\bigr)=1$, then
\[ y_{\max}(\beta)\;:=\;\sum_{p^e\,\|\,\beta}\frac1{p^e}\;\ge\;\frac1x .\]
Moreover $y_{\max}(\beta)$ is computable without forming or factorising $\beta$: for every prime
power $p^e$ with $p\nmid\alpha$, Lemma~\ref{C-lem:symbolfact} gives
$\beta\equiv\alpha\sum_{a\in A}a^{-1}\pmod{p^e}$, so each $v_p(\beta)$ is decided by arithmetic
modulo $p^e$ alone.
\end{proposition}
\begin{proof}
Write $y=\sum_B1/b=N/\gamma$ with $\gamma=\prod B$ and $N=\sum_{b\in B}\gamma/b$. As in
Remark~\ref{rem:coprimeedge}, $\gcd(\beta,\alpha)=\gcd(N,\gamma)=1$, so $xy=1$ reads
$\beta N=\alpha\gamma$ with $\beta\mid\gamma$ and $N\mid\alpha$, forcing $\gamma=\beta$ and
$N=\alpha$. Hence $B$ is a coprime factorisation of $\beta$ and $\sum_B1/b=\alpha/\beta=1/x$. Every coprime factorisation is obtained
from the factorisation into prime powers by merging blocks, and merging strictly decreases the
reciprocal sum, since $1/(mn)<1/m+1/n$ for $m,n\ge2$ \textup{(}$0$ \texttt{sorry},
\texttt{Erdos307.Coprime.merge\_loses}\textup{)}. Hence the prime--power factorisation is the
maximiser and $1/x=\sum_B1/b\le y_{\max}(\beta)$. The congruence is
Lemma~\ref{C-lem:symbolfact} with $r=p^e$, legitimate because $p\nmid\alpha$.
\end{proof}

The test is therefore free: primes above a trial bound $Z$ contribute below $1/Z$ apiece, so trial
moduli up to $10^5$ already pin $y_{\max}$ to within $10^{-5}$, using no integer larger than $10^{10}$.
Verified against direct multiprecision evaluation of $v_p(\beta)$ over $200$ random families and all
$p\le500$: $13{,}065$ valuations, $0$ mismatches \textup{(\texttt{code/beta\_smoothness.rs},
checked by the script in the same file's header)}.

What the sieve says is stark. Over $5{,}272$ splits of the first $60$ primes lying in the mass window
$|xy-1|\le10^{-2}$, the mean of $y_{\max}(\beta)$ is $0.1795$ and its maximum is $0.8371$, against a
requirement of $1/x\approx0.95$: \emph{not one candidate survives}
\textup{(\texttt{code/beta\_smoothness.rs})}. The mechanism is a tension between the two constraints.
Reaching $x\approx1$ obliges $A$ to hold small primes; but $\gcd(\alpha,\beta)=1$ bars every prime of
$A$ from dividing $\beta$, and $y_{\max}$ is carried almost entirely by the small primes. The mass
window and the smoothness requirement pull on the same scarce resource.

The same tension defeats the complementary route, which uses Lemma~\ref{C-lem:symbolfact} in the
forward direction: $b\mid\beta$ is equivalent to $\sum_{a\in A}a^{-1}\equiv0\pmod b$, a subset--sum
condition on $A$ that again needs no factorisation, so one may try to \emph{impose} the divisibilities
$\beta$ must satisfy. Annealing on the forced divisor $D(A)=\prod\{b\in B:\ b\mid\beta\}$ over ground
sets of $60$ to $70$ primes reaches $\log_{10}D\approx10$ against $\log_{10}\beta\approx50$--$75$ in a
window $|xy-1|\le10^{-2}$, and degrades to $\log_{10}D\approx3$ as the window tightens to $10^{-7}$,
below which no split is found at all \textup{(\texttt{code/split\_forcing.rs})}. Since a random split
already satisfies $\sum_{b\in B}1/b\approx1$ of the congruences by chance, the search inside the mass
window is performing at the random baseline: the freedom the congruence route needs is exactly the
freedom the mass window removes.

\subsection{Candidates that are smooth by construction}\label{ssec:smoothbuilt}

Both failures above share an order of operations: fix a ground set $U$, split it, and hope $\beta$
comes out smooth. Reversing the order removes the obstruction. Choose the small primes $B_0$ that
are to divide $\beta$, draw $A$ from primes above $\max B_0$, and \emph{impose}
$\sum_{a\in A}a^{-1}\equiv0\pmod q$ for each $q\in B_0$. Two points make this work where the earlier
routes did not. First, the conditions are congruences on $A$ alone, so they can be solved exactly
rather than sampled --- we use meet--in--the--middle over disjoint swap pairs drawn from adjacent
primes, so the mass moves by less than $10^{-4}$. Second, and decisively, $B$ is no longer the
complement of a fixed ground set but whatever factors $\beta$; the rigid window $x+y=T(U)$ of
Section~\ref{ssec:coprimesieve} therefore does not apply, and the only surviving requirement is
$x>1$.

Two obstructions appear and both are removable. Every pool prime is odd, so $a^{-1}\equiv1\pmod2$ and
the residue mod $2$ is pinned to $|A|\bmod2$: no swap can move it, and $|A|$ must be even. And
forcing $q\mid\beta$ does not prevent $q^2\mid\beta$, in which case the block contributes $1/q^2$ in
place of $1/q$ --- at $B_0=\{2,\dots,41\}$ an unconstrained solution came out with $2^5\,\|\,\beta$ and
$3^3\,\|\,\beta$, which alone costs $0.76$ and fails the sieve. Imposing $q\,\|\,\beta$, tested by the
same identity taken modulo $q^2$, repairs it.

\begin{proposition}[The necessary condition is constructible]\label{prop:smoothbuilt}
For every $B_0$ among the first $8$ to $12$ primes there are coprime families $A$ of primes above
$\max B_0$ with $x=\sum_A1/a>1$ such that each $q\in B_0$ divides $\beta$ exactly once. Such $A$
satisfy $y_{\max}(\beta)\ge\sum_{q\in B_0}1/q>1>1/x$, hence pass
Proposition~\ref{prop:coprimesieve}.
\end{proposition}

Verified in exact multiprecision arithmetic at $B_0=\{2,\dots,37\}$: $|A|=3290$, largest member
$31{,}063$, $\log_{10}\alpha=\log_{10}\beta=13{,}163.1$, $\gcd(\alpha,\beta)=1$, every forced prime
exactly dividing $\beta$, and
\[ y_{\max}(\beta)=1.5927\ldots\ \ge\ 1/x=0.99597\ldots \]
\textup{(\texttt{code/smooth\_by\_construction.rs})}. Against $0$ of $5{,}272$ for the sampled
search, the construction passes at every depth from $8$ to $12$; depth $13$ needs modulus
$3.0\times10^{14}$ and more swap pairs than the $2^{24}$ table cap admits.

\begin{remark}\label{rem:coprimeverdict}
The two routes are worth separating. The \emph{factorisation} obstacle is genuinely gone --- both
tests run in arithmetic modulo $10^{10}$ and never form $\beta$ --- and
Proposition~\ref{prop:coprimesieve} is a sound, essentially free first pass that any future search
should run. What replaces it is a harder obstacle: inside the mass window the candidates are not
merely hard to test, they are uniformly far from the target, by a margin of $0.77$ in $y_{\max}$ at
the observed best. A search for a $1$--free coprime example must therefore generate candidates that
are smooth \emph{by construction} rather than filtered for smoothness after the fact.
Section~\ref{ssec:smoothbuilt} gives such a construction. What it does not give is a test of
sufficiency: deciding whether some coprime factorisation of $\beta$ sums to exactly $\alpha/\beta$
needs the factorisation of $\beta$, whose unfactored part in the verified instance has $13{,}151$
digits. Factoring returns, but only at the last step, and it no longer blocks the generation of
candidates.
\end{remark}

\subsection{The two--member side, and what every known example really is}\label{ssec:twoside}

The searches above treat both sides symmetrically. They need not be. Suppose one side has exactly
two members, $A=\{u,v\}$. Then $\alpha=uv$ and $\beta=u+v$, so $u$ and $v$ are the roots of
$z^2-\beta z+\alpha=0$: \emph{$A$ is recovered from $(\alpha,\beta)$ by a quadratic, with no
factorisation at all}. The roots are automatically coprime, since $\gcd(u,v)$ divides
$\gcd(u+v,uv)=\gcd(\beta,\alpha)=1$.

\begin{proposition}[The two--member reduction]\label{prop:twoside}
Let $B$ be a coprime family, $\beta=\prod B$, and let $u\ge1$ be an integer with $u<\beta$. Then
$\{u,\beta-u\}$ and $B$ satisfy $\bigl(\sum_A1/a\bigr)\bigl(\sum_B1/b\bigr)=1$ if and only if
\[ D(B)\;=\;u(\beta-u),\qquad\text{that is}\qquad \sum_{b\in B}\frac1b=\frac{u(\beta-u)}{\beta}. \]
$0$ \texttt{sorry} \textup{(\texttt{Erdos307.Coprime.two\_element\_identity})}.
\end{proposition}
\begin{proof}
$\sum_{a\in A}1/a=1/u+1/(\beta-u)=\beta/\bigl(u(\beta-u)\bigr)$, and $\sum_B1/b=D(B)/\beta$. The
product is $D(B)/\bigl(u(\beta-u)\bigr)$, which is $1$ exactly when $D(B)=u(\beta-u)$.
\end{proof}

The value of $u$ decides everything, and it identifies the known examples.

\emph{$u=1$.} The condition is $D(B)=\beta-1$, which is precisely the \emph{primary pseudoperfect}
condition, and $A=\{1,\beta-1\}$. This is $(1+\frac15)(\frac12+\frac13)$ from $\beta=6$ and
$(1+\frac1{1805})(\frac12+\frac13+\frac17+\frac1{43})$ from $\beta=1806$; the Sylvester products
$2,6,42,1806$ are pseudoperfect because Sylvester's recursion maintains $\beta-D(B)=1$ identically.
So the whole recorded family of solutions is the case $u=1$, and it contains the member $1$ for the
structural reason that $u$ \emph{is} that member.

This is an equivalence, not merely a supply of examples, and it links two of Erd\H{o}s's problems.
That a primary pseudoperfect number yields a solution of the relaxation is an observation of Cambie
on the problem page; Proposition~\ref{prop:twoside} supplies the converse, since it derives
$D(B)=u(\beta-u)$ as a \emph{criterion}.

\begin{corollary}[\#307's relaxation at $u=1$ is \#313]\label{cor:threethirteen}
The two--member solutions of the coprime relaxation of \#307 that contain the element $1$ are exactly
the primary pseudoperfect numbers: $\bigl(\{1,\beta-1\},B\bigr)$ solves the relaxation if and only if
$\beta=\prod B$ satisfies $\sum_{q\mid\beta}1/q+1/\beta=1$. Consequently Erd\H{o}s Problem \#313,
the infinitude of primary pseudoperfect numbers \cite{ButskeJajeMayernik}, is equivalent to the
statement that the relaxation of \#307 has infinitely many two--member solutions containing $1$.
\end{corollary}
\begin{proof}
Proposition~\ref{prop:twoside} at $u=1$ reads $D(B)=\beta-1$, and dividing by $\beta$ gives
$\sum_{q\mid\beta}1/q=1-1/\beta$, which is the primary pseudoperfect condition. Both directions are
the same computation.
\end{proof}

The identity $D(B)=u(\beta-u)$ also names the right generalisation. Writing $\partial$ for the
arithmetic derivative, \cite{PortFillings2026} calls $(R,c)$ a \emph{port} and says a squarefree $B$
\emph{fills} it when $cB-R\,\partial(B)=1$; the case $R=1$, $c=1$ is \#313. Our condition is
$u\beta-\partial(\beta)=u^{2}$, that is, the port $(1,u)$ filled at value $u^{2}$ rather than $1$.
The $1$--free case of the relaxation is therefore a port--filling problem one step outside the range
that \#313 occupies, which is where any machinery developed for \#313 would have to be pushed.

\begin{proposition}[Only the value one sustains itself]\label{prop:portfixed}
For squarefree $B$ and $c\ge1$ put $\Delta_c(B)=cB-\partial(B)$. Then $\gcd(\Delta_c(B),B)=1$, and
appending a prime $q\nmid B$ gives $\Delta_c(qB)=q\,\Delta_c(B)-B$. Consequently a step returns to
the same value, $\Delta_c(qB)=\Delta_c(B)$, only when $\Delta_c(B)=1$, and then $q=B+1$. $0$
\texttt{sorry} \textup{(\texttt{Erdos307.port\_fixed\_value})}.
\end{proposition}
\begin{proof}
A prime $p\mid B$ dividing $\Delta_c(B)$ would divide $\partial(B)$, contradicting
Theorem~\ref{thm:structure}; so $\gcd(\Delta_c(B),B)=1$. The recursion is
$\partial(qB)=B+q\,\partial(B)$. Fixing the value gives $\Delta_c(B)(q-1)=B$, so $\Delta_c(B)\mid B$,
and coprimality forces $\Delta_c(B)=1$.
\end{proof}

\begin{remark}[Why \#313 has a family and the $1$--free case has none]\label{rem:portvalue}
Proposition~\ref{prop:portfixed} explains the asymmetry between
Corollary~\ref{cor:threethirteen}'s two sides. At $u=1$ the target value is $1$, which is a fixed
point of the port step: from $\Delta=1$ the choice $q=B+1$ returns $\Delta=1$, and iterating from
$B=2$ gives $2,6,42,1806,\dots$, Sylvester's sequence, whose partial products are exactly the known
primary pseudoperfect numbers. The family costs nothing because the recursion sits still.

At $u\ge2$ the target is $u^{2}$, and by the proposition \emph{no} value other than $1$ is
self-sustaining, so there is no orbit to sit on: every step must be steered, and the deficit
$\Delta$ wanders. This is the structural content of the failed descent recorded earlier --- over
$399$ seeds the recursion $\Delta\mapsto q\Delta-B$ never reached a divisor of $B+4$, the smallest
$\Delta$ attained being $2.7\times10^{6}$ \textup{(\texttt{code/two\_side\_descent.gp})}. The search
was not unlucky. Verified: $23{,}916$ pairs $(B,c)$ with $\gcd(\Delta_c(B),B)=1$ and $0$ failures,
$14{,}920$ checks of the recursion with $0$ failures, and among all self-sustaining steps found in a
sweep, every one had $\Delta=1$ \textup{(\texttt{code/port\_recursion.gp})}.

The consequence for transfer is specific. Machinery built for \#313 targets the value $1$ and may
use its stability; anything of that kind does not carry to $u\ge2$, where the stability is provably
absent. What does carry is the recursion itself, which is value--blind.

Proposition~\ref{prop:portfixed} is sharp at period one only. A two--step orbit
$\Delta\mapsto q_1\Delta-B\mapsto\Delta$ is possible above the value $1$, and requires exactly
$\Delta\mid q_1+q_2$ \textup{(}$0$ \texttt{sorry}, \texttt{Erdos307.port\_period\_two}\textup{)};
at $c=1$ such orbits occur, for instance $B=33$ with $\Delta:19\mapsto5\mapsto19$ at
$q_1=2$, $q_2=17$. They are single events rather than families, since returning a second time needs
a fresh pair: the condition ties the primes to $B$. In the regime that the $1$--free case actually
requires --- $c=2$ with $B$ and both primes odd --- a sweep of $162{,}110$ states found none
\textup{(\texttt{code/port\_orbit.gp})}.
\end{remark}

\begin{proposition}[The barrier survives relaxing the equations to divisibilities]\label{prop:divrelax}
Let $P,Q$ be disjoint finite sets of primes and suppose only that
\[ \textstyle\prod Q\ \bigm|\ D(P)\qquad\text{and}\qquad \prod P\ \bigm|\ D(Q), \]
in place of the two equalities a two--cycle requires. Then
$\sigma(P)\sigma(Q)\ge1$, hence $\sigma(P)+\sigma(Q)\ge2$, and therefore $|P\cup Q|\ge60$ and
$\min(\prod P,\prod Q)\ge2.09\times10^{56}$ as in Theorem~\ref{thm:barrier}. $0$ \texttt{sorry}
\textup{(\texttt{Erdos307.mass\_from\_divisibility})}.
\end{proposition}
\begin{proof}
Divisibility between positive integers gives $\prod Q\le D(P)=\sigma(P)\prod P$ and
$\prod P\le D(Q)=\sigma(Q)\prod Q$. Multiplying and cancelling the positive quantity
$\prod P\prod Q$ leaves $1\le\sigma(P)\sigma(Q)$, and the arithmetic--geometric mean inequality
gives $\sigma(P)+\sigma(Q)\ge2$, which is the only input Theorem~\ref{thm:barrier} consumes.
\end{proof}

\begin{remark}[What the relaxation localises]\label{rem:divrelax}
The relaxation is worth recording because it is natural to expect the mass condition to be an
artefact of the \emph{exact} equations, so that weakening them would open room below level $60$. It
does not: the divisibilities alone already force the mass, because a divisibility between positive
integers is an inequality, and the two inequalities point in opposite directions. Consequently the
difficulty in \#307 is not located in the passage from ``$\prod Q$ divides $D(P)$'' to
``$\prod Q$ equals $D(P)$''; both carry the same barrier.

At $|P|=|Q|=2$ the relaxed system has no solutions at all among the primes below $283$
\textup{(}$2{,}925{,}810$ pairs, $0$ solutions, \texttt{code/relax\_divisibility.gp}\textup{)}, and
Proposition~\ref{prop:divrelax} says why: four primes cannot carry mass $2$. A random search for
relaxed solutions is therefore vacuous rather than evidential, and the proposition is checked
instead on the arithmetic it actually asserts, over $200{,}000$ random positive quadruples with $0$
counterexamples.
\end{remark}

\begin{corollary}[The balanced splits are exactly the forbidden ones]\label{cor:splitbalance}
In the situation of Corollary~\ref{cor:level61height}, a split $S=P_0\sqcup Q_0$ can support a
solution only if $\sigma(P_0)\sigma(Q_0)<1$, hence only if $\sigma(P_0)$ lies strictly \emph{outside}
the base's own mass window $[1/t_0,t_0]$, where $t_0+1/t_0=\sigma(S)$. In particular the
mass--balanced splits, $\sigma(P_0)\approx\sigma(S)/2$, are precisely excluded. $0$ \texttt{sorry}
\textup{(\texttt{Erdos307.split\_outside\_window})}.
\end{corollary}
\begin{proof}
Theorem~\ref{thm:frame} gives $\alpha=(MN'+N^2)/\Delta$ with $\Delta=D(1-\sigma(P_0)\sigma(Q_0))$, and
$\alpha>0$ forces $\Delta>0$. With $\sigma(P_0)+\sigma(Q_0)=\sigma(S)$ fixed, the product is the
downward parabola $z\mapsto z(\sigma(S)-z)$, which exceeds $1$ exactly between the roots of
$z^2-\sigma(S)z+1$ and is maximal at the midpoint.
\end{proof}

\begin{remark}[Why level $61$ is still not enumerable]\label{rem:level61search}
This is the complement of Proposition~\ref{C-prop:window}, which places the mass of a \emph{solution}
inside the window of the full support; here the split of the \emph{base} must sit outside its own.
It is a genuine constraint and it points a search away from where one would naturally begin, but it
does not make level $61$ enumerable, and it is worth recording the arithmetic rather than leaving the
size vague. Sampling $2\times10^{6}$ splits of the first $59$ primes, $89.4\%$ satisfy
$\sigma(P_0)\sigma(Q_0)<1$, so the condition removes about one split in nine
\textup{(\texttt{code/split\_balance.gp})}. Against $49{,}961$ bases and $2^{59}$ splits each, some
$2.6\times10^{22}$ candidates remain, and the product--balance $M\approx N\approx\sqrt D$ that
Theorem~\ref{thm:frame} prefers trims this by a further factor of order $10$, not of order $10^{10}$.
Each surviving candidate still needs two divisibility tests on $113$--digit integers. We therefore
record level $61$ as out of computational reach by roughly ten orders of magnitude, and note that no
constraint in this paper closes that gap.
\end{remark}

\begin{remark}[What actually obstructs, in both problems]\label{rem:portobstruction}
It is worth being exact about where the difficulty sits, because the port dynamics is not where it
sits. Along Sylvester's chain the value is $\Delta\equiv1$ \emph{identically}, so for \#313 the value
dynamics never obstructs at all; what obstructs is that the next element $B/\Delta+1$ must be
\emph{prime}, and the chain $2,6,42,1806$ stops because $1807=13\cdot139$. That single composite is
the whole reason \#313 is open rather than settled by Sylvester.

The comparison is then clean, and it is the honest conclusion of the transfer. \#313 carries one
obstruction: the primality of a constructed element. The $1$--free case of \#307's relaxation
carries that same obstruction \emph{and}, by Proposition~\ref{prop:portfixed}, a second one that
\#313 does not have --- no self--sustaining value to run a chain along in the first place. So the two
problems are not of equal difficulty under this correspondence: \#307's relaxation is strictly harder
in an identifiable way, and a solution of \#313 would not by itself supply one. This is a negative
conclusion, but it is a proved one rather than a failed search, and it retires the transfer as a
route.
\end{remark}

\emph{$u=2$.} The first $1$--free case. The condition is $D(B)=2\beta-4$, i.e.\
$\sum_{b\in B}1/b=2-4/\beta$, and then $A=\{2,\beta-2\}$, which needs $\beta$ odd so that
$\gcd(2,\beta-2)=1$ \textup{(}$0$ \texttt{sorry}, \texttt{Erdos307.Coprime.two\_free\_case}\textup{)}.
Reducing $D(B)=2\beta-4$ modulo each $b\in B$ gives the Znám--type system
$\prod_{b'\ne b}b'\equiv-4\pmod b$, and conversely those congruences force
$\beta\mid D(B)-2\beta+4$, so with $1<\sum_B1/b<2$ the quotient is an integer in $(-1,0]$ and must
vanish: \emph{the local conditions imply the global one}.

\begin{proposition}[$u$ is the mass, so $u=1$ is the only cheap case]\label{prop:umass}
In the situation of Proposition~\ref{prop:twoside}, $\sum_{b\in B}1/b=u(\beta-u)/\beta=u-u^2/\beta$.
Hence the reciprocal sum of $B$ is $u$ up to $O(u^2/\beta)$: the parameter $u$ \emph{is} the mass
$B$ must carry. Consequently the cost is monotone in $u$, and $u=1$ --- the case that contains the
member $1$ --- is the only value for which $B$ needs mass $\approx1$ and can therefore be short.
\end{proposition}

This settles what one might hope for from larger $u$. Since $\sum_A1/a=1/u+1/(\beta-u)$ and
$\sum_B1/b$ is its reciprocal, asking for $1\notin P\cup Q$ is asking for $u\ge2$, which is asking
for $B$ to carry mass at least $2-4/\beta$; $u=3$ demands $3-9/\beta$, and so on. There is no
cheaper $u$ hiding above $2$, and no analogue of Sylvester's invariant to be had there: the
invariant $\beta-D(B)=1$ works at $u=1$ precisely because mass $1$ is reachable in four terms.
Verified exhaustively: over every $\beta\le4\times10^{7}$ and \emph{every} coprime factorisation of
it --- all set partitions of its prime powers --- the discriminant $\beta^2-4D(B)$ is a square with
admissible root in exactly $7$ cases, and every one has $u=1$. The largest $u$ attained anywhere in
that range is $1$ \textup{(\texttt{code/two\_member\_scan.rs})}.

\begin{remark}[The price of dropping the member $1$]\label{rem:twosidecost}
The reduction is exact and factoring--free, but it fixes the size. Since
$\sum_B1/b=2-4/\beta$ and $B$ consists of odd pairwise coprime integers, the reciprocal sum must
exceed $2$, and $\sum_{3\le p\le X}1/p$ first exceeds $2$ at the $1412$nd odd prime $p=11789$
\textup{(}sum $2.0000206\ldots$\textup{)}. Hence $|A\cup B|\ge1414$, against $60$ for the prime
problem: passing from $u=1$ to $u=2$ costs a factor of $23$ in the number of members, because the
member $1$ was carrying a full unit of mass that must now be found among odd reciprocals.
This is the exact price of the gap on the problem page, and it is why the Sylvester descent, which
produces the $u=1$ family for free, does not reach $u=2$: driving $e=2\prod B-D(B)$ to a divisor of
$\beta+4$ by the recursion $e\mapsto be-\prod B$ failed over $399$ seeds, the smallest $e$ reached
being $2.7\times10^{6}$, while the products outrun any bound long before $1412$ members
\textup{(\texttt{code/two\_side\_descent.gp})}.
\end{remark}

\section{Status and questions}\label{sec:status}

Erd\H{o}s~\#307 remains \textbf{open}. It is Conjecture~4 of Ufnarovski--\AA hlander
\cite{UfnarovskiAhlander2003}, so it is equivalent to the existence of a non--trivial two--cycle of
the arithmetic derivative and has been open in that literature since $2003$ under another name. We
establish rigidity and a barrier, $|P\cup Q|\ge60$ with $\min(\prod P,\prod Q)>3.50\times10^{57}$,
which voids all direct search; it appears to be the first valid unconditional bound of its kind,
the only prior candidate resting on an invalid step (Section~\ref{sec:prior}). The rigidity result
and the barrier are Lean--verified, extremality proved rather than assumed.

The barrier is one instance of a single inequality. Writing $n(c)=\min\{n:T_n\ge c\}$,
\[ \sum_{p\in U}1/p\ \ge\ c\ \Longrightarrow\ |U|\ge n(c),\qquad \prod_{p\in U}p\ \ge\ \Pi_{n(c)}, \]
and the arithmetic of each result enters only in producing its $c$ (Remark~\ref{C-rem:dictionary}).
Read in the size it gives the double--logarithmic frontier law
(Proposition~\ref{C-prop:frontier}); in the near--miss ratio, the cost family
(Corollary~\ref{C-cor:cost}), whose endpoint at $r=1$ is the barrier and which shows $a''\ge2a$ needs
$37{,}963$ primes; in the period, that the barrier is uniform in $k$, so it bounds the whole of
Ufnarovski--\AA hlander's Conjecture~3 and not only \#307 (Proposition~\ref{prop:kuniform}); in the
orbit length, that the derivative cannot outgrow its own support
(Proposition~\ref{C-prop:growth}); in the split, that $\min(\omega(a),\omega(b))\ge3$ whenever
$|U|\le68$ (Proposition~\ref{C-prop:split}); and in the level, the forced core and the ceiling that
pin level $59$ from both sides and yield its admissible count $49{,}961$ in closed characterisation
(Propositions~\ref{C-prop:window},~\ref{C-prop:band}, Corollary~\ref{C-cor:count59}). The algebraic cores
of these are Lean--verified (\texttt{Dictionary.lean}).

On the constructive side we give a Th\=abit/Euler--type calculus and a bilinear breeder whose
candidate $q$ is always integral (Proposition~\ref{C-prop:bdeint}), a conditional reduction
$(\mathrm{H})\Rightarrow$\textsc{yes}, and a quantitative ledger showing $(\mathrm{H})$ is not a
feasible finite search. No fixed--shape family of any growth rate can parametrise \#307
(Proposition~\ref{C-prop:nofamily}), which closes the Th\=abit/Euler programme rather than only its
polynomial shadow. We also give a parity dichotomy, an Erd\H{o}s--Wintner density theorem, and a
function--field theorem localising the obstruction to the archimedean place. The heuristics favour
\textsc{yes}, and the model behind them survives a second--moment test
(Proposition~\ref{C-prop:anticorr}); we claim no proof in either direction.

Open problems: (i) close the residual $3{,}478$ single--tail families at level $60$, and close the
pair sector of Proposition~\ref{prop:sectors}. What this buys should be read with
Proposition~\ref{C-prop:nottree}: the admissible supports do not form a tree under adjunction, so
completing every one--new--prime family at level $60$ would still not close level $60$; for which Proposition~\ref{prop:pairlocal} rules
out every congruence kill, so only global, per--family input (factorisation or direct search) can
serve (the canonical tail family
is closed, Proposition~\ref{prop:tailkill}; the finite verification of
Proposition~\ref{prop:close59} is Lean--checked, \texttt{erdos307\_sixty}; the conditional $3/4$
law is explained; it is the reciprocity switch of Remark~\ref{rem:fiveeighths}, made exact by
the class--vector functional there); (ii) decide
the parity dichotomy's all--odd case (we expect it empty); (iii) sharpen the certified enclosure
$0.04202\le\delta\le0.04223$ of Proposition~\ref{C-prop:density}, whose width is carried almost
entirely by the tail window $a$; three digits are certified, and a longer envelope convolution
with a smaller window certifies more; (iv) settle the density hypothesis $(\mathrm{H})$ of
Section~\ref{C-sec:breeder}, or exhibit a non--circular frame family escaping the $\tau(K)/G$ ceiling;
(v) find a $\mathbb{Z}$--side structure detecting the archimedean obstruction of
Section~\ref{C-sec:ff}; one that must be non--local, since by Proposition~\ref{C-prop:nolocal} no single
modulus obstructs; (vi) settle \#307.

\begin{remark}[The place of the problem: a local--global dictionary]\label{rem:place}
The pieces above assemble into a dictionary. \textup{(i)} For a fixed support $U$, the subset form of
Proposition~\ref{C-prop:subset}, $A(\Sigma-A)=D^2$, satisfies the \emph{Hasse principle as an equation}:
an integer solution $A$ exists iff one exists over $\mathbb{R}$ and over every $\mathbb{Z}_\ell$ (the
discriminant $\Sigma^2-4D^2$ must be a nonnegative square, and squares obey local--global). Its
real--place condition is $\Sigma\ge2D$, i.e.\ $\sigma(U)\ge2$: \emph{the barrier is precisely the
real--place solvability condition of the subset form}. \textup{(ii)} The barrier never needed
primality: pairwise coprime integers have distinct smallest prime factors, so
$\sum 1/a_i\le\sum1/p_i$, and \emph{any} pairwise--coprime support of mass $\ge2$ has at least $59$
elements and product at least $\Pi_{59}$. Coprimality and the ordering of $\mathbb{R}$ carry the wall;
primality enters only through rigidity. \textup{(iii)} What remains unsolved is the \emph{realization}
of the root $A^*=\bigl(\Sigma+\sqrt{\Sigma^2-4D^2}\bigr)/2$ as a sub--multiset sum of the cofactors
$D/p$. Its local conditions are exactly the local anatomy of Lemma~\ref{C-lem:anatomy}: modulo $\ell^2$
for $\ell\in U$ the subset sums occupy precisely the two cosets
$\ell\mathbb{Z}\cup(D/\ell+\ell\mathbb{Z})$ ($2\ell$ residues, observed exactly ($6$ of $9$, $10$ of
$25$, $14$ of $49$, $22$ of $121$)) and these are the same constraints the equation itself forces:
no completion separates the equation from its realization. In the language of integral points, then:
\#307 is everywhere locally soluble, with the barrier as its archimedean height condition, and its
resolution is a \emph{strong approximation} question for the thin, prime--structured set of subset sums
a failure mode invisible to Brauer--Manin--type obstructions, consistent with every no--go above
and with the existence of cycles over the neighbouring rings. We record this as a dictionary, not a
method.
\end{remark}

\paragraph{Attribution.}
We are careful to separate prior, concurrent, and present contributions (details in
Section~\ref{sec:prior}). \emph{Prior/concurrent:} the two--cycle reformulation, and its
equivalence with the arithmetic derivative, are due to Ufnarovski--\AA hlander
\cite{UfnarovskiAhlander2003} (2003, Conjecture~4); the naming of it as \#307 to Seva
\cite{Seva323194} (2019); an independent rederivation to Bado \cite{Bado2026}; the bound $|P\cup Q|\ge59$ to
J.~Rosen and the disjointness valuation $v_p(\sum1/p_i)=-1$ to R.~Israel \cite{MO320838}; the source
identification to ``Woett'' \cite{MO320838}. Bado \cite{Bado2026} independently obtained the rigidity
cross--matching, the local anatomy, the parity restriction, a one--sided exact test, approximate
solvability, and the union criterion of Proposition~\ref{prop:pyth}. \emph{Present note:} the identification of \#307 with
Ufnarovski--\AA hlander's Conjecture~4 and the resulting two--way bridge to the $n''=n$ literature; the product barrier $2.09\times10^{56}$; the
function--field theorem and archimedean localisation; the constructive deficit analysis and seed
obstruction; the Erd\H{o}s--Wintner density theorem, and the $k$--cycle barriers. Both external
sources \cite{Seva323194,Bado2026} have been examined directly and the overlaps above reflect their
actual contents.

\paragraph{Acknowledgments.}
This work was carried out by the author with substantial assistance from Anthropic's Claude. The
AI assistant contributed to the derivations, the exact and heuristic computations, the Lean~4
formalisation, and the drafting of this note; it was also used to stress-test and attempt to
refute each claim. All statements, proofs, attributions, and the final text were reviewed and are
the responsibility of the author. The author thanks the contributors to the MathOverflow thread
\cite{MO320838}; in particular R.~Israel and J.~Rosen, whose observations are credited above.

\begin{remark}[Where a proof of \textsc{no} would have to come from]\label{rem:whereleft}
Four results of this note combine into a single statement about the shape any negative proof must
take, and it is worth assembling because each of the four is easy to read as an isolated obstruction.

Proposition~\ref{prop:localcomplete} shows the cycle system is soluble modulo every $m$, so it has
solutions in $\widehat{\mathbb{Z}}$ and there is no congruence obstruction to find, at any modulus or
any finite combination of moduli. Theorem~\ref{C-thm:ff} shows the analogue over $\mathbb{F}_q[t]$ is
false, having no cycles at all, so no argument that survives the passage to function fields can
decide \#307 negatively: whatever a proof uses, $\mathbb{F}_q[t]$ must lack it.
Proposition~\ref{C-prop:noplace} shows no place carries a descent. Together these localise the
obstruction at the archimedean place, which is the reading already recorded.

The fourth is the constraint. The only archimedean tool in play is the mass, and
Proposition~\ref{prop:barthreshold} shows it is exactly tight: raising the barrier one level requires
$\sigma(a)>t_n$, the only unconditional bound available is $\sigma(a)-1\ge1/a$, and
$\{2,3,7,41\}$ realises $\sigma=1.00058\ldots$ on four primes, so nothing forbids $a'-a=1$. So a
proof of \textsc{no} needs an archimedean argument that is not the mass, and we know of none. This is
not a claim that none exists; it is a statement of what the four results jointly exclude, recorded so
that the search is not repeated inside the region they cover.
\end{remark}

\section{Prior and concurrent work}\label{sec:prior}

We record the overlaps with four items, each examined directly.

\paragraph{Ufnarovski--\AA hlander (2003) \cite{UfnarovskiAhlander2003}.}
Their Theorem~5 gives $p^{k}\|n$, $0<k<p$ $\Rightarrow$ $p^{k-1}\|n'$, and their Theorem~4 gives
$n^{(j)}\to\infty$ when $p^{p}\mid n$ properly; together with Corollary~2 ($n$ squarefree
$\Leftrightarrow$ $(n,n')=1$) these force a period--two point to have squarefree members with
disjoint supports. The argument is sound: for $2\le k<p$ two applications of Theorem~5 give
$v_p(n'')=k-2\ne k$, while $k\ge p$ falls to Theorem~4 or to the fixed point $n=p^p$, and the edge
cases $p=2$, $k\in\{2,3\}$ are covered by the latter because $4=p^p$. They then state the
period--two case as \textbf{Conjecture~4}: $\bigl(\sum_{i\le k}1/p_i\bigr)\bigl(\sum_{j\le l}
1/q_j\bigr)=1$ has no solution in distinct primes. This is \#307, and priority for the two--cycle
reformulation belongs here rather than to 2019: what \cite{Seva323194} adds is the identification of
that equation as \#307, which \cite{UfnarovskiAhlander2003} does not make --- they cite
\cite{Barbeau1961} for the derivative, not \cite{Barbeau1976} for the problem. The consequence is
that \#307 has been an open conjecture in the arithmetic--derivative literature since 2003 under
another name, and that the bounds of Section~\ref{sec:barrier} are bounds on Conjecture~4.

\paragraph{Seva (MathOverflow 323194, 2019) \cite{Seva323194}.}
Defines $s(Q)=Q\sum_{p\mid Q}1/p$, derives the valuation property $v_p(s(Q))=v_p(Q)-1$, asks for
$s(P)=Q$, $s(Q)=P$, and states this is Problem~\#307 ``in disguise.'' Priority for \emph{naming}
the two--cycle question as \#307 belongs here. He further asks whether every $s$--trajectory stabilises at $0$; the
$s$--analogue of the Ufnarovski--\AA hlander conjecture, posed independently of it. (Note $s\ne n'$
globally, e.g.\ $s(p^\alpha)=p^{\alpha-1}$, but $s=n'$ on squarefree integers, where both two--cycle
questions live.)

\paragraph{Kovi\v{c} (2012), Propositions~17 and~18 \cite{Kovic2012}.}
Proposition~18 claims $r+s\ge34$ unconditionally ($\ge57$ if $\min\{p_1,q_1\}\ge3$, $\ge110$ if
$\ge5$), together with $p_1<r$, $q_1<s$ and $\max\{m,n\}\ge\prod_{i\le17}p_i\approx1.92\times
10^{21}$. Its proof does not establish them. Substituting $n'=m$, $m'=n$ into its two correct pairs
of bounds gives
\[ \text{(A)}\ \ \frac{rn}{p_r}<m<\frac{rn}{p_1}, \qquad
   \text{(B)}\ \ \frac{sm}{q_s}<n<\frac{sm}{q_1}, \]
after which it asserts ``Hence: $p_1q_s<rs$ and $q_1p_r<rs$'', and deduces everything from these via
$2\max\{p_r,q_s\}<N^2/4$ against $\max\{p_r,q_s\}\ge P_N$. But products of inequalities are valid
only in matching directions, and (A) with (B) admits exactly two: upper$\,\times\,$upper gives
$p_1q_1<rs$, lower$\,\times\,$lower gives $rs<p_rq_s$. The asserted $p_1q_s<rs$ pairs an upper bound
on $m$, carrying $p_1$, with a lower bound on $n$, carrying $q_s$. Nor is it recoverable: the cycle
says precisely $\sigma(n)\in(r/p_r,\,r/p_1)$ and $\sigma(m)=1/\sigma(n)\in(s/q_s,\,s/q_1)$, so
$\sigma(n)$ lies in $(r/p_r,\,r/p_1)\cap(q_1/s,\,q_s/s)$; nonemptiness of that intersection is
equivalent to those two products and to nothing further, whereas $p_1q_s<rs$ reads $q_s/s<r/p_1$, a
comparison of the two \emph{upper} endpoints. In
$n=2\cdot5\cdot13\cdot37\cdot53\cdot73\cdot79\cdot89\cdot97\cdot101$,
$m=3\cdot7\cdot19\cdot31\cdot43\cdot59\cdot67\cdot71\cdot83\cdot103$ --- squarefree, disjoint
supports, $r=s=10$ --- all four inequalities hold while $p_1q_s=206>100=rs$ and $q_1p_r=303>100$
(\texttt{code/kovic\_prop18\_audit.gp}; the valid products and the invalidity are machine--checked in
\texttt{lean/Erdos307/KovicProp18.lean}). Two arithmetic slips sit alongside: Table~2 lists
$P_{11}=27$, and the ``product of the first $17$ primes'' is printed with $27$ for $29$.

What fails is the \emph{inference}, not the \emph{statement}: no two--cycle is known, so no
counterexample to $r+s\ge34$ exists to be given, and none is claimed. The conclusion is unproven
rather than false, and is in fact true, since Theorem~\ref{thm:barrier} gives $|U|\ge60$ by a route
touching none of this. Propositions~16 and~19 are sound, and so is~17, but~17 assumes the
\emph{smaller} member is odd and \cite{Kovic2012} notes that no estimate follows when it is even. The
strongest valid unconditional bound in the prior literature is therefore that paper's computer search
$x>10^4$; Theorem~\ref{thm:barrier} improves it by $52$ orders of magnitude, and $|U|\ge60$ is the
first valid cardinality bound for arithmetic--derivative two--cycles rather than an improvement on an
earlier one. None of this diminishes \cite{Kovic2012}, which remains a standard reference and whose
Propositions~16, 17 and~19 we use.

\paragraph{Bado (preprint, Aix--Marseille, 2026) \cite{Bado2026}.}
Working with $A(S)=\sum_{p\in S}M(S)/p$ and $M(S)=\prod_{p\in S}p$ (our $N_S,D_S$), Bado independently
obtains: the lowest--terms lemma (his Lemma~3.1, our Lemma~\ref{thm:rigidity}); the forcing identities
$A(P)=M(Q)$, $A(Q)=M(P)$ with $P\cap Q=\varnothing$ (his Thm.~A, our Theorem~\ref{thm:structure}); the
two--cycle reformulation $T(T(P))=P$ (his Thm.~B); the valuation signature (his Thm.~6.1, = Israel's);
the local zero--sum congruences (his Prop.~8.1, our Lemma~\ref{C-lem:anatomy}); the parity restrictions
(his Prop.~8.2, the parity part of Theorem~\ref{thm:parity}); an exact one--sided test (his Thm.~9.1); the deviation parameter, i.e.\ the integer $t$ with
$A(P)=n+tm$, $A(Q)=m-tn$ and solubility exactly at $t=0$ (his Thm.~7.6, our rung $k$ of
Proposition~\ref{C-prop:ladder}); the intrinsic union criterion (his Thm.~7.3, our
Proposition~\ref{prop:pyth} in \emph{iff} form); uniqueness of the splitting (his Cor.~7.5);
the approximation theorem (his Thm.~10.2, sharper than our greedy version in that all primes exceed a
given bound), and the union discriminant / Pythagorean structure (his Thms.~7.1, 7.4). He cites only
Barbeau, Bloom, and Erd\H{o}s--Graham: crucially, he does \emph{not} identify the map with the
arithmetic derivative, and his note contains none of the present note's $n''=n$ bridge, explicit
barrier constant, function--field theorem, density theorem, $k$--cycle bounds, or constructive
deficit analysis. We regard the overlap on the shared structural layer as strong independent
validation, and import his union criterion with credit. Bado has confirmed this account in
correspondence (August 2026): he agrees that the results listed above were present in his
preprint and that the overlap is concurrent and independent, and he has confirmed the
identification $A(S)=M(S)'$, which he had not made, and is revising his historical discussion
accordingly. The concurrency is therefore recorded on both sides rather than asserted on one:

\begin{proposition}[Single--integer test; derivative form of Bado's Thms.~7.1, 7.4]\label{prop:pyth}
If $(a,b)$ is a two--cycle of the arithmetic derivative, then $N=ab$ is squarefree and satisfies
$N'=a^2+b^2$, whence $N'^2-4N^2=(a^2-b^2)^2$ is a perfect square. Conversely, a squarefree $N$ with
$N'^2-4N^2=\square$ determines a unique unordered factorisation $N=ab$ with $N'=a^2+b^2$, on which the
cycle condition reduces to a single equation (Bado's deviation parameter $t=0$, his Thm.~7.6). This
gives a one--integer necessary test;
computationally, no squarefree $N<2\times10^5$ with $\omega(N)\ge2$ passes it ($0$ of $103{,}596$,
consistent with Theorem~\ref{thm:barrier}, which forces $N>4\times10^{112}$).
\end{proposition}
\begin{proof}
$N=ab$ with $a,b$ coprime squarefree gives $N$ squarefree and, by Leibniz with $a'=b,\,b'=a$,
$N'=a'b+ab'=b^2+a^2$; then $N'^2-4N^2=(a^2+b^2)^2-4(ab)^2=(a^2-b^2)^2$. The converse is Bado's
\cite{Bado2026}; the emptiness of the test below $10^{112}$ is upgraded from a sieve to a theorem in
Corollary~\ref{cor:emptytest}.
\end{proof}

The Pythagorean equation $N'=a^2+b^2$ deserves to be isolated from the (harder) cycle condition.

\begin{proposition}[Pair form; the one--equation relaxation]\label{prop:pairform}
For coprime squarefree $a,b$, the single equation $a^2+b^2=(ab)'$ holds if and only if there is an
integer $k$ with $a'=b+ak$ and $b'=a-bk$; a two--cycle is exactly the case $k=0$. We call a solution
of $a^2+b^2=(ab)'$ (any $k$) a \emph{Pythagorean pair}.
\end{proposition}
\begin{proof}
The Leibniz identity $a^2+b^2-(ab)'=a(a-b')-b(a'-b)$ (a polynomial identity, verified symbolically
and on $2000$ random pairs) makes the equation equivalent to $a(a-b')=b(a'-b)$; since $\gcd(a,b)=1$
this forces $a\mid a'-b$ and $b\mid a-b'$ with equal quotients $k$, i.e.\ $a'=b+ak$, $b'=a-bk$.
Setting $k=0$ recovers $a'=b$, $b'=a$. (Equivalently, in $\mathbb{Z}[i]$: \#307 asks for $z=a+bi$ with
coprime coordinates and $ab$ squarefree satisfying $|z|^2=(\operatorname{Im}z^2/2)'$ and $k=0$.)
\end{proof}

\begin{corollary}[The single--integer test is provably empty below $10^{112}$]\label{cor:emptytest}
Every Pythagorean pair satisfies $\sum_{p\mid a}1/p+\sum_{q\mid b}1/q=a/b+b/a\ge2$ (divide
$a^2+b^2=(ab)'$ by $ab$ and apply AM--GM), so its support carries at least $59$ primes and
$N=ab\ge\prod(\text{first }59\text{ primes})>7.9\times10^{112}$. Hence no squarefree $N<10^{112}$ with
$\omega(N)\ge2$ satisfies $N'^2-4N^2=\square$: the empirical ``$0$ of $103{,}596$'' of
Proposition~\ref{prop:pyth} is a theorem, and the barrier of Theorem~\ref{thm:barrier} extends from
two--cycles to the strictly weaker one--equation (Pythagorean) problem.

\textup{Formalised.} \texttt{Erdos307.pyth\_sum\_of\_squares} is $N'=a^2+b^2$ from Leibniz and the
cycle, \texttt{pyth\_discriminant} the resulting square, \texttt{split\_identity} and
\texttt{split\_product} the two--layer form of Proposition~\ref{C-prop:split}, and
\texttt{mass\_ge\_two\_of\_pythagorean} with \texttt{card\_ge\_59\_of\_pythagorean} this corollary,
glued by \texttt{mass\_ge\_two\_of\_pyth\_support} into \texttt{emptytest}, which takes the
Pythagorean equation itself rather than an assumed mass bound. Earlier editions proved the two ends
and not the link, so this tag overstated what was present.
The AM--GM step of Proposition~\ref{prop:kcycles} is \texttt{sum\_ge\_card\_of\_prod\_eq\_one}, valid
for every $k$ at once and proved from $\log x\le x-1$ rather than a $k$--term AM--GM. Bado's converse
is cited, not reproved \textup{(\texttt{lean/Erdos307/Pythagorean.lean})}.
\end{corollary}

\begin{remark}
(a) As $N'=a^2+b^2$ is a sum of two coprime squares, every odd prime dividing $N'$ is
$\equiv1\pmod4$; a cheap union--side pre--filter complementing Bado's congruence filters. (When
exactly one of $a,b$ is even this is the hypotenuse of the primitive triple $(|a^2-b^2|,2ab,a^2+b^2)$.)
(b) Since Seva's $s$ agrees with $n'$ on squarefree integers, and his valuation property forces
two--cycle members squarefree, his question is \emph{exactly} \#307, not merely an analogue.
(c) Primitive Pythagorean triples form the Barning--Hall ternary tree (the orbit structure of a
reflection group in $O(2,1;\mathbb{Z})$), and a \#307 solution's triple $(|a^2-b^2|,2ab,a^2+b^2)$
occupies one node with generator $(a,b)=(\prod P,\prod Q)$. The tree nonetheless yields no descent on
the problem: its moves $(a,b)\mapsto(2a\mp b,\,a),\,(a+2b,\,b)$ are linear, preserving coprimality but
not squarefreeness (over random squarefree coprime generators only $37\%$ of children keep both
members squarefree) so the locus of \#307--eligible nodes (squarefree generators of disjoint prime
support meeting $a^2+b^2=(ab)'$) is not tree--invariant. The additive tree is transverse to the
multiplicative constraint; the arithmetically natural recursion is instead the prime--append
$\delta$--tree of Proposition~\ref{C-prop:delta}, $\delta(Mp)=p\,\delta(M)+M$, which preserves
squarefreeness and tracks the line parameter, which is why the classical members organise on
\emph{that} tree, not Barning--Hall.
\end{remark}

The deviation parameter of Proposition~\ref{prop:pairform} organises the one--equation problem into a
short graded family.


\begin{thebibliography}{9}
\bibitem{Barbeau1961} E.~J.~Barbeau, \emph{Remarks on an arithmetic derivative}, Canad.\ Math.\
Bull.\ \textbf{4} (1961), 117--122.
\bibitem{Barbeau1976} E.~J.~Barbeau, \emph{Computer challenge corner: Problem 477}, J.\ Recreational
Math.\ \textbf{9} (1976/77), 30.
\bibitem{ErdosGraham1980} P.~Erd\H{o}s and R.~L.~Graham, \emph{Old and New Problems and Results in
Combinatorial Number Theory}, Monogr.\ Enseign.\ Math.\ \textbf{28}, Geneva, 1980, p.~38.
\bibitem{GrahamEgyptian} R.~L.~Graham, \emph{Paul Erd\H{o}s and Egyptian fractions}, in \emph{Erd\H{o}s
Centennial}, Bolyai Soc.\ Math.\ Stud.\ \textbf{25}, Springer, 2013, pp.~289--309.
\bibitem{Watanabe2020} T.~Watanabe, \emph{New examples of the representation of $1$ by the sum of
reciprocals of semiprime numbers}, arXiv:2009.03275 (2020).
\bibitem{BloomShifted} T.~F.~Bloom, \emph{Unit fractions with shifted prime denominators}, Proc.\
Roy.\ Soc.\ Edinburgh Sect.\ A (2024); arXiv:2305.02689.
\bibitem{UfnarovskiAhlander2003} V.~Ufnarovski and B.~\AA hlander, \emph{How to differentiate a
number}, J.\ Integer Seq.\ \textbf{6} (2003), Article 03.3.4.
\bibitem{Kovic2012} J.~Kovi\v{c}, \emph{The arithmetic derivative and antiderivative}, J.\ Integer
Seq.\ \textbf{15} (2012), Article 12.3.8.
\bibitem{MO320838} MathOverflow, \emph{Can the product of two sums of prime reciprocals be $1$?},
question 320838 (comments by R.~Israel and J.~Rosen, answer by ``Woett''), 2019.
\url{https://mathoverflow.net/questions/320838}.
\bibitem{ErdosProblems307} T.~F.~Bloom, \emph{Erd\H{o}s problem \#307},
\url{https://www.erdosproblems.com/307}.
\bibitem{OEIS_A003415} OEIS Foundation, \emph{Sequence A003415 (arithmetic derivative)},
\url{https://oeis.org/A003415}.
\bibitem{teRiele1984} H.~J.~J.~te~Riele, \emph{Computation of all the amicable pairs below
$10^{10}$}, Math.\ Comp.\ \textbf{47} (1986), 361--368.
\bibitem{BHP2001} R.~C.~Baker, G.~Harman, and J.~Pintz, \emph{The difference between consecutive
primes, II}, Proc.\ London Math.\ Soc.\ (3) \textbf{83} (2001), 532--562.
\bibitem{Iwaniec1978} H.~Iwaniec, \emph{Almost-primes represented by quadratic polynomials},
Invent.\ Math.\ \textbf{47} (1978), 171--188.
\bibitem{LemkeOliver2012} R.~J.~Lemke Oliver, \emph{Almost-primes represented by quadratic
polynomials}, Acta Arith.\ \textbf{151} (2012), 241--261.
\bibitem{HardyWright} G.~H.~Hardy and E.~M.~Wright, \emph{An Introduction to the Theory of Numbers},
6th ed., Oxford Univ.\ Press, 2008.
\bibitem{Tenenbaum} G.~Tenenbaum, \emph{Introduction to Analytic and Probabilistic Number Theory},
3rd ed., Grad.\ Stud.\ Math.\ \textbf{163}, Amer.\ Math.\ Soc., 2015 (Erd\H{o}s--Wintner theorem).
\bibitem{Stothers} W.~W.~Stothers, \emph{Polynomial identities and Hauptmoduln}, Quart.\ J.\ Math.\
Oxford (2) \textbf{32} (1981), 349--370; R.~C.~Mason, \emph{Diophantine Equations over Function
Fields}, Cambridge Univ.\ Press, 1984.
\bibitem{Seva323194} Seva (\texttt{mathoverflow.net} user), \emph{A recursion with a number--theoretic
function}, MathOverflow question 323194 (2019). \url{https://mathoverflow.net/questions/323194}.
\bibitem{PortFillings2026} \emph{Port fillings for primary pseudoperfect numbers}, arXiv:2605.21518
(May 2026).
\bibitem{Semiprime2026} \emph{Every natural number is a sum of distinct semiprime unit fractions},
arXiv:2606.15159 (June 2026).
\bibitem{ButskeJajeMayernik} W.~Butske, L.~M.~Jaje and D.~R.~Mayernik, \emph{On the equation
$\sum_{p\mid N}1/p+1/N=1$, pseudoperfect numbers, and perfectly weighted graphs}, Math.\ Comp.\
\textbf{69} (2000), 407--420.
\bibitem{Bado2026} I.~O.~Bado, \emph{Structural constraints for an Erd\H{o}s unit--fraction problem
over primes}, preprint, Aix--Marseille University (May 2026).
\url{https://doi.org/10.13140/RG.2.2.32245.54245}.
\bibitem{Giuga1950} G.~Giuga, \emph{Su una presumibile propriet\`a caratteristica dei numeri primi},
Ist.\ Lombardo Sci.\ Lett.\ Rend.\ Cl.\ Sci.\ Mat.\ Nat.\ (3) \textbf{14} (1950), 511--528.
\bibitem{BBBG1996} D.~Borwein, J.~M.~Borwein, P.~B.~Borwein, and R.~Girgensohn, \emph{Giuga's conjecture
on primality}, Amer.\ Math.\ Monthly \textbf{103} (1996), no.~1, 40--50.
\bibitem{BorweinWong1997} J.~M.~Borwein and E.~Wong, \emph{A survey of results relating to Giuga's
conjecture on primality}, CRM Proc.\ Lecture Notes \textbf{11} (1997), 13--27.
\bibitem{BJM2000} W.~Butske, L.~M.~Jaje, and D.~R.~Mayernik, \emph{On the equation
$\sum_{p\mid N}1/p+1/N=1$, pseudoperfect numbers, and perfectly weighted graphs}, Math.\ Comp.\
\textbf{69} (2000), no.~229, 407--420.
\bibitem{GrauOllerMarcen2012} J.~M.~Grau and A.~M.~Oller-Marc\'en, \emph{Giuga numbers and the
arithmetic derivative}, J.\ Integer Seq.\ \textbf{15} (2012), Article 12.4.1; arXiv:1103.2298.
\bibitem{SondowMacMillan2017} J.~Sondow and K.~MacMillan, \emph{Primary pseudoperfect numbers,
arithmetic progressions, and the Erd\H{o}s--Moser equation}, Amer.\ Math.\ Monthly \textbf{124} (2017),
no.~3, 232--240; arXiv:1812.06566.
\bibitem{GrauOllerMarcenSadornil2021} J.~M.~Grau, A.~M.~Oller-Marc\'en, and D.~Sadornil,
\emph{On $\mu$-Sondow numbers}, arXiv:2111.14211 (2021); Acta Math.\ Hungar.\ (2022).
\bibitem{Cox} D.~A.~Cox, \emph{Primes of the Form $x^2+ny^2$: Fermat, Class Field Theory, and Complex
Multiplication}, 2nd ed., Wiley, 2013.
\bibitem{OEIS_A054377} OEIS Foundation, \emph{Sequence A054377 (primary pseudoperfect numbers)},
\url{https://oeis.org/A054377}.
\bibitem{OEIS_A007850} OEIS Foundation, \emph{Sequence A007850 (Giuga numbers)},
\url{https://oeis.org/A007850}.
\bibitem{OEIS_A396915} OEIS Foundation, \emph{Sequence A396915 (composite squarefree $k$ with
$k'+2k$ a perfect square)}, \url{https://oeis.org/A396915}.
\bibitem{Evertse1984} J.-H.~Evertse, \emph{On sums of $S$-units and linear recurrences}, Compositio
Math.\ \textbf{53} (1984), 225--244.
\bibitem{FI1998} J.~Friedlander and H.~Iwaniec, \emph{The polynomial $X^2+Y^4$ captures its
primes}, Ann. of Math. \textbf{148} (1998), 945--1040.
\bibitem{HB2001} D.~R.~Heath-Brown, \emph{Primes represented by $x^3+2y^3$}, Acta Math.
\textbf{186} (2001), 1--84.
\bibitem{Stewart2013} C.~L.~Stewart, \emph{On divisors of Lucas and Lehmer numbers}, Acta Math.
\textbf{211} (2013), 291--314.
\bibitem{BCZ2003} Y.~Bugeaud, P.~Corvaja, and U.~Zannier, \emph{An upper bound for the G.C.D. of
$a^n-1$ and $b^n-1$}, Math. Z. \textbf{243} (2003), 79--84.
\bibitem{Merikoski2019} J.~Merikoski, \emph{On the largest prime factor of $n^2+1$}, J. Eur. Math.
Soc. \textbf{25} (2023), 1253--1284.
\bibitem{BGS2010} J.~Bourgain, A.~Gamburd, and P.~Sarnak, \emph{Affine linear sieve, expanders,
and sundry applications}, Invent. Math. \textbf{179} (2010), 559--644.
\bibitem{ESS2002} J.-H.~Evertse, H.~P.~Schlickewei, and W.~M.~Schmidt, \emph{Linear equations in
variables which lie in a multiplicative group}, Ann.\ of Math.\ (2) \textbf{155} (2002), 807--836.
\bibitem{Erdos1952} P.~Erd\H{o}s, \emph{On the sum $\sum_{k=1}^x d(f(k))$}, J.\ London Math.\ Soc.\
\textbf{27} (1952), 7--15.
\bibitem{BHV2001} Y.~Bilu, G.~Hanrot, and P.~M.~Voutier, \emph{Existence of primitive divisors of
Lucas and Lehmer numbers} (with an appendix by M.~Mignotte), J.\ Reine Angew.\ Math.\ \textbf{539}
(2001), 75--122.
\end{thebibliography}
\end{document}
